% Le Cameroun et les structures d'intégration sous-régionales — ECM Tle Tech — Cours signé OnBuch+
% Programme officiel MINESEC. Compiler avec : tectonic N05-cameroun-structures-integration-sous-regionales.tex
\newcommand{\DOCMATIERE}{ECM}
\newcommand{\DOCNIVEAU}{Tle Tech}
\newcommand{\DOCMODULE}{ECM – classe de Terminale technique}
\newcommand{\DOCLECON}{N05}
\newcommand{\DOCTITRE}{Le Cameroun et les structures d'intégration sous-régionales}
% OnBuch+ — préambule « cours signés » ENRICHI (compilé avec tectonic / XeLaTeX)
% Système visuel « Pop » d'OnBuch+ : fond crème (jamais blanc), bordures encre
% épaisses, ombres « dures » décalées, titres Archivo Black en capitales.
% Version MATHÉMATIQUES : graphes (pgfplots/tikz), boîtes pédagogiques
% (définition, propriété, méthode, attention, le savais-tu, exercice résolu,
% exercice, TP, bilan), page de garde et signature OnBuch+ ancrée en bas.
%
% Macros attendues dans main.tex AVANT \input{preamble} :
%   \DOCMATIERE  \DOCNIVEAU  \DOCMODULE  \DOCLECON (numéro)  \DOCTITRE
\documentclass[11pt]{article}

\usepackage{fontspec}
\usepackage[french]{babel}
\usepackage[a4paper,margin=2cm,top=3.4cm,bottom=2.8cm,headheight=1.4cm,footskip=1.3cm]{geometry}
\usepackage{amsmath,amssymb,amsfonts}
\usepackage{mathtools} % \xmapsto, \coloneqq... souvent utilisés par les modèles
\usepackage{cancel} % simplifications barrées (bilans de forces)
% \abs{z} (module, valeur absolue) et \norm{u} (norme), utilisés par les modèles
\providecommand{\abs}{}\let\abs\relax\DeclarePairedDelimiter{\abs}{\lvert}{\rvert}
\providecommand{\norm}{}\let\norm\relax\DeclarePairedDelimiter{\norm}{\lVert}{\rVert}
\usepackage[table]{xcolor} % [table] : fournit aussi \rowcolors (alternance de lignes)
\usepackage{pagecolor}
\usepackage{tikz}
\usetikzlibrary{arrows.meta,positioning,calc,shapes.geometric,shapes.misc,shapes.arrows,decorations.pathmorphing,decorations.pathreplacing,decorations.markings,patterns,shadows,mindmap,trees,fit,backgrounds,matrix,intersections,angles,quotes}
\usepackage{pgfplots}
\usepackage[version=4]{mhchem} % équations nucléaires \ce{^{235}_{92}U}
\usepackage[european,siunitx]{circuitikz} % schémas électriques
\pgfplotsset{compat=1.18}
\usepgfplotslibrary{fillbetween,groupplots} % aires entre courbes (calcul intégral)
\usepackage{enumitem}
\usepackage{fancyhdr}
% [most] charge toutes les bibliothèques tcolorbox (listings, minted, raster,
% documentation...) et gonfle inutilement la mémoire TeX sur un document long
% (~300 boîtes) : on ne charge que ce qui est réellement utilisé.
\usepackage[skins,breakable]{tcolorbox}
\usepackage{graphicx}
\usepackage{colortbl}
\usepackage{booktabs}
\usepackage{tabularx}
\usepackage{multirow}
\usepackage{diagbox} % case d'en-tête diagonale des tableaux à double entrée
% Colonnes extensibles centrée / gauche / droite (C, L, R), souvent utilisées par les modèles
\newcolumntype{C}{>{\centering\arraybackslash}X}
\newcolumntype{L}{>{\raggedright\arraybackslash}X}
\newcolumntype{R}{>{\raggedleft\arraybackslash}X}
\usepackage{array}
\usepackage{multicol}
\usepackage{siunitx}
% parse-numbers=false : accepte aussi \SI{2,5 \times 10^{-3}}{...}
\sisetup{locale=FR,output-decimal-marker={,},group-minimum-digits=4,parse-numbers=false}
\DeclareSIUnit\ohm{\ensuremath{\symup{\Omega}}} % U+2126 absent de la police
\DeclareSIUnit\division{div} % réglages d'oscilloscope (ms/div, V/div)
\DeclareSIUnit\tr{tr} % tours (vitesse de rotation, ex. \SI{1500}{\tr\per\minute})
\DeclareSIUnit\molar{M} % concentration molaire (ex. \SI{3}{\molar}, \SI{1}{\micro\molar})
\DeclareSIUnit\an{an} % année (le modèle confond parfois avec \year, un compteur TeX) — ex. \SI{5}{\kilo\meter\per\an}
\DeclareSIUnit\FCFA{FCFA} % franc CFA (ex. \SI{500}{\FCFA\per\kilo\gram})
\DeclareSIUnit\degreeBrix{\textdegree Brix} % teneur en sucre d'un sirop/jus (réfractométrie)
\DeclareSIUnit\month{mois} % le modèle utilise parfois \month, un compteur TeX, comme unité de durée
\DeclareSIUnit\microliter{\micro\liter} % le modèle écrit parfois \microliter en un seul token
\DeclareSIUnit\million{millions} % dénombrement (ex. \SI{15}{\million\per\mL} spermatozoïdes)
\usepackage{qrcode}
% \degree : commande fréquemment utilisée par les modèles pour un angle ou une
% température en dehors de \SI{}{} (non fournie par les paquets chargés).
\providecommand{\degree}{^{\circ}}
% \qed (fin de démonstration), utilisé par les modèles sans amsthm
\providecommand{\qed}{\hfill$\square$}
% \barre : double barre des valeurs interdites dans un tableau de variations (tkz-tab)
\providecommand{\barre}{\ensuremath{\|}}
% environnement proof (amsthm non chargé) utilisé par les modèles
\ifdefined\proof\else\newenvironment{proof}[1][Démonstration]{\par\noindent\textit{#1.}\ }{\qed\par}\fi
\usepackage{lastpage}
\usepackage{hyperref}
\hypersetup{hidelinks}

% ============================================================
% Palette « Pop » OnBuch+ — mêmes hex que la classe P (plus_ui.dart)
% ============================================================
\definecolor{poporange}{HTML}{F59321}
\definecolor{popdark}{HTML}{E07A0C}
\definecolor{popcream}{HTML}{FFF6E8}
\definecolor{popink}{HTML}{1C1714}
\definecolor{poppurple}{HTML}{7A5AE0}
\definecolor{popgreen}{HTML}{2E9E6A}
\definecolor{poppink}{HTML}{E0457B}
\definecolor{popblue}{HTML}{2F7DE1}
\definecolor{popgold}{HTML}{C98A12}
\definecolor{popmuted}{HTML}{8A7F76}
\definecolor{popline}{HTML}{EADFD2}
\definecolor{popgray}{HTML}{8A7F76}
\colorlet{popgrey}{popgray}
% Alias tolérants (le modèle nomme parfois les couleurs différemment)
\colorlet{popwhite}{white}
\colorlet{popblack}{popink}
\colorlet{popred}{poppink}
\colorlet{popyellow}{popgold}
% Teintes claires (fonds de tableaux/encadrés) — le modèle nomme parfois
% les couleurs avec un suffixe L (ex. popblueL) sans les avoir définies.
\colorlet{poporangeL}{poporange!20}
\colorlet{popdarkL}{popdark!20}
\colorlet{poppurpleL}{poppurple!20}
\colorlet{popgreenL}{popgreen!20}
\colorlet{poppinkL}{poppink!20}
\colorlet{popblueL}{popblue!20}
\colorlet{popgoldL}{popgold!20}
\colorlet{popredL}{popred!20}
\colorlet{popgrayL}{popgray!20}
% Teintes claires pour fonds de tableaux / zones
\colorlet{popblueL}{popblue!10!white}
\colorlet{poporangeL}{poporange!14!white}
\colorlet{popgreenL}{popgreen!12!white}
\colorlet{poppurpleL}{poppurple!12!white}
\colorlet{poppinkL}{poppink!12!white}
\colorlet{popgoldL}{popgold!14!white}

\pagecolor{popcream}

% ============================================================
% Typographie
% ============================================================
\defaultfontfeatures{Ligatures=TeX}
\setmainfont{PlusJakartaSans-Medium.ttf}[Path=../../../fonts/,BoldFont=PlusJakartaSans-Bold.ttf]
% Maths en sans-serif (Fira Math) avec chiffres et lettres droites de Plus
% Jakarta, pour que les formules se fondent dans le texte.
\usepackage[math-style=ISO,bold-style=ISO]{unicode-math}
\setmathfont{FiraMath-Regular.otf}[Path=../../../fonts/]
\setmathfont{PlusJakartaSans-Medium.ttf}[Path=../../../fonts/,range={up/num,up/{Latin,latin},\mathup}]
\usepackage{icomma} % virgule décimale sans espace en mode math
% Caractères mathématiques Unicode absents des polices de texte (Plus Jakarta,
% Archivo Black) : rendus par leur équivalent mathématique, en texte comme en
% formule — sinon ils s'affichent en carré vide (ex. « Arithmétique dans ℤ »).
\usepackage{newunicodechar}
\newunicodechar{ℝ}{\ensuremath{\mathbb{R}}}
\newunicodechar{ℤ}{\ensuremath{\mathbb{Z}}}
\newunicodechar{ℂ}{\ensuremath{\mathbb{C}}}
\newunicodechar{ℕ}{\ensuremath{\mathbb{N}}}
\newunicodechar{ℚ}{\ensuremath{\mathbb{Q}}}
\newunicodechar{𝔻}{\ensuremath{\mathbb{D}}}
\newunicodechar{α}{\ensuremath{\alpha}}
\newunicodechar{β}{\ensuremath{\beta}}
\newunicodechar{γ}{\ensuremath{\gamma}}
\newunicodechar{δ}{\ensuremath{\delta}}
\newunicodechar{ε}{\ensuremath{\varepsilon}}
\newunicodechar{θ}{\ensuremath{\theta}}
\newunicodechar{λ}{\ensuremath{\lambda}}
\newunicodechar{μ}{\ensuremath{\mu}}
\newunicodechar{σ}{\ensuremath{\sigma}}
\newunicodechar{τ}{\ensuremath{\tau}}
\newunicodechar{φ}{\ensuremath{\varphi}}
\newunicodechar{ψ}{\ensuremath{\psi}}
\newunicodechar{ω}{\ensuremath{\omega}}
\newunicodechar{Σ}{\ensuremath{\Sigma}}
% majuscules produites par \MakeUppercase dans les titres (α-aminés -> Α)
\newunicodechar{Α}{\ensuremath{\alpha}}
\newunicodechar{Β}{\ensuremath{\beta}}
\newunicodechar{Γ}{\ensuremath{\Gamma}}
\newunicodechar{Δ}{\ensuremath{\Delta}}
\newunicodechar{Ε}{\ensuremath{\varepsilon}}
\newunicodechar{Θ}{\ensuremath{\Theta}}
\newunicodechar{Λ}{\ensuremath{\Lambda}}
\newunicodechar{Μ}{\ensuremath{\mu}}
\newunicodechar{Τ}{\ensuremath{\tau}}
\newunicodechar{Φ}{\ensuremath{\Phi}}
\newunicodechar{Ψ}{\ensuremath{\Psi}}
\newunicodechar{Ω}{\ensuremath{\Omega}}
\newunicodechar{↦}{\ensuremath{\mapsto}}
\newunicodechar{∈}{\ensuremath{\in}}
\newunicodechar{∉}{\ensuremath{\notin}}
\newunicodechar{∧}{\ensuremath{\wedge}}
\newunicodechar{⇔}{\ensuremath{\Leftrightarrow}}
\newunicodechar{⟺}{\ensuremath{\Longleftrightarrow}}
\newunicodechar{⇒}{\ensuremath{\Rightarrow}}
\newunicodechar{⟹}{\ensuremath{\Longrightarrow}}
\newunicodechar{∘}{\ensuremath{\circ}}
\newunicodechar{∪}{\ensuremath{\cup}}
\newunicodechar{∩}{\ensuremath{\cap}}
\newunicodechar{≡}{\ensuremath{\equiv}}
\newunicodechar{⊂}{\ensuremath{\subset}}
\newunicodechar{⊥}{\ensuremath{\perp}}
\newunicodechar{∃}{\ensuremath{\exists}}
\newunicodechar{∀}{\ensuremath{\forall}}
\newunicodechar{∛}{\ensuremath{\sqrt[3]{\,}}}
\newunicodechar{∜}{\ensuremath{\sqrt[4]{\,}}}
\newunicodechar{⃗}{\ensuremath{{}^{\to}}}
\newunicodechar{ⁿ}{\ifmmode^{n}\else\textsuperscript{\ensuremath{n}}\fi}
\newunicodechar{ˣ}{\ifmmode^{x}\else\textsuperscript{\ensuremath{x}}\fi}
\newunicodechar{ᵇ}{\ifmmode^{b}\else\textsuperscript{\ensuremath{b}}\fi}
\newunicodechar{ʳ}{\ifmmode^{r}\else\textsuperscript{\ensuremath{r}}\fi}
\newunicodechar{ᵘ}{\ifmmode^{u}\else\textsuperscript{\ensuremath{u}}\fi}
\newunicodechar{ʸ}{\ifmmode^{y}\else\textsuperscript{\ensuremath{y}}\fi}
\newunicodechar{ᵃ}{\ifmmode^{a}\else\textsuperscript{\ensuremath{a}}\fi}
\newunicodechar{ᵉ}{\ifmmode^{e}\else\textsuperscript{\ensuremath{e}}\fi}
\newunicodechar{ᵏ}{\ifmmode^{k}\else\textsuperscript{\ensuremath{k}}\fi}
\newunicodechar{ᵖ}{\ifmmode^{p}\else\textsuperscript{\ensuremath{p}}\fi}
\newunicodechar{ᴺ}{\ifmmode^{N}\else\textsuperscript{\ensuremath{N}}\fi}
\newunicodechar{ᶜ}{\ifmmode^{c}\else\textsuperscript{\ensuremath{c}}\fi}
\newunicodechar{ʰ}{\ifmmode^{h}\else\textsuperscript{\ensuremath{h}}\fi}
\newunicodechar{ᵠ}{\ifmmode^{q}\else\textsuperscript{\ensuremath{q}}\fi}
\newunicodechar{ₙ}{\ifmmode_{n}\else\textsubscript{\ensuremath{n}}\fi}
\newunicodechar{ₐ}{\ifmmode_{a}\else\textsubscript{\ensuremath{a}}\fi}
\newunicodechar{ᵢ}{\ifmmode_{i}\else\textsubscript{\ensuremath{i}}\fi}
\newunicodechar{ᵣ}{\ifmmode_{r}\else\textsubscript{\ensuremath{r}}\fi}
\newunicodechar{ₖ}{\ifmmode_{k}\else\textsubscript{\ensuremath{k}}\fi}
\newunicodechar{ᵦ}{\ifmmode_{\beta}\else\textsubscript{\ensuremath{\beta}}\fi}
\newunicodechar{ₓ}{\ifmmode_{x}\else\textsubscript{\ensuremath{x}}\fi}
\newfontfamily\popdisplay{ArchivoBlack-Regular.ttf}[Path=../../../fonts/]
\newfontfamily\poplogo{SpaceGrotesk-Bold.ttf}[Path=../../../fonts/]
\newfontfamily\popsemi{PlusJakartaSans-SemiBold.ttf}[Path=../../../fonts/]
\newfontfamily\popxbold{PlusJakartaSans-ExtraBold.ttf}[Path=../../../fonts/]
\newfontfamily\popxboldls{PlusJakartaSans-ExtraBold.ttf}[Path=../../../fonts/,LetterSpace=3.0]

\setlist[enumerate]{leftmargin=1.7em,itemsep=5pt,topsep=4pt}
\setlist[itemize]{leftmargin=1.7em,itemsep=4pt,topsep=4pt,label={\color{poporange}\textbullet}}
\setlength{\parindent}{0pt}
\setlength{\parskip}{5pt}
\renewcommand{\arraystretch}{1.25}

% pgfplots : style Pop par défaut
\pgfplotsset{
  every axis/.append style={axis line style={popink,line width=1pt},tick style={popink},
    grid style={popline,line width=0.5pt},label style={font=\small},tick label style={font=\footnotesize}},
  popcurve/.style={poporange,line width=1.8pt,no markers,smooth},
}
% Même style disponible dans un \draw TikZ simple (hors pgfplots)
\tikzset{popcurve/.style={poporange,line width=1.8pt}}
\tikzset{popgrid/.style={popline,very thin}}
\colorlet{popcurve}{poporange} % le nom du style est parfois utilisé comme couleur

% ============================================================
% Pill / squiggle
% ============================================================
\newcommand{\poppill}[3]{%
  \tikz[baseline=-0.55ex]\node[draw=popink,line width=1.2pt,rounded corners=2.6pt,%
    fill=#2,inner xsep=6.5pt,inner ysep=3pt]{%
    \popxboldls\fontsize{7}{7}\selectfont\color{#3}\MakeUppercase{#1}};%
}
\newcommand{\popsquiggle}[1]{%
  \tikz[baseline=-0.4ex]{
    \draw[#1,line width=2.6pt,line cap=round]
      plot [smooth] coordinates {(0,0.09) (0.34,0.01) (0.68,0.21) (1.02,0.01) (1.36,0.10)};
  }%
}
% Mot-clé surligné (marqueur orange sous le texte)
\newcommand{\cle}[1]{{\popxbold\color{popdark}#1}}

% ============================================================
% En-tête / pied
% ============================================================
\pagestyle{fancy}
\fancyhf{}
\renewcommand{\headrulewidth}{0pt}
\renewcommand{\footrulewidth}{0pt}
\lhead{\raisebox{-0.15ex}{\poplogo\fontsize{13}{13}\selectfont\color{popink}OnBuch\color{poporange}+}}
\rhead{\poppill{\DOCMATIERE\ \textbullet\ \DOCNIVEAU\ \textbullet\ Leçon \DOCLECON}{white}{popink}}
\lfoot{\popxbold\fontsize{7.6}{7.6}\selectfont\color{popmuted}\MakeUppercase{Cours signé OnBuch+}\ \textbullet\ {\popsemi\DOCTITRE}}
\rfoot{\tikz[baseline=-0.6ex]\node[rounded corners=8.5pt,fill=popink,minimum height=17pt,minimum width=17pt,inner xsep=5pt,inner sep=0]{\popxbold\fontsize{8}{8}\selectfont\color{popcream}\thepage\,/\,\pageref*{LastPage}};}

% ============================================================
% Titres
% ============================================================
\newcommand{\chaptitre}[2]{%
  \begin{tcolorbox}[enhanced,colback=poporange,colframe=popink,boxrule=2.6pt,arc=11pt,
      drop shadow=popink,left=15pt,right=15pt,top=12pt,bottom=11pt,before skip=3pt,after skip=18pt]
    \poppill{#2}{popink}{popcream}\par\vspace{9pt}
    {\raggedright\hyphenpenalty=10000\exhyphenpenalty=10000\popdisplay\fontsize{22}{24}\selectfont\color{popcream}\MakeUppercase{#1}\par}
  \end{tcolorbox}
}
% Section numérotée : pastille orange avec le numéro + titre Archivo
\newcounter{popsec}
\newcommand{\coursec}[1]{%
  \stepcounter{popsec}\setcounter{popsub}{0}%
  \par\vspace{16pt}\noindent
  \begin{minipage}{\linewidth}
  \tikz[baseline=-0.7ex]\node[draw=popink,line width=1.6pt,fill=poporange,rounded corners=4pt,minimum size=22pt,inner sep=0]{\popdisplay\fontsize{12}{12}\selectfont\color{popcream}\thepopsec};%
  \hspace{8pt}{\popdisplay\fontsize{15}{17}\selectfont\color{popink}\MakeUppercase{#1}}\par
  \vspace{4pt}\noindent{\color{poporange}\rule{36pt}{3.6pt}}
  \end{minipage}\par\nopagebreak\vspace{8pt}
}
\newcounter{popsub}
\newcommand{\courssub}[1]{%
  \stepcounter{popsub}%
  \par\vspace{9pt}\noindent{\popxbold\fontsize{11.5}{13}\selectfont\color{popink}{\color{poporange}\thepopsec.\thepopsub}\hspace{6pt}#1}\par\nopagebreak\vspace{4pt}
}

% ============================================================
% Boîtes pédagogiques — même recette : fond blanc, bordure épaisse colorée,
% ombre dure encre, badge plein. Titre optionnel [#1].
% ============================================================
\tcbset{popbase/.style={breakable,enhanced,colback=white,boxrule=2.3pt,arc=9pt,
  shadow={3pt}{-3pt}{0pt}{popink},left=14pt,right=14pt,top=11pt,bottom=11pt,before skip=11pt,after skip=15pt,
  fontupper=\popsemi\fontsize{10.6}{14.5}\selectfont\color{popink},
  fontlower=\popsemi\fontsize{10.6}{14.5}\selectfont\color{popink}}}
% Coche dessinée (le glyphe ✓ n'existe pas dans les polices du document)
\newcommand{\popcheck}[1]{\tikz[baseline=-0.5ex]\draw[#1,line width=1.6pt,line cap=round,line join=round](-0.13,0.0)--(-0.04,-0.09)--(0.13,0.1);}
\providecommand{\coche}{\popcheck{popgreen}}
\providecommand{\resultat}[1]{\par\smallskip\noindent\fcolorbox{popgreen}{popgreenL}{\parbox{\dimexpr\linewidth-2\fboxsep-2\fboxrule\relax}{\textbf{Résultat.} #1}}\par\smallskip}
\newcommand{\popboxhead}[3]{\poppill{#1}{#2}{white}\ifx\relax#3\relax\else\hspace{6pt}{\popxbold\fontsize{10.6}{12}\selectfont\color{popink}#3}\fi\par\vspace{7pt}}

\newtcolorbox{aretenir}[1][]{popbase,colframe=popgreen,before upper={\popboxhead{À retenir}{popgreen}{#1}}}
\newtcolorbox{exemplebox}[1][]{popbase,colframe=poporange,before upper={\popboxhead{Exemple}{poporange}{#1}}}
\newtcolorbox{definition}[1][]{popbase,colframe=popblue,before upper={\popboxhead{Définition}{popblue}{#1}}}
\newtcolorbox{propriete}[1][]{popbase,colframe=poppurple,before upper={\popboxhead{Propriété}{poppurple}{#1}}}
\newtcolorbox{methode}[1][]{popbase,colframe=popgold,before upper={\popboxhead{Méthode}{popgold}{#1}}}
\newtcolorbox{attention}[1][]{popbase,colframe=poppink,before upper={\popboxhead{Attention}{poppink}{#1}}}
\newtcolorbox{experience}[1][]{popbase,colframe=popblue,colback=popblueL,before upper={\popboxhead{Expérience}{popblue}{#1}}}
% « Le savais-tu ? » : carte violette pleine (contexte, culture, Cameroun)
\newtcolorbox{savaistu}[1][]{popbase,colback=poppurple,colframe=popink,
  fontupper=\popsemi\fontsize{10.4}{14.2}\selectfont\color{white},
  before upper={\poppill{Le savais-tu\,?}{white}{poppurple}\ifx\relax#1\relax\else\hspace{6pt}{\popxbold\color{white}#1}\fi\par\vspace{7pt}}}
% Exercice résolu : énoncé en haut, \tcblower puis la solution sur fond crème
\newcounter{popexo}\newcounter{popexr}
\newtcolorbox{exoresolu}[1][]{popbase,colframe=poporange,colbacklower=popcream,
  segmentation style={popink,line width=1.2pt,dashed},
  before upper={\stepcounter{popexr}\popboxhead{Exercice résolu \thepopexr}{poporange}{#1}},
  before lower={\poppill{Solution}{popink}{popcream}\par\vspace{6pt}}}
% Exercice (non résolu en place) — la correction est dans la partie Corrigés
\newtcolorbox{exercice}[1][]{popbase,colframe=popink,
  before upper={\stepcounter{popexo}\popboxhead{Exercice \thepopexo}{popink}{#1}}}
% Corrigé
\newtcolorbox{corrige}[1][]{popbase,colframe=popgreen,colback=popgreenL,
  before upper={\popboxhead{Corrigé}{popgreen}{#1}}}
% Objectifs / prérequis / bilan
\newtcolorbox{objectifs}{popbase,colframe=popink,colback=poporangeL,
  before upper={\popboxhead{Objectifs de la leçon}{popink}{}}}
\newtcolorbox{prerequis}{popbase,colframe=popink,colback=popblueL,
  before upper={\popboxhead{Prérequis}{popblue}{}}}
\newtcolorbox{bilan}{popbase,colframe=popink,colback=white,boxrule=2.8pt,
  before upper={\popboxhead{Fiche bilan}{poporange}{L'essentiel à connaître pour l'examen}}}
% Figure encadrée avec légende
\newtcolorbox{popfigure}[1][]{enhanced,breakable=false,colback=white,colframe=popline,boxrule=1.4pt,arc=8pt,
  left=10pt,right=10pt,top=10pt,bottom=8pt,before skip=10pt,after skip=12pt,halign=center,
  lower separated=false,halign lower=center,
  fontlower=\popsemi\fontsize{9}{11.5}\selectfont\color{popmuted},
  #1}
\newcommand{\legende}[1]{\par\vspace{6pt}{\popsemi\fontsize{9}{11.5}\selectfont\color{popmuted}#1}}

% ============================================================
% Signature OnBuch+
% ============================================================
\newcommand{\poplogoinline}[1]{{\poplogo\fontsize{#1}{#1}\selectfont\color{popink}OnBuch\color{poporange}+}}
\newcommand{\signatureonbuch}{%
  \begin{tcolorbox}[enhanced,colback=popink,colframe=popink,boxrule=2.2pt,arc=10pt,
      drop shadow=poporange,left=14pt,right=14pt,top=10pt,bottom=10pt,before skip=0pt,after skip=0pt]
    \tikz[baseline=-0.7ex]\node[circle,fill=popgreen,draw=popcream,line width=1.3pt,minimum size=22pt,inner sep=0]{\popcheck{white}};%
    \hspace{9pt}\begin{minipage}[c]{0.62\linewidth}
      {\popxbold\fontsize{12}{13}\selectfont\color{popcream}Signé\ \textbullet\ {\poplogo OnBuch\color{poporange}+}}\par\vspace{2pt}
      {\raggedright\popsemi\fontsize{8}{10.5}\selectfont\color{popline}\MakeUppercase{Équipe pédagogique OnBuch+}\\\MakeUppercase{Conforme au programme officiel MINESEC}\par}
    \end{minipage}\hfill
    {\poplogo\fontsize{22}{22}\selectfont\color{popcream}OnBuch\color{poporange}+}
  \end{tcolorbox}%
}

% ============================================================
% Page de garde
% #1 titre · #2 matière · #3 niveau · #4 module · #5 n° leçon · #6 durée ·
% #7 description / objectifs (texte ou itemize)
% ============================================================
\newcommand{\pagegarde}[7]{%
  \thispagestyle{empty}
  \newgeometry{a4paper,margin=1.7cm,top=1.6cm,bottom=1.4cm}%
  \begin{tikzpicture}[remember picture,overlay]
    \fill[poporange!18] ([xshift=-3.2cm,yshift=-3.6cm]current page.north east) circle (4.6cm);
    \fill[poppurple!14] ([xshift=2.4cm,yshift=7.6cm]current page.south west) circle (3.8cm);
  \end{tikzpicture}%
  {\poplogo\fontsize{30}{30}\selectfont\color{popink}OnBuch\color{poporange}+}\hfill
  \raisebox{0.6ex}{\poppill{Cours signé \textbullet\ Premium}{popink}{popcream}}\par
  \vspace{1pt}\popsquiggle{poppurple}\par
  \vspace{8pt}
  \begin{tcolorbox}[enhanced,colback=poporange,colframe=popink,boxrule=2.8pt,arc=13pt,
      drop shadow=popink,left=18pt,right=18pt,top=12pt,bottom=13pt,after skip=10pt]
    \poppill{#2 \textbullet\ #3}{popink}{popcream}\hspace{5pt}\poppill{#4}{white}{popink}\par\vspace{8pt}
    {\popdisplay\fontsize{14}{14}\selectfont\color{popink}LEÇON #5}\par\vspace{4pt}
    {\raggedright\hyphenpenalty=10000\exhyphenpenalty=10000\popdisplay\fontsize{25}{27}\selectfont\color{popcream}\MakeUppercase{#1}\par}
    \vspace{8pt}
    {\popxbold\fontsize{9.5}{9.5}\selectfont\color{popink}\MakeUppercase{Durée conseillée : #6}}
  \end{tcolorbox}
  \begin{tcolorbox}[enhanced,colback=white,colframe=popink,boxrule=2.2pt,arc=10pt,
      drop shadow=popink,left=16pt,right=16pt,top=10pt,bottom=10pt,after skip=8pt]
    \poppill{Dans cette leçon}{poppurple}{white}\par\vspace{5pt}
    \popsemi\fontsize{9.3}{12.6}\selectfont\color{popink}#7
  \end{tcolorbox}
  \noindent\begin{minipage}[t]{0.487\linewidth}
    \begin{tcolorbox}[enhanced,colback=popgreen,colframe=popink,boxrule=2.2pt,arc=10pt,
        drop shadow=popink,left=11pt,right=11pt,top=10pt,bottom=10pt,height=3.1cm,valign=center]
      \tcbox[colback=white,colframe=popink,boxrule=1.3pt,arc=5pt,boxsep=0pt,nobeforeafter,
        left=3pt,right=3pt,top=3pt,bottom=3pt,valign=center]{\qrcode[height=1.9cm]{https://play.google.com/store/apps/details?id=com.luvvix.onbuch&hl=fr}}%
      \hspace{8pt}\begin{minipage}[c]{0.5\linewidth}
        {\popdisplay\fontsize{11}{12.5}\selectfont\color{white}\MakeUppercase{Télécharge OnBuch}}\par\vspace{4pt}
        {\raggedright\popsemi\fontsize{8.4}{11}\selectfont\color{white}Gratuit \textbullet\ Google Play\\Tuteur IA, annales, concours, résultats\par}
      \end{minipage}
    \end{tcolorbox}
  \end{minipage}\hfill
  \begin{minipage}[t]{0.487\linewidth}
    \begin{tcolorbox}[enhanced,colback=poppurple,colframe=popink,boxrule=2.2pt,arc=10pt,
        drop shadow=popink,left=11pt,right=11pt,top=10pt,bottom=10pt,height=3.1cm,valign=center]
      \tikz[baseline=-0.7ex]\node[circle,fill=white,draw=popink,line width=1.3pt,minimum size=1.6cm,inner sep=0]{\popdisplay\fontsize{24}{24}\selectfont\color{poppurple}+};%
      \hspace{8pt}\begin{minipage}[c]{0.6\linewidth}
        {\popdisplay\fontsize{11}{12.5}\selectfont\color{white}\MakeUppercase{Passe à OnBuch+}}\par\vspace{4pt}
        {\raggedright\popsemi\fontsize{8.4}{11}\selectfont\color{white}Tous les cours signés, Tuteur IA illimité, fiches PDF — onglet OnBuch+ de l'app\par}
      \end{minipage}
    \end{tcolorbox}
  \end{minipage}
  \vspace{8pt}
  \popcontenu
  \vfill
  \signatureonbuch
  \restoregeometry
  \newpage
}

% Rangée « Ce cours contient » (page de garde)
\newcommand{\poptile}[3]{% #1 couleur #2 gros texte #3 légende
  \begin{tcolorbox}[enhanced,colback=#1,colframe=popink,boxrule=2pt,arc=9pt,shadow={2.5pt}{-2.5pt}{0pt}{popink},
      left=6pt,right=6pt,top=9pt,bottom=9pt,halign=center,valign=center,height=1.75cm,width=\linewidth,nobeforeafter]
    {\popdisplay\fontsize{12.5}{13}\selectfont\color{white}\MakeUppercase{#2}}\par\vspace{3pt}
    {\centering\popsemi\fontsize{7.8}{9.5}\selectfont\color{white}#3\par}
  \end{tcolorbox}}
\newcommand{\popcontenu}{%
  \par\noindent{\popxboldls\fontsize{7.5}{7.5}\selectfont\color{popmuted}CE COURS SIGNÉ CONTIENT}\par\vspace{7pt}
  \noindent\begin{minipage}[t]{0.235\linewidth}\poptile{popblue}{Cours}{complet, illustré, pas à pas}\end{minipage}\hfill
  \begin{minipage}[t]{0.235\linewidth}\poptile{poporange}{Méthodes}{et exercices résolus}\end{minipage}\hfill
  \begin{minipage}[t]{0.235\linewidth}\poptile{poppink}{Exercices}{type examen + corrigés}\end{minipage}\hfill
  \begin{minipage}[t]{0.235\linewidth}\poptile{popgreen}{Fiche bilan}{l'essentiel à retenir}\end{minipage}\par}

% Sommaire
\newtcolorbox{sommaire}{enhanced,colback=white,colframe=popink,boxrule=2.2pt,arc=10pt,
  drop shadow=popink,left=18pt,right=18pt,top=13pt,bottom=13pt,before skip=6pt,after skip=20pt,
  before upper={\poppill{Sommaire}{poppurple}{white}\par\vspace{9pt}\popsemi\fontsize{10.6}{14.5}\selectfont\color{popink}}}

% Signature de fin de cours — poussée tout en bas de la dernière page
\newcommand{\findecours}{%
  \par\vfill\nopagebreak
  \begin{center}\popsquiggle{poporange}\end{center}\vspace{4pt}
  \signatureonbuch
}
\AtBeginDocument{\def\llbracket{\mathopen{[\mkern-3.5mu[}}\def\rrbracket{\mathclose{]\mkern-3.5mu]}}} % intervalles d'entiers (glyphe ⟦ absent de la police)
% Glyphes absents de Fira Math : redéfinis à partir de glyphes disponibles
\AtBeginDocument{%
  \def\setminus{\mathbin{\backslash}}\let\smallsetminus\setminus
  \def\vdots{\mathord{\vbox{\baselineskip3.2pt\lineskiplimit0pt\kern4pt\hbox{.}\hbox{.}\hbox{.}}}}%
  \def\ddots{\mathinner{\mkern1mu\raise6.5pt\hbox{.}\mkern2mu\raise3.6pt\hbox{.}\mkern2mu\raise0.7pt\hbox{.}\mkern1mu}}%
  \def\triangle{\mathord{\scalebox{0.9}{$\symup{\Delta}$}}}%
  \def\nabla{\mathord{\raisebox{\depth}{\scalebox{1}[-1]{$\symup{\Delta}$}}}}%
  \def\checkmark{\popcheck{popdark}}%
  \def\star{\mathbin{\ast}}\def\bigstar{\mathord{\ast}}%
  \def\diamond{\mathbin{\scalebox{0.6}{$\Diamond$}}}%
  \def\bigcup{\mathop{\mathchoice{\raisebox{-0.3em}{\scalebox{1.7}{$\cup$}}}{\scalebox{1.25}{$\cup$}}{\cup}{\cup}}\displaylimits}%
  \def\bigcap{\mathop{\mathchoice{\raisebox{-0.3em}{\scalebox{1.7}{$\cap$}}}{\scalebox{1.25}{$\cap$}}{\cap}{\cap}}\displaylimits}%
  \def\lesseqqgtr{\mathrel{\lessgtr}}\def\bigoplus{\mathop{\oplus}\displaylimits}%
}
\providecommand{\ssi}{si et seulement si} % abréviation fréquente
\AtBeginDocument{\providecommand{\wideparen}{\overparen}} % arc de cercle

\begin{document}

\pagegarde{Le Cameroun et les structures d'intégration sous-régionales}{ECM}{Tle Tech}{ECM – classe de Terminale technique}{N05}{2 séances (fiche de progression harmonisée)}{Cette leçon présente l'engagement du Cameroun au sein de l'Union Africaine et des structures d'intégration sous-régionales (CEEAC, CEMAC, BEAC, CEMAC...). Elle montre comment l'unité africaine et sous-régionale se traduit concrètement dans la vie des Camerounais : libre circulation, monnaie commune, coopération économique et sécuritaire.
\vspace{4pt}
\begin{multicols}{2}\small
\begin{itemize}
\item À la fin de cette leçon, je sais définir l'intégration sous-régionale et citer les principales organisations auxquelles le Cameroun appartient
\item À la fin de cette leçon, je sais présenter l'Union Africaine (histoire, organes, objectifs) et le rôle du Cameroun en son sein
\item À la fin de cette leçon, je sais décrire les structures sous-régionales du Cameroun (CEEAC, CEMAC, BEAC, COBAC, BDEAC) et leurs missions
\item À la fin de cette leçon, je sais localiser sur une carte les États membres de la CEEAC et de la CEMAC
\item À la fin de cette leçon, je sais identifier les avantages concrets de l'intégration pour les Camerounais (libre circulation, monnaie unique, marché commun)
\item À la fin de cette leçon, je sais analyser les obstacles à l'intégration sous-régionale en Afrique centrale
\end{itemize}\end{multicols}}

\begin{sommaire}
\begin{multicols}{2}
\begin{enumerate}
\item Comprendre l'intégration sous-régionale
\item Le Cameroun et l'Union Africaine (UA)
\item Le Cameroun dans la CEEAC
\item Le Cameroun dans la CEMAC et les institutions communes
\item Bilan : acquis et obstacles de l'intégration
\item TP guidé : cartographier l'intégration du Cameroun
\item Activité d'intégration
\item Méthodes et exercices résolus
\item Exercices
\item Corrigés des exercices
\item Fiche bilan
\end{enumerate}
\end{multicols}
\end{sommaire}

\begin{objectifs}
\begin{itemize}
\item À la fin de cette leçon, je sais définir l'intégration sous-régionale et citer les principales organisations auxquelles le Cameroun appartient
\item À la fin de cette leçon, je sais présenter l'Union Africaine (histoire, organes, objectifs) et le rôle du Cameroun en son sein
\item À la fin de cette leçon, je sais décrire les structures sous-régionales du Cameroun (CEEAC, CEMAC, BEAC, COBAC, BDEAC) et leurs missions
\item À la fin de cette leçon, je sais localiser sur une carte les États membres de la CEEAC et de la CEMAC
\item À la fin de cette leçon, je sais identifier les avantages concrets de l'intégration pour les Camerounais (libre circulation, monnaie unique, marché commun)
\item À la fin de cette leçon, je sais analyser les obstacles à l'intégration sous-régionale en Afrique centrale
\item À la fin de cette leçon, je sais appliquer les notions d'unité et de solidarité à des situations concrètes de la vie courante
\item À la fin de cette leçon, je sais rédiger et justifier une conclusion argumentée sur l'intérêt de l'intégration pour le Cameroun
\end{itemize}
\end{objectifs}

\begin{prerequis}
\begin{itemize}
\item Les organisations internationales (ONU) vues en classe de Première
\item La localisation du Cameroun en Afrique centrale et ses pays voisins (géographie de 4ème/3ème)
\item La notion de citoyenneté nationale et de solidarité (leçons précédentes de Terminale)
\item Les grandes étapes de l'indépendance du Cameroun et de l'unité nationale
\end{itemize}
\end{prerequis}

\newpage

\coursec{Comprendre l'intégration sous-régionale}

\textbf{Situation d'accroche.} En 2023, un commerçant de Douala exporte des bananes plantains vers Libreville sans visa, tandis qu'un étudiant camerounais part étudier à N'Djamena avec sa seule carte nationale d'identité. Comment est-ce possible ? Par ailleurs, le Cameroun accueille le siège de la Banque des États de l'Afrique Centrale (BEAC) à Yaoundé et participe chaque année aux sommets de l'Union Africaine à Addis-Abeba. Pourquoi le Cameroun s'engage-t-il dans ces organisations régionales et continentales ? Quels avantages concrets en tirent les Camerounais dans leur vie quotidienne ?

\textbf{Problématique.} Comment comprendre le phénomène d'intégration sous-régionale, et pourquoi le Cameroun y participe-t-il à plusieurs niveaux, de la sous-région au continent ?

\courssub{Qu'est-ce que l'intégration ?}

\textbf{L'intuition d'abord.} Imaginez un petit commerçant seul face à un grand distributeur étranger : il n'a aucun pouvoir de négociation. Maintenant, imaginez que dix commerçants se regroupent pour acheter ensemble leurs marchandises : ils obtiennent de meilleurs prix, se défendent mieux, et parlent d'une seule voix. C'est exactement la logique de l'\cle{intégration} entre États : seul, un pays est faible ; ensemble, plusieurs pays pèsent davantage, économiquement et diplomatiquement.

\begin{definition}[L'intégration]
L'\cle{intégration} est le processus par lequel des États voisins mettent en commun des moyens et des politiques pour renforcer leur unité et favoriser leur développement. Elle peut conduire à la création d'institutions communes, d'un marché unique ou d'une monnaie partagée.
\end{definition}

\begin{definition}[L'intégration sous-régionale]
L'\cle{intégration sous-régionale} est le regroupement d'États d'une même zone géographique qui coopèrent pour mutualiser leurs ressources et harmoniser leurs politiques (économiques, monétaires, sécuritaires, culturelles).
\end{definition}

\textbf{Pourquoi s'intégrer ?} Le raisonnement est simple et peut être vérifié par l'observation. Un petit pays seul, avec un marché de quelques millions d'habitants, attire peu les investisseurs et pèse peu dans les négociations internationales. En revanche, un ensemble de plusieurs pays offre un marché plus vaste, une voix plus forte dans les instances mondiales, des coûts partagés (une seule banque centrale au lieu de six, par exemple) et une meilleure sécurité collective. C'est le principe : \emph{l'union fait la force}.

\courssub{Les formes d'intégration}

L'intégration n'est pas un bloc uniforme : elle prend plusieurs formes, qui peuvent se combiner.

\begin{itemize}
\item \textbf{Intégration économique} : création d'un \cle{marché commun} (les marchandises circulent sans droits de douane entre les pays membres), d'un \cle{tarif extérieur commun} (tous les pays membres appliquent les mêmes droits de douane aux produits venant de l'extérieur) et parfois d'une \cle{monnaie unique}.
\item \textbf{Intégration politique} : mise en place d'\cle{institutions communes} (parlement, commission, conférence des chefs d'État) qui prennent des décisions valables pour tous les membres.
\item \textbf{Intégration sécuritaire} : coopération en matière de défense et de maintien de la paix (forces conjointes, lutte commune contre le terrorisme ou la piraterie).
\item \textbf{Intégration culturelle} : échanges scolaires, universitaires et artistiques, reconnaissance mutuelle des diplômes, promotion de langues et de valeurs partagées.
\end{itemize}

\begin{exemplebox}[Une situation concrète : le camion de Douala à Bangui]
Un transporteur camerounais charge un camion de marchandises à Douala à destination de Bangui (République centrafricaine). Grâce à l'intégration sous-régionale :
\begin{itemize}
\item les \textbf{formalités} sont simplifiées : les deux pays appartiennent à la CEMAC, les marchandises bénéficient de règles douanières harmonisées ;
\item la \textbf{monnaie} est la même : le franc CFA (XAF), émis par la BEAC, est accepté dans les deux pays, ce qui évite les opérations de change coûteuses ;
\item la \textbf{sécurité} est renforcée par la coopération entre les États de la sous-région sur les corridors routiers.
\end{itemize}
Sans intégration, ce même trajet coûterait plus cher, prendrait plus de temps et comporterait plus de risques.
\end{exemplebox}

\begin{aretenir}[Vocabulaire clé]
\begin{itemize}
\item \cle{Sous-région} : ensemble de pays voisins d'une même zone géographique (ex. : l'Afrique centrale).
\item \cle{Espace communautaire} : territoire des États membres d'une organisation, à l'intérieur duquel s'appliquent des règles communes.
\item \cle{Libre circulation} : droit pour les personnes et les biens de se déplacer entre les pays membres sans visa ni droits de douane internes.
\item \cle{Tarif extérieur commun} : droits de douane identiques appliqués par tous les pays membres aux importations venant des pays extérieurs.
\end{itemize}
\end{aretenir}

\courssub{Les cercles d'intégration du Cameroun}

Le Cameroun n'appartient pas à une seule organisation, mais à plusieurs ensembles emboîtés, comme des cercles concentriques : son quartier immédiat (la CEMAC), sa sous-région au sens large (la CEEAC), son continent (l'Union Africaine) et le monde (l'ONU).

\begin{popfigure}
\begin{center}
\begin{tikzpicture}
\fill[popgold!25] (0,0) circle (5.2);
\fill[popgreen!25] (0,0) circle (4.1);
\fill[popblue!25] (0,0) circle (3.0);
\fill[poporange!35] (0,0) circle (1.9);
\fill[poppink!45] (0,0) circle (0.85);
\node[font=\bfseries\small] at (0,0) {Cameroun};
\node[font=\small\bfseries, popdark] at (0,1.35) {CEMAC (6 États)};
\node[font=\small\bfseries, popdark] at (0,2.45) {CEEAC (11 États)};
\node[font=\small\bfseries, popdark] at (0,3.55) {Union Africaine (55 États)};
\node[font=\small\bfseries, popdark] at (0,4.65) {ONU (193 États)};
\end{tikzpicture}
\end{center}
\legende{Schéma en cercles concentriques des niveaux d'intégration du Cameroun : 1. Cameroun ; 2. CEMAC (union économique et monétaire) ; 3. CEEAC (organisation sous-régionale élargie) ; 4. Union Africaine (continent) ; 5. ONU (monde). Chaque niveau ajoute des partenaires et des règles communes.}
\end{popfigure}

\begin{exemplebox}[Quelques chiffres à retenir]
La \cle{CEMAC} (Communauté Économique et Monétaire de l'Afrique Centrale) compte \textbf{6 États} (Cameroun, Congo, Gabon, Guinée équatoriale, République centrafricaine, Tchad) et \textbf{plus de 55 millions d'habitants}. La \cle{CEEAC} (Communauté Économique des États de l'Afrique Centrale) en compte \textbf{11} : elle inclut les six pays de la CEMAC plus l'Angola, le Burundi, la République démocratique du Congo, le Rwanda et Sao Tomé-et-Principe.
\end{exemplebox}

\begin{attention}[Erreur fréquente : CEEAC ou CEMAC ?]
Il ne faut pas confondre la \textbf{CEEAC} et la \textbf{CEMAC}. La \textbf{CEEAC} est une organisation politico-économique \emph{large} (11 États) qui vise la coopération générale et la paix en Afrique centrale. La \textbf{CEMAC} est une union économique et monétaire \emph{plus restreinte} (6 États) : c'est elle qui gère la monnaie commune (le franc CFA) et le marché commun. À l'examen, précisez toujours de quelle organisation vous parlez.
\end{attention}

\begin{popfigure}
\begin{center}
\begin{tikzpicture}[scale=1.05]
\fill[popcream] (-5.5,-4.2) rectangle (5.5,4.4);
\fill[poppink!50] (-0.6,0.4) -- (0.4,0.9) -- (1.0,0.2) -- (0.9,-1.2) -- (0.2,-2.1) -- (-0.8,-2.0) -- (-1.3,-0.9) -- (-1.4,0.0) -- cycle;
\fill[popblue!30] (-1.4,0.0) -- (-0.6,0.4) -- (-0.9,1.6) -- (-1.8,2.6) -- (-2.9,2.4) -- (-2.6,1.2) -- (-2.0,0.3) -- cycle;
\fill[popgreen!30] (-2.9,2.4) -- (-1.8,2.6) -- (-2.2,3.8) -- (-3.6,3.9) -- (-4.2,3.0) -- cycle;
\fill[poporange!30] (0.4,0.9) -- (1.6,1.5) -- (2.4,0.8) -- (2.0,-0.2) -- (1.0,0.2) -- cycle;
\fill[popgold!35] (1.0,0.2) -- (2.0,-0.2) -- (2.3,-1.4) -- (1.4,-2.3) -- (0.9,-1.2) -- cycle;
\fill[poppurple!25] (0.9,-1.2) -- (1.4,-2.3) -- (0.6,-3.3) -- (-0.4,-3.2) -- (0.2,-2.1) -- cycle;
\fill[popgreen!20] (-0.8,-2.0) -- (0.2,-2.1) -- (-0.4,-3.2) -- (-1.3,-3.0) -- cycle;
\fill[popblue!20] (-1.3,-0.9) -- (-0.8,-2.0) -- (-1.3,-3.0) -- (-2.2,-2.4) -- (-2.1,-1.2) -- cycle;
\node[font=\scriptsize\bfseries] at (-0.35,-0.55) {CAMEROUN};
\node[font=\scriptsize] at (-1.9,1.4) {Tchad};
\node[font=\scriptsize] at (-3.1,3.3) {Nigeria};
\node[font=\scriptsize] at (1.6,0.9) {RCA};
\node[font=\scriptsize] at (1.6,-1.3) {Congo};
\node[font=\scriptsize] at (0.5,-2.8) {Gabon};
\node[font=\scriptsize] at (-1.7,-2.6) {Guinée éq.};
\draw[->, thick, popdark] (4.6,-3.6) -- (4.6,-2.6) node[above, font=\scriptsize\bfseries]{N};
\end{tikzpicture}
\end{center}
\legende{Carte très simplifiée (croquis non à l'échelle) du Cameroun et de ses voisins d'Afrique centrale : Nigeria, Tchad, République centrafricaine (RCA), Congo, Gabon, Guinée équatoriale. Le Cameroun est au carrefour de l'Afrique de l'Ouest et de l'Afrique centrale.}
\end{popfigure}

\courssub{Les principales organisations auxquelles appartient le Cameroun}

\begin{itemize}
\item \textbf{À l'échelle du continent} : l'\cle{Union Africaine} (UA), créée en 2002 à Durban en succédant à l'OUA (1963), réunit 55 États africains. Son siège est à Addis-Abeba (Éthiopie). Elle œuvre pour l'unité, la paix et le développement du continent.
\item \textbf{À l'échelle de la sous-région} : la \cle{CEEAC} (11 États, coopération politique et économique, siège à Libreville) et la \cle{CEMAC} (6 États, union économique et monétaire, siège à Bangui ; la BEAC, banque centrale commune, a son siège à Yaoundé).
\item \textbf{En complément (utile à connaître pour le Bac)} : le Cameroun est aussi membre de la \cle{CEN-SAD} (Communauté des États sahélo-sahariens) et du \cle{Commonwealth}, héritage de sa double histoire coloniale française et britannique. Ces appartenances complètent son ancrage africain.
\end{itemize}

\begin{savaistu}[Un pays carrefour]
Le Cameroun est le seul pays membre à la fois de la CEEAC, de la CEMAC, du Commonwealth et de la Francophonie. Cette position unique, héritée de son histoire (double colonisation allemande puis française et britannique), fait de lui un véritable pont entre l'Afrique anglophone et l'Afrique francophone : on l'appelle parfois « l'Afrique en miniature ».
\end{savaistu}

\courssub{Méthode : analyser un document en ECM}

\begin{methode}[Analyser une carte ou un texte sur l'intégration]
\begin{enumerate}
\item \textbf{Lire la légende et les sources} : que représentent les couleurs, les symboles, les chiffres ? Qui a produit le document et quand ?
\item \textbf{Identifier les acteurs} : quels États, quelles organisations, quelles institutions sont cités ?
\item \textbf{Dégager l'idée principale} : que veut montrer le document (une coopération, un obstacle, un progrès) ?
\item \textbf{Relier à la notion d'unité} : en quoi ce document illustre-t-il l'intégration sous-régionale, ses avantages ou ses limites ?
\item \textbf{Conclure par une phrase rédigée} qui répond à la question posée et justifie votre réponse.
\end{enumerate}
\end{methode}

\begin{exoresolu}[Analyser une situation d'intégration]
Énoncé. Un élève lit cette affirmation : « Grâce à la CEMAC, un étudiant camerounais peut partir étudier à N'Djamena (Tchad) avec sa seule carte nationale d'identité. » En vous appuyant sur la méthode d'analyse, montrez en quoi cette situation illustre l'intégration sous-régionale.
\tcblower
Solution détaillée.
\begin{enumerate}
\item \textbf{Acteurs identifiés} : le Cameroun et le Tchad, deux États membres de la CEMAC ; un citoyen camerounais (l'étudiant).
\item \textbf{Idée principale} : entre les pays membres de la CEMAC, les personnes circulent librement, sans visa.
\item \textbf{Lien avec la notion d'unité} : cette situation illustre la \cle{libre circulation} des personnes, l'un des acquis concrets de l'intégration sous-régionale. En harmonisant leurs politiques, les États de la CEMAC créent un \cle{espace communautaire} où les citoyens jouissent de droits communs.
\item \textbf{Conclusion rédigée} : la possibilité pour un étudiant camerounais d'étudier au Tchad sans visa montre que l'intégration sous-régionale n'est pas une idée abstraite : elle améliore concrètement la vie quotidienne des citoyens en facilitant les études, le commerce et les déplacements.
\end{enumerate}
\end{exoresolu}

\begin{exercice}[À vous de jouer]
À partir du schéma en cercles concentriques et de la carte de l'Afrique centrale étudiés dans cette leçon :
\begin{enumerate}
\item Citez les quatre niveaux d'intégration du Cameroun, du plus proche au plus large.
\item Expliquez en cinq lignes la différence entre la CEEAC et la CEMAC.
\item Donnez deux avantages concrets de l'intégration sous-régionale pour un commerçant camerounais.
\end{enumerate}
\end{exercice}

\begin{corrige}[Exercice]
\begin{enumerate}
\item Les quatre niveaux sont : la CEMAC (6 États), la CEEAC (11 États), l'Union Africaine (55 États) et l'ONU (193 États).
\item La CEEAC est une organisation politico-économique large de 11 États d'Afrique centrale, chargée de la coopération générale et de la paix. La CEMAC est une union économique et monétaire plus restreinte (6 États) qui gère la monnaie commune (le franc CFA), la banque centrale (BEAC) et le marché commun. Tous les pays de la CEMAC sont membres de la CEEAC, mais l'inverse est faux.
\item Deux avantages possibles : la libre circulation des marchandises sans droits de douane internes, ce qui réduit les coûts ; l'utilisation d'une monnaie commune (le franc CFA), qui évite les frais de change. On pouvait aussi citer le tarif extérieur commun et la sécurité renforcée sur les corridors commerciaux.
\end{enumerate}
\end{corrige}

\courssub{Bilan : acquis et obstacles de l'intégration}

L'intégration sous-régionale a produit des résultats tangibles pour les populations, mais elle reste inachevée. Il est indispensable, pour le Bac, de savoir nuancer son propos en distinguant les \textbf{acquis} des \textbf{obstacles} persistants.

\begin{propriete}[Principaux acquis de l'intégration (CEMAC / CEEAC / UA)]
\begin{itemize}
\item \textbf{Libre circulation des personnes} : suppression du visa pour les ressortissants de la CEMAC (carte d'identité biométrique suffisante) ; passeport CEEAC en cours de généralisation.
\item \textbf{Libre circulation des biens et marché commun} : suppression des droits de douane internes (CEMAC), mise en place du \cle{Tarif Extérieur Commun} (TEC) à 4 bandes (5\%, 10\%, 20\%, 30\%) pour protéger la production locale.
\item \textbf{Intégration monétaire et financière} : monnaie unique, le \cle{Franc CFA (XAF)}, gérée par la \cle{BEAC} (siège à Yaoundé) ; harmonisation du droit bancaire (COBAC) et des marchés financiers (BVMAC).
\item \textbf{Coopération sécuritaire} : création de la \cle{Force Multinationale Mixte (FMM)} contre Boko Haram (bassin du Lac Tchad) ; stratégie commune de lutte contre la piraterie dans le Golfe de Guinée (Yaoundé Code of Conduct).
\item \textbf{Poids diplomatique} : position commune aux Nations Unies et à l'UA ; candidature unique pour les postes internationaux ; Agenda 2063 de l'UA comme feuille de route continentale.
\end{itemize}
\end{propriete}

\begin{attention}[Obstacles majeurs à l'intégration (à connaître pour l'argumentation)]
\begin{itemize}
\item \textbf{Application insuffisante des textes} : harcèlement aux frontières (rackets, barrières illicites), non-reconnaissance mutuelle de certains diplômes ou titres de séjour, lenteur de la mise en œuvre du passeport CEEAC.
\item \textbf{Déficit d'infrastructures} : routes transfrontalières dégradées (ex. corridor Douala-Bangui-N'Djamena), absence de réseau ferroviaire régional, interconnexion électrique embryonnaire.
\item \textbf{Instabilités politiques et sécuritaires} : crises récurrentes en RCA, au Tchad, au Nord-Cameroun ; transitions politiques qui fragilisent la continuité des engagements communautaires.
\item \textbf{Chevauchement des appartenances} : certains pays (ex. RDC, Angola) appartiennent à plusieurs CER (CEEAC, SADC, COMESA), ce qui dilue la loyauté communautaire et complexifie l'harmonisation des politiques commerciales.
\item \textbf{Financement communautaire précaire} : la \cle{Prélèvement Communautaire de Solidarité (PCS)} (0,4\% sur les importations extra-CEMAC) est mal recouvrée ; arriérés de cotisations des États aux budgets de la CEEAC et de l'UA.
\item \textbf{Question monétaire} : débat persistant sur la réforme du Franc CFA (parité fixe avec l'euro, centralisation des réserves à 50\% en France) qui divise les opinions publiques et certains gouvernements.
\end{itemize}
\end{attention}

\begin{exemplebox}[Cas concret : le corridor Douala -- N'Djamena]
Ce corridor est la vitrine des acquis et des blocages.
\begin{itemize}
\item \textbf{Acquis} : monnaie unique, TEC, textes sur la libre circulation, existence de la FMM pour sécuriser l'axe.
\item \textbf{Réalité vécue par le transporteur} : multiples barrages de police/gendarmerie/douane (parfois 50+ sur le trajet), tracasseries financières, état de la route (bourbiers en saison des pluies), délais imprévisibles. L'intégration \emph{de droit} (textes) ne coïncide pas encore avec l'intégration \emph{de fait} (pratiques).
\end{itemize}
\end{exemplebox}

\courssub{TP guidé : cartographier l'intégration du Cameroun}

\begin{methode}[Réaliser un croquis d'intégration sous-régionale (type Bac)]
\begin{enumerate}
\item \textbf{Fond de carte} : tracer le contour simplifié du Cameroun et de ses 6 voisins immédiats (Nigeria, Tchad, RCA, Congo, Gabon, Guinée équatoriale). Positionner Yaoundé (capitale) et Douala (port principal).
\item \textbf{Zonage (coloration)} :
      \begin{itemize}
      \item Hachurer en \textbf{vert} les 5 pays membres de la CEMAC (hors Cameroun) : Tchad, RCA, Congo, Gabon, Guinée équatoriale.
      \item Hachurer en \textbf{bleu clair} les 5 pays membres de la CEEAC \emph{non membres de la CEMAC} : Angola, Burundi, RDC, Rwanda, Sao Tomé-et-Principe (les placer en périphérie ou par flèches s'ils sont hors cadre).
      \end{itemize}
\item \textbf{Symboles ponctuels} :
      \item $\bullet$ \textbf{Carré rouge} à Yaoundé : Siège de la BEAC (Banque centrale).
      \item $\bullet$ \textbf{Triangle bleu} à Bangui : Siège de la Commission CEMAC.
      \item $\bullet$ \textbf{Losange vert} à Libreville : Siège du Secrétariat général de la CEEAC.
      \item $\bullet$ \textbf{Étoile jaune} à Addis-Abeba (hors cadre, flèche) : Siège de l'UA.
\item \textbf{Flux (flèches proportionnelles)} :
      \item Flèches \textbf{oranges} (import/export) : Douala $\leftrightarrow$ N'Djamena / Bangui (corridor routier/fluvial).
      \item Flèches \textbf{violettes} (pétrole/bois) : Kribi / Douala $\rightarrow$ Monde ; Chad-Cameroon Pipeline (Kome $\rightarrow$ Kribi).
\item \textbf{Légende organisée} : structurer par thèmes (Zonage, Sièges institutionnels, Flux économiques, Capitales). Indiquer l'échelle (ex. 1 cm $\approx$ 300 km) et le Nord.
\item \textbf{Titre complet} : \emph{Le Cameroun dans les dynamiques d'intégration sous-régionale (CEMAC / CEEAC)}.
\end{enumerate}
\end{methode}

\begin{exercice}[TP d'application : Croquis commenté]
Réalisez le croquis décrit dans la méthode ci-dessus sur une feuille A4 (ou dans votre cahier). Puis, rédigez un commentaire structuré de 15 lignes maximum répondant à cette consigne : \emph{« Montrez que le Cameroun est un acteur central et un carrefour de l'intégration en Afrique centrale, tout en relevant les limites qui freinent l'aboutissement de cet espace communautaire. »}
\end{exercice}

\begin{corrige}[TP d'application -- Éléments attendus pour le commentaire]
\textbf{Introduction} : Localisation du Cameroun (carrefour Afrique centrale / Afrique de l'Ouest), appartenance multiple (CEMAC, CEEAC, UA, Commonwealth, Francophonie).
\textbf{I. Un acteur central (Atouts)} :
\begin{itemize}
\item Siège de la BEAC (Yaoundé) = cœur monétaire de la CEMAC.
\item Premier PIB de la CEMAC ($\approx$ 40\% du total), premier port (Douala) et premier corridor d'exportation pour les pays enclavés (Tchad, RCA).
\item Rôle diplomatique : médiations régionales, contribution majeure aux troupes de la FMM et aux contingents UA/ONU.
\end{itemize}
\textbf{II. Des limites structurelles (Freins)} :
\begin{itemize}
\item Infrastructures de transport insuffisantes (route unique Douala-Bangui-N'Djamena, pas de rail transfrontalier).
\item Persistance des barrières non tarifaires (rackets, lenteurs douanières) qui renchérissent le transport de 30 à 50\%.
\item Dépendance aux exportations brutes (bois, pétrole, coton) : faible intégration industrielle régionale (peu d'échanges de produits transformés intra-CEMAC).
\item Instabilités sécuritaires aux frontières (Extrême-Nord, Est, Nord-Ouest/Sud-Ouest) qui perturbent les flux.
\end{itemize}
\textbf{Conclusion} : Le Cameroun est le \emph{moteur} indispensable de l'intégration, mais l'effectivité de l'espace communautaire suppose un saut qualitatif : infrastructures lourdes, discipline aux frontières, diversification économique.
\end{corrige}

\begin{aretenir}[L'essentiel de la leçon -- Le Cameroun et les structures d'intégration]
L'\cle{intégration sous-régionale} est le regroupement d'États voisins qui mutualisent leurs ressources et harmonisent leurs politiques (économiques, monétaires, sécuritaires, culturelles). Le Cameroun s'y inscrit à plusieurs niveaux emboîtés : \textbf{CEMAC} (6 États, union monétaire, marché commun, BEAC à Yaoundé), \textbf{CEEAC} (11 États, coopération politique, paix, siège à Libreville), \textbf{Union Africaine} (55 États, Agenda 2063, siège à Addis-Abeba), \textbf{ONU}. Il cumule des appartenances originales (Commonwealth, Francophonie, CEN-SAD).
\textbf{Acquis concrets} : libre circulation (carte d'identité CEMAC), monnaie unique (FCFA), TEC, coopération sécuritaire (FMM).
\textbf{Obstacles majeurs} : non-application des textes aux frontières (rackets), déficit d'infrastructures, instabilités politiques, chevauchement des CER, financement communautaire fragile, débat sur le Franc CFA.
\textbf{Enjeu} : passer de l'intégration \emph{de droit} (textes signés) à l'intégration \emph{de fait} (réalité vécue par les citoyens et les entreprises) pour transformer le potentiel démographique et naturel en développement partagé.
\end{aretenir}

\coursec{Le Cameroun et l'Union Africaine (UA)}

Dans la section précédente, nous avons compris ce qu'est l'\cle{intégration sous-régionale} : le regroupement volontaire d'États voisins pour mettre en commun une partie de leurs moyens et de leurs décisions. Nous avons vu que l'intégration peut se vivre à plusieurs étages : le quartier, la sous-région, mais aussi le continent entier. C'est ce dernier étage que nous explorons maintenant. Avant d'être membre de la CEEAC ou de la CEMAC, le Cameroun est d'abord un État \cle{africain}, engagé dans la grande organisation qui rassemble tout le continent : l'\cle{Union Africaine}. Comprendre cette organisation, c'est comprendre le cadre le plus large dans lequel le Cameroun agit, coopère et défend ses intérêts.

\courssub{De l'OUA à l'UA : une histoire de l'unité africaine}

\textbf{L'intuition de départ.} Imaginez un village où chaque famille cultive seule son champ, garde seule sa récolte et règle seule ses conflits. Face à un voisin puissant ou à une sécheresse, chaque famille est fragile. Mais si toutes les familles se réunissent, décident ensemble et s'entraident, le village devient fort. C'est exactement cette idée qui a poussé les États africains, au sortir des indépendances, à se regrouper.

\textbf{La création de l'OUA.} Au début des années 1960, la plupart des colonies africaines viennent d'accéder à l'indépendance. Ces jeunes États se sentent faibles face aux anciennes puissances coloniales et aux grandes puissances de la guerre froide. Le 25 mai 1963, les dirigeants de 32 États africains indépendants se réunissent à Addis-Abeba, capitale de l'Éthiopie, et signent la charte créant l'\cle{Organisation de l'Unité Africaine} (OUA). Le \cle{Cameroun}, indépendant depuis 1960, fait partie des \cle{États fondateurs} : il signe la charte dès l'origine. Les objectifs de l'OUA sont alors simples : achever la décolonisation du continent, défendre la souveraineté des États membres et promouvoir l'unité africaine. L'OUA joue un rôle important dans la lutte contre l'apartheid en Afrique du Sud et la colonisation portugaise (Angola, Mozambique, Guinée-Bissau).

\textbf{La transformation en Union Africaine.} À la fin du XX\textsuperscript{e} siècle, l'OUA apparaît dépassée : elle défendait les États, mais restait silencieuse face aux conflits internes, aux coups d'État et aux violations des droits humains. Les dirigeants africains décident alors de la moderniser. L'acte constitutif de la nouvelle organisation est adopté à Lomé (Togo) en 2000, puis, au sommet de Durban (Afrique du Sud), en juillet 2002, l'OUA est officiellement remplacée par l'\cle{Union Africaine} (UA).

\begin{definition}[L'Union Africaine]
L'\cle{Union Africaine} (UA) est une organisation continentale qui regroupe les 55 États africains. Lancée en 2002 au sommet de Durban pour succéder à l'OUA (fondée en 1963), elle a pour but de promouvoir l'unité, la solidarité, l'intégration politique et économique, la paix et les droits humains en Afrique. Son siège est à Addis-Abeba (Éthiopie).
\end{definition}

\begin{attention}[Erreur fréquente à éviter]
Ne jamais écrire, dans une copie de Probatoire ou de Baccalauréat, que l'OUA existe encore ou que le Cameroun « est membre de l'OUA » aujourd'hui. L'OUA a été \textbf{remplacée} par l'Union Africaine en 2002. Dire « le Cameroun est membre fondateur de l'OUA » est exact pour 1963 ; dire « le Cameroun est membre de l'UA » est exact pour aujourd'hui. Confondre les deux coûte des points.
\end{attention}

\begin{savaistu}[Le 25 mai, Journée de l'Afrique]
Le 25 mai, date de la création de l'OUA à Addis-Abeba en 1963, est célébré chaque année dans tout le continent comme la \cle{Journée de l'Afrique}. Dans de nombreux établissements scolaires camerounais, des conférences, des expositions et des activités culturelles sont organisées ce jour-là pour célébrer l'unité africaine.
\end{savaistu}

\begin{popfigure}
\begin{center}
\begin{tikzpicture}[scale=1]
% Ligne du temps
\draw[line width=2pt,popblue] (0,0) -- (13.5,0);
\draw[-{Stealth[length=3mm]},line width=2pt,popblue] (13.5,0) -- (14.3,0);
% Jalons
\foreach \x/\year/\txt/\col in {1.2/1963/{Création de l'OUA\\à Addis-Abeba ;\\le Cameroun est\\membre fondateur}/popgreen,
5.2/2002/{Sommet de Durban :\\l'OUA devient\\l'Union Africaine}/poporange,
9.2/2019/{Entrée en vigueur\\de la ZLECAf (libre-échange\\continental)}/poppurple,
12.6/2063/{Horizon de\\l'Agenda 2063 :\\l'Afrique que\\nous voulons}/popgold}{
\fill[\col] (\x,0) circle (0.14);
\draw[\col,line width=1pt] (\x,0) -- (\x,0.7);
\node[above,align=center,font=\small\bfseries,text=\col] at (\x,0.75) {\year};
\node[below,align=center,font=\scriptsize,text=popdark] at (\x,-0.35) {\txt};
}
\end{tikzpicture}
\end{center}
\legende{Frise chronologique : de l'OUA (1963) à l'Agenda 2063, les grandes étapes de l'unité africaine.}
\end{popfigure}

\courssub{Les objectifs de l'Union Africaine}

L'acte constitutif de l'UA fixe des objectifs plus ambitieux que ceux de l'ancienne OUA. On peut les regrouper en quatre grandes missions :

\begin{enumerate}
\item \textbf{L'unité et la solidarité africaines} : renforcer les liens entre les États et les peuples du continent, défendre la souveraineté et l'intégrité territoriale des membres.
\item \textbf{L'intégration politique et économique} : accélérer l'union économique du continent (marché commun, monnaie, libre circulation), en s'appuyant sur les communautés sous-régionales comme la CEEAC et la CEMAC, considérées comme des « piliers » de l'intégration continentale.
\item \textbf{La paix et la sécurité} : prévenir et gérer les conflits, y compris par des missions de maintien de la paix. Grande nouveauté par rapport à l'OUA : l'UA peut \emph{intervenir} dans un État membre en cas de crimes graves (génocide, crimes de guerre, crimes contre l'humanité).
\item \textbf{La promotion des droits humains et de la démocratie} : encourager la bonne gouvernance, l'égalité hommes-femmes et le respect des droits des citoyens, notamment à travers la Charte africaine des droits de l'homme et des peuples.
\end{enumerate}

\begin{exemplebox}[L'UA en chiffres]
L'Union Africaine regroupe \textbf{55 États membres} — c'est-à-dire la quasi-totalité des États du continent — et plus de \textbf{1,3 milliard d'habitants}, soit environ 17 \% de la population mondiale. C'est la plus grande organisation régionale du monde par le nombre d'États membres. Le Cameroun, avec environ 28 millions d'habitants, représente à lui seul près de 2 \% de la population de l'Union.
\end{exemplebox}

\courssub{Les organes de l'Union Africaine}

Pour fonctionner, l'UA s'est dotée d'institutions qui ressemblent à celles d'un État : un sommet de dirigeants (le pouvoir de décision), une administration (la Commission), un parlement, etc. Retenons les cinq organes principaux :

\begin{itemize}
\item \textbf{La Conférence des chefs d'État et de gouvernement} : c'est l'organe suprême de l'Union. Elle se réunit au moins une fois par an et prend les grandes décisions (admission de membres, orientations générales). Le Président de la République du Cameroun y représente notre pays.
\item \textbf{Le Conseil exécutif} : il réunit les ministres des Affaires étrangères des États membres. Il prépare les sommets et veille à l'application des décisions de la Conférence.
\item \textbf{La Commission de l'UA} : c'est l'administration permanente de l'Union, son « gouvernement » au quotidien. Son siège est à \cle{Addis-Abeba}. Elle est dirigée par un président de la Commission, élu pour quatre ans.
\item \textbf{Le Parlement panafricain} : installé à Midrand (Afrique du Sud), il regroupe des délégués des parlements nationaux. Il donne des avis et fait des recommandations pour rapprocher l'Union des citoyens africains.
\item \textbf{Le Conseil de paix et de sécurité} (CPS) : véritable « conseil de sécurité africain », il prévient et gère les conflits, autorise les missions de paix et peut recommander des sanctions contre un État ou des individus.
\end{itemize}

\begin{popfigure}
\begin{center}
\begin{tikzpicture}[
font=\small,
box/.style={draw=popblue,fill=popblueL,rounded corners=2pt,align=center,text=popdark,minimum height=0.9cm,inner sep=4pt},
topbox/.style={draw=poporange,fill=poporangeL,rounded corners=2pt,align=center,text=popdark,font=\bfseries,minimum height=1cm,inner sep=4pt},
link/.style={draw=popmuted,line width=1pt}]
\node[topbox, align=center] (conf) at (0,0){Conférence des chefs d'État\\et de gouvernement (organe suprême)};
\node[box, align=center] (cexec) at (-4.6,-2.2){Conseil exécutif\\(ministres des Affaires\\étrangères)};
\node[box, align=center] (comm) at (-1.5,-2.2){Commission de l'UA\\(siège : Addis-Abeba)\\administration permanente};
\node[box, align=center] (parl) at (1.8,-2.2){Parlement\\panafricain\\(Midrand)};
\node[box, align=center] (cps) at (4.9,-2.2){Conseil de paix\\et de sécurité\\(gestion des conflits)};
\draw[link] (conf.south) -- (cexec.north);
\draw[link] (conf.south) -- (comm.north);
\draw[link] (conf.south) -- (parl.north);
\draw[link] (conf.south) -- (cps.north);
\end{tikzpicture}
\end{center}
\legende{Organigramme simplifié des principaux organes de l'Union Africaine. La Conférence des chefs d'État, organe suprême, chapeaute les autres institutions.}
\end{popfigure}

\begin{aretenir}[L'essentiel sur l'UA]
\begin{itemize}
\item 25 mai 1963 : création de l'\cle{OUA} à Addis-Abeba ; le Cameroun est \cle{membre fondateur}.
\item 2002 (sommet de Durban) : l'OUA devient l'\cle{Union Africaine}, 55 États membres.
\item Objectifs : unité, intégration politique et économique, paix et sécurité, droits humains.
\item Organes : Conférence des chefs d'État, Conseil exécutif, Commission (Addis-Abeba), Parlement panafricain, Conseil de paix et de sécurité.
\end{itemize}
\end{aretenir}

\courssub{La participation du Cameroun à l'Union Africaine}

Être membre d'une organisation, ce n'est pas seulement payer une contribution : c'est \emph{participer}. Le Cameroun s'implique dans la vie de l'UA de plusieurs manières concrètes.

\textbf{La présence aux sommets.} Le Président de la République, ou son représentant, participe aux sommets de la Conférence des chefs d'État. Le Cameroun y vote les décisions, signe les traités (comme celui de la ZLECAf) et défend ses positions, par exemple sur la sécurité dans le bassin du Lac Tchad ou dans le golfe de Guinée.

\textbf{La contribution aux opérations de paix.} Le Cameroun fournit des soldats, des policiers et des observateurs à des missions de paix africaines et onusiennes sur le continent (par exemple en République centrafricaine, voisine du Cameroun). Cette contribution renforce l'image du pays comme acteur de la stabilité en Afrique centrale.

\textbf{L'accueil de réunions africaines à Yaoundé.} La capitale camerounaise accueille régulièrement des conférences, des réunions d'experts et des sommets sous-régionaux ou continentaux (santé, sécurité, économie). Accueillir ces rencontres, c'est à la fois un honneur diplomatique et une manière de peser dans les décisions africaines.

\textbf{Des Camerounais dans les institutions africaines.} Des ressortissants camerounais occupent ou ont occupé des postes de responsabilité dans les organes de l'UA et dans d'autres institutions africaines, ce qui témoigne de l'influence du pays sur la scène continentale.

\courssub{Des réalisations utiles au quotidien : l'Agenda 2063 et la ZLECAf}

On entend parfois : « l'Union Africaine, c'est loin de nos vies ». C'est faux. Deux projets récents montrent que l'UA touche directement la vie des citoyens, y compris celle d'un jeune technicien camerounais.

\textbf{L'Agenda 2063.} Adopté en 2013 (cinquantenaire de l'OUA), l'\cle{Agenda 2063} est le plan de développement du continent pour les cinquante ans suivants. Sa devise : « l'Afrique que nous voulons ». Il vise une Afrique prospère, intégrée, pacifique et dirigée par ses propres citoyens, avec des projets concrets : réseaux de routes et de chemins de fer transafricains, énergie, industrialisation, éducation des jeunes.

\textbf{La ZLECAf.} La \cle{Zone de libre-échange continentale africaine} (ZLECAf), dont l'accord a été signé à Kigali (Rwanda) en 2018, est entrée en vigueur en 2019 et les échanges ont effectivement démarré en 2021. Elle crée un immense marché unique : les droits de douane entre pays africains sont progressivement supprimés pour faciliter les échanges. Pour un technicien camerounais, cela signifie par exemple qu'une entreprise de Douala pourra vendre plus facilement ses produits au Tchad, au Gabon ou au Kenya, et qu'un jeune entrepreneur pourra chercher des clients dans tout le continent. La ZLECAf est un sujet d'actualité très apprécié des correcteurs au Baccalauréat.

\begin{exoresolu}[Exercice type Bac : analyser l'engagement du Cameroun dans l'UA]
\textbf{Énoncé.} Un élève déclare : « L'Union Africaine ne sert à rien pour le Cameroun, c'est seulement des réunions de chefs d'État. » À l'aide de deux arguments précis et documentés, montrez les limites de cette affirmation, puis rédigez une conclusion de quatre à six lignes.
\tcblower
\textbf{Solution détaillée.}

\emph{Argument 1 — La paix et la sécurité.} L'UA, à travers son Conseil de paix et de sécurité, coordonne la lutte contre l'insécurité sur le continent. Le Cameroun, confronté à la menace de Boko Haram dans l'Extrême-Nord et à l'instabilité en République centrafricaine, bénéficie de la coopération africaine en matière de sécurité et participe lui-même aux opérations de paix. L'UA sert donc directement la sécurité des Camerounais.

\emph{Argument 2 — L'intégration économique.} Avec la ZLECAf (entrée en vigueur en 2019), l'UA construit un marché continental de plus de 1,3 milliard de consommateurs. Les entreprises camerounaises pourront exporter plus facilement, ce qui peut créer des emplois, notamment pour les jeunes diplômés techniques.

\emph{Conclusion rédigée.} L'affirmation de l'élève est donc excessive. Certes, l'UA connaît des limites réelles (décisions parfois peu appliquées, dépendance financière), mais elle rend des services concrets au Cameroun : coopération sécuritaire face à Boko Haram et ouverture d'un marché continental grâce à la ZLECAf. Loin d'être un simple club de chefs d'État, l'Union Africaine est un outil que le Cameroun doit utiliser davantage pour son développement. Reste à souhaiter que les décisions prises à Addis-Abeba soient mieux appliquées sur le terrain, au bénéfice des citoyens.
\end{exoresolu}

\courssub{Les limites de l'Union Africaine}

Être honnête intellectuellement, c'est aussi voir les difficultés. Trois limites majeures freinent l'action de l'UA :

\begin{enumerate}
\item \textbf{La dépendance financière} : une part importante du budget de l'Union, surtout pour les opérations de paix, provient de partenaires extérieurs (Union européenne, Nations unies, pays donateurs). Une organisation qui dépend de l'argent des autres peine à décider librement.
\item \textbf{Les conflits persistants} : malgré le Conseil de paix et de sécurité, plusieurs crises durent (Sahel, Soudan, est de la République démocratique du Congo). L'UA manque de moyens militaires et logistiques propres pour imposer la paix rapidement.
\item \textbf{La faible application des décisions} : beaucoup de traités et de décisions des sommets ne sont pas ratifiés ou appliqués par les États membres. La libre circulation des personnes, par exemple, reste très limitée en pratique.
\end{enumerate}

Ces limites ne signifient pas que l'UA est inutile : elles montrent que l'intégration est un \emph{processus} long, qui exige la volonté constante des États et des citoyens.

\courssub{Méthode : rédiger une conclusion en ECM}

Au Probatoire et au Baccalauréat, une dissertation ou une question de synthèse en ECM se termine toujours par une conclusion soignée. Voici la démarche attendue.

\begin{methode}[Construire une conclusion en trois temps]
\begin{enumerate}
\item \textbf{Rappeler la thèse} : reformulez en une phrase la position que vous avez défendue (sans copier mot pour mot l'introduction).
\item \textbf{Présenter deux arguments justifiés} : résumez vos deux meilleurs arguments, chacun appuyé sur un fait précis (date, texte, exemple camerounais ou africain).
\item \textbf{Ouvrir sur une perspective} : terminez par une ouverture — un souhait, une question, un élargissement — qui montre votre réflexion de citoyen (par exemple : « Reste à renforcer l'application des décisions de l'UA pour que l'intégration profite à tous les Africains »).
\end{enumerate}
\end{methode}

\begin{exercice}[Pour s'entraîner]
\begin{enumerate}
\item Complétez : l'OUA a été créée le \ldots{} à \ldots{} ; elle est devenue l'Union Africaine en \ldots{} au sommet de \ldots{}.
\item Citez trois organes de l'UA et donnez le rôle de chacun.
\item En cinq lignes, expliquez en quoi la ZLECAf peut être utile à un jeune technicien camerounais.
\end{enumerate}
\end{exercice}

\begin{corrige}[Exercice n°1]
\begin{enumerate}
\item L'OUA a été créée le \textbf{25 mai 1963} à \textbf{Addis-Abeba} (Éthiopie) ; elle est devenue l'Union Africaine en \textbf{2002} au sommet de \textbf{Durban} (Afrique du Sud).
\item Trois organes (au choix) : la \textbf{Conférence des chefs d'État} (organe suprême, prend les grandes décisions) ; la \textbf{Commission} (administration permanente, siège à Addis-Abeba) ; le \textbf{Conseil de paix et de sécurité} (prévention et gestion des conflits) ; le \textbf{Parlement panafricain} (avis et recommandations) ; le \textbf{Conseil exécutif} (ministres des Affaires étrangères, prépare les sommets).
\item Réponse attendue : la ZLECAf supprime progressivement les droits de douane entre pays africains. Un jeune technicien camerounais peut donc vendre ses produits ou services plus facilement dans d'autres pays africains, trouver des débouchés dans un marché de plus de 1,3 milliard de consommateurs et bénéficier des emplois créés par les entreprises qui exportent.
\end{enumerate}
\end{corrige}

Ainsi, le Cameroun, membre fondateur de l'OUA et membre actif de l'Union Africaine, est pleinement inséré dans la dynamique continentale. Mais l'UA est un cadre vaste, parfois éloigné des réalités quotidiennes. C'est pourquoi le Cameroun s'engage aussi, plus près de chez lui, dans sa sous-région : c'est l'objet de la section suivante, consacrée au Cameroun dans la CEEAC.

\coursec{Le Cameroun dans la CEEAC}

Après avoir étudié le rôle du Cameroun au sein de l'Union Africaine, organisation continentale qui regroupe l'ensemble des pays africains, nous descendons maintenant d'un niveau : celui de la \cle{sous-région}. En effet, entre le continent africain tout entier et l'État national, il existe des regroupements intermédiaires de pays voisins, appelés \cle{communautés économiques régionales}. Pour l'Afrique centrale, cette organisation s'appelle la CEEAC. Comprendre la place du Cameroun dans la CEEAC, c'est comprendre comment notre pays coopère chaque jour avec ses voisins les plus proches pour le commerce, la paix et la sécurité.

\courssub{Présentation générale de la CEEAC}

\textbf{Intuition.} Imaginez un quartier où onze familles décident de s'entraider : elles ouvrent leurs clôtures, organisent ensemble la sécurité de la rue et créent un marché commun. Chaque famille garde sa maison, mais toutes gagnent en sécurité et en prospérité. La CEEAC fonctionne exactement selon cette logique, à l'échelle de onze États d'Afrique centrale.

\begin{definition}[La CEEAC]
La \cle{CEEAC} (Communauté Économique des États de l'Afrique Centrale) est une organisation sous-régionale créée le \textbf{18 octobre 1983 à Libreville (Gabon)}, où se trouve son siège. Elle regroupe \textbf{11 États membres} et vise l'\cle{intégration} politique, économique et sécuritaire de l'Afrique centrale.
\end{definition}

Les onze États membres de la CEEAC sont : le \textbf{Cameroun}, le \textbf{Gabon}, le \textbf{Congo}, la \textbf{République Démocratique du Congo (RDC)}, le \textbf{Tchad}, la \textbf{République Centrafricaine (RCA)}, la \textbf{Guinée équatoriale}, \textbf{São Tomé-et-Príncipe}, l'\textbf{Angola}, le \textbf{Burundi} et le \textbf{Rwanda}.

\begin{exemplebox}[Une communauté de grande taille]
La CEEAC couvre un espace de 11 États réunissant plus de \textbf{190 millions d'habitants}. C'est presque sept fois la population du Cameroun. Un marché commun à cette échelle offre des débouchés considérables aux entreprises camerounaises : un fabricant de savon de Douala ou un atelier de transformation agroalimentaire de Yaoundé peut, en principe, vendre ses produits de Libreville à Kinshasa.
\end{exemplebox}

\begin{attention}[Erreur fréquente]
Beaucoup d'élèves croient que la CEEAC ne regroupe que les petits pays autour du Cameroun. C'est faux : la \textbf{RDC} et l'\textbf{Angola}, bien que situés plus au sud et comptant parmi les plus grands pays d'Afrique, sont membres de la CEEAC. De même, le Burundi et le Rwanda, en Afrique de l'Est, en font partie. Au Probatoire et au Baccalauréat technique, citez les 11 membres sans en oublier.
\end{attention}

\courssub{Les objectifs de la CEEAC}

Les objectifs de la CEEAC peuvent être retenus autour de trois axes complémentaires :

\begin{itemize}
\item \textbf{L'intégration économique} : promouvoir la coopération économique entre les États membres et réaliser à terme un \cle{marché commun}, c'est-à-dire un espace où les biens, les services, les capitaux et les personnes circulent librement.
\item \textbf{La paix et la sécurité} : prévenir et gérer les conflits en Afrique centrale, car sans paix, aucun développement économique n'est possible.
\item \textbf{La solidarité politique} : harmoniser les positions des États membres sur les grandes questions qui touchent la sous-région.
\end{itemize}

\begin{aretenir}[L'essentiel sur les objectifs]
La CEEAC poursuit un double pari : \textbf{la paix} comme condition, et \textbf{l'intégration économique} comme finalité. Un marché commun suppose que les frontières cessent d'être des murs pour devenir de simples repères.
\end{aretenir}

\courssub{Les organes de la CEEAC}

Pour fonctionner, la CEEAC s'est dotée d'institutions, à l'image d'un État qui possède un gouvernement, un parlement et des tribunaux.

\begin{center}
\begin{tabularx}{\linewidth}{|l|X|}
\hline
\rowcolor{popblueL} \textbf{Organe} & \textbf{Rôle principal} \\
\hline
Conférence des chefs d'État & Organe suprême : elle fixe les grandes orientations de la Communauté. \\
\hline
Conseil des ministres & Prépare les décisions et veille à leur application dans chaque domaine (économie, sécurité, transports). \\
\hline
Secrétariat général & Organe exécutif permanent, installé au siège de Libreville : il administre la Communauté au quotidien. \\
\hline
Cour de justice & Règle les litiges entre États membres et veille au respect des textes communautaires. \\
\hline
Parlement (siège à Malabo) & Organe consultatif représentant les peuples ; il siège à Malabo, en Guinée équatoriale. \\
\hline
\end{tabularx}
\end{center}

\begin{popfigure}
\begin{center}
\begin{tikzpicture}[scale=0.62, font=\small]
% Contour très simplifié de l'Afrique centrale (repères approximatifs)
% Cameroun
\fill[poporange] (2.2,4.6) -- (3.4,5.6) -- (4.6,5.4) -- (5.2,4.4) -- (4.6,3.2) -- (3.4,2.6) -- (2.4,3.2) -- (2.0,4.0) -- cycle;
% Tchad
\fill[popblueL] (3.6,5.8) -- (5.0,5.6) -- (5.4,7.4) -- (3.8,7.6) -- cycle;
% RCA
\fill[popblueL] (4.8,4.6) -- (6.6,4.8) -- (6.8,3.6) -- (5.0,3.4) -- cycle;
% Gabon
\fill[popblueL] (2.2,2.4) -- (3.2,2.4) -- (3.2,1.0) -- (2.2,1.0) -- cycle;
% Congo
\fill[popblueL] (3.4,2.4) -- (4.4,2.2) -- (4.2,0.6) -- (3.4,0.8) -- cycle;
% Guinée équatoriale (petit carré côtier)
\fill[popblueL] (1.9,2.6) rectangle (2.15,2.85);
% São Tomé (point insulaire)
\fill[popblueL] (1.2,1.9) circle (0.09);
% RDC
\fill[popblueL] (4.6,3.2) -- (6.8,3.4) -- (7.6,2.4) -- (7.2,0.2) -- (5.2,-0.2) -- (4.4,0.6) -- (4.6,2.0) -- cycle;
% Angola
\fill[popblueL] (4.6,-0.4) -- (7.0,-0.4) -- (6.8,-2.4) -- (4.8,-2.4) -- cycle;
% Rwanda
\fill[popblueL] (7.0,2.6) rectangle (7.3,2.9);
% Burundi
\fill[popblueL] (7.0,2.1) rectangle (7.3,2.4);
% Toponymes
\node[font=\small\bfseries, popdark] at (3.4,4.3) {CAMEROUN};
\node[font=\scriptsize] at (4.5,6.6) {Tchad};
\node[font=\scriptsize] at (5.8,4.1) {RCA};
\node[font=\scriptsize] at (2.7,1.7) {Gabon};
\node[font=\scriptsize] at (3.9,1.5) {Congo};
\node[font=\scriptsize] at (1.0,2.7) {G. équat.};
\node[font=\scriptsize] at (0.6,1.9) {STP};
\node[font=\scriptsize] at (6.0,1.6) {RDC};
\node[font=\scriptsize] at (5.8,-1.4) {Angola};
\node[font=\scriptsize] at (7.9,2.75) {Rwanda};
\node[font=\scriptsize] at (7.9,2.25) {Burundi};
% Villes
\fill[popdark] (2.5,3.4) circle (0.09); \node[font=\scriptsize, left] at (2.5,3.4) {Douala};
\fill[popdark] (3.1,3.7) circle (0.09); \node[font=\scriptsize, above] at (3.1,3.75) {Yaoundé};
\fill[popdark] (2.3,1.6) circle (0.09); \node[font=\scriptsize, right] at (2.35,1.55) {Libreville};
% Corridor Douala - Bangui
\draw[->, very thick, poppink] (2.5,3.4) .. controls (3.6,4.4) and (4.6,4.6) .. (5.6,4.2);
\fill[popdark] (5.6,4.2) circle (0.09); \node[font=\scriptsize, above] at (5.6,4.3) {Bangui};
% Nord
\draw[->, thick, popdark] (0.4,6.6) -- (0.4,7.4); \node[font=\scriptsize] at (0.4,7.7) {N};
% Légende
\fill[poporange] (0.2,-2.3) rectangle (0.6,-1.9); \node[font=\scriptsize, right] at (0.7,-2.1) {Cameroun (membre mis en évidence)};
\fill[popblueL] (4.6,-2.9) rectangle (5.0,-2.5); \node[font=\scriptsize, right] at (5.1,-2.7) {Autres membres de la CEEAC};
\draw[->, very thick, poppink] (0.2,-2.9) -- (1.0,-2.9); \node[font=\scriptsize, right] at (1.1,-2.9) {Corridor Douala--Bangui};
\end{tikzpicture}
\end{center}
\legende{Carte schématique de l'Afrique centrale : les 11 États membres de la CEEAC, le Cameroun en évidence, et le corridor Douala--Bangui. Représentation simplifiée, sans précision cartographique.}
\end{popfigure}

\begin{methode}[Construire une carte de localisation]
Pour réaliser une carte de localisation correcte au Baccalauréat technique, respecte cinq étapes :
\begin{enumerate}
\item Tracer le \textbf{fond de carte} : le contour simplifié de l'espace étudié (ici l'Afrique centrale).
\item Dessiner les \textbf{frontières} des États concernés et colorier les membres de l'organisation.
\item Placer les \textbf{toponymes} : noms des pays, capitales, villes importantes (Douala, Libreville, Bangui).
\item Ajouter la \textbf{légende} : chaque couleur ou symbole doit être expliqué.
\item Ne pas oublier le \textbf{titre} et l'\textbf{orientation} (flèche du nord), sans lesquels la carte est incomplète.
\end{enumerate}
\end{methode}

\courssub{Les actions concrètes de la CEEAC}

Une organisation ne se juge pas seulement à ses textes, mais à ses actes. Trois réalisations illustrent l'action de la CEEAC.

\textbf{1. La libre circulation des personnes.} En \textbf{2017}, la Conférence des chefs d'État a adopté une décision instaurant la libre circulation des personnes et des biens dans l'espace CEEAC. Son application est \emph{progressive} : certains pays l'appliquent déjà partiellement, d'autres tardent. Pour un jeune technicien camerounais, cela signifie à terme la possibilité d'aller travailler à Libreville, Brazzaville ou Kinshasa avec des formalités simplifiées.

\textbf{2. La FOMAC.} La \cle{FOMAC} (Force Multinationale de l'Afrique Centrale) est l'outil militaire de la CEEAC. Composée de contingents fournis par les États membres, elle peut être déployée pour des missions de paix, d'observation électorale ou de sécurisation. Elle montre que la sécurité de la sous-région est une affaire commune.

\textbf{3. La lutte contre Boko Haram.} Dans le \textbf{bassin du Lac Tchad}, la secte terroriste Boko Haram menace le Cameroun, le Tchad, le Nigeria et le Niger. La coopération militaire sous-régionale, soutenue par les mécanismes de la CEEAC et de l'Union Africaine, a permis des opérations conjointes qui ont considérablement affaibli le groupe. C'est la preuve qu'aucun État ne peut vaincre seul une menace transfrontalière.

\begin{savaistu}[Le corridor Douala--Bangui]
La République Centrafricaine est un pays \cle{enclavé} : elle n'a aucun accès à la mer. Pour exporter son bois, son coton ou ses diamants, elle emprunte le \textbf{corridor Douala--Bangui}, axe routier qui relie sa capitale au \textbf{port de Douala}. Ainsi, une grande partie du commerce extérieur centrafricain transite par le Cameroun. Le port de Douala est donc un poumon économique non seulement pour le Cameroun, mais pour toute la sous-région.
\end{savaistu}

\courssub{La place du Cameroun dans la CEEAC}

Le Cameroun n'est pas un membre ordinaire de la CEEAC : il y occupe une position de premier plan, pour trois raisons.

\begin{itemize}
\item \textbf{Un membre actif et contributeur majeur} : le Cameroun participe à toutes les instances de la Communauté et figure parmi les principaux contributeurs au budget de l'organisation. Cette contribution financière lui donne un poids politique réel dans les décisions.
\item \textbf{Une façade atlantique} : le Cameroun possède un littoral sur l'océan Atlantique et le principal port de la sous-région, Douala. C'est la porte d'entrée et de sortie des marchandises pour plusieurs pays enclavés.
\item \textbf{Une position stratégique de corridor} : l'axe \textbf{Douala--N'Djamena--Bangui} fait du Cameroun le trait d'union entre la mer et les pays de l'intérieur (Tchad, RCA). Qui contrôle les corridors contrôle les flux commerciaux de la sous-région.
\end{itemize}

\begin{exoresolu}[Le Cameroun, pivot de la CEEAC]
\textbf{Énoncé.} À partir de tes connaissances, dégage deux atouts qui font du Cameroun un acteur majeur de la CEEAC et montre comment chacun profite à la sous-région.
\tcblower
\textbf{Solution détaillée.} Premier atout : la \textbf{façade atlantique} et le port de Douala. Les pays enclavés comme la RCA et le Tchad exportent et importent leurs marchandises via Douala ; le Cameroun offre donc à la sous-région son accès à la mer, tout en tirant des revenus de ce transit. Deuxième atout : la \textbf{contribution financière et diplomatique}. En tant que contributeur majeur au budget de la CEEAC, le Cameroun finance le fonctionnement des institutions communes (Secrétariat général, FOMAC) et pèse dans les décisions de la Conférence des chefs d'État. Conclusion : le Cameroun est à la fois la \emph{porte maritime} et l'un des \emph{piliers financiers} de l'intégration en Afrique centrale.
\end{exoresolu}

\courssub{CEEAC et CEMAC : ne pas confondre}

Le Cameroun appartient à la fois à la CEEAC et à la CEMAC. Ces deux organisations se recouvrent partiellement, ce qui est une source fréquente de confusion. Le tableau suivant les compare ; la CEMAC sera étudiée en détail dans la section suivante.

\begin{center}
\begin{tabularx}{\linewidth}{|l|X|X|}
\hline
\rowcolor{popblueL} \textbf{Critère} & \textbf{CEEAC} & \textbf{CEMAC} \\
\hline
Date de création & 1983 (Libreville) & 1994 (N'Djamena), héritière de l'UDEAC (1964) \\
\hline
Nombre de membres & 11 États & 6 États \\
\hline
Membres & Cameroun, Gabon, Congo, RDC, Tchad, RCA, Guinée équatoriale, São Tomé-et-Príncipe, Angola, Burundi, Rwanda & Cameroun, Gabon, Congo, Tchad, RCA, Guinée équatoriale \\
\hline
Siège & Libreville (Gabon) & Bangui (RCA) \\
\hline
Missions principales & Intégration économique, paix et sécurité en Afrique centrale & Union économique et monétaire (monnaie commune : le franc CFA) \\
\hline
\end{tabularx}
\end{center}

\begin{attention}[Piège classique à l'examen]
Ne confonds jamais CEEAC et CEMAC : la \textbf{CEEAC} est plus large (11 membres) et s'occupe aussi de \textbf{paix et de sécurité} ; la \textbf{CEMAC} est plus restreinte (6 membres) et centrée sur l'\textbf{économie et la monnaie}. Écrire que la RDC est membre de la CEMAC, ou que la CEEAC émet le franc CFA, sont des erreurs éliminatoires dans une copie.
\end{attention}

\begin{aretenir}[À retenir de cette section]
La CEEAC, créée le 18 octobre 1983 à Libreville, regroupe 11 États d'Afrique centrale et plus de 190 millions d'habitants. Elle œuvre pour le marché commun, la paix et la sécurité (FOMAC, lutte contre Boko Haram, libre circulation depuis 2017). Le Cameroun y joue un rôle moteur grâce à sa contribution financière, à son port de Douala et au corridor Douala--N'Djamena--Bangui.
\end{aretenir}

\begin{exercice}[Pour s'entraîner]
\begin{enumerate}
\item Donne la date et le lieu de création de la CEEAC, puis cite cinq de ses États membres.
\item Explique en quelques lignes pourquoi le port de Douala est un atout stratégique pour la CEEAC.
\item Construis un tableau à deux colonnes distinguant la CEEAC et la CEMAC selon leur nombre de membres et leurs missions.
\end{enumerate}
\end{exercice}

\begin{corrige}[Exercice n°1]
1. La CEEAC a été créée le 18 octobre 1983 à Libreville (Gabon). Cinq membres : Cameroun, Gabon, Congo, Tchad, RDC (toute autre combinaison parmi les 11 est acceptée). 2. Le port de Douala est la principale ouverture maritime de la sous-région : les pays enclavés (RCA, Tchad) y transmettent leurs importations et exportations via les corridors Douala--Bangui et Douala--N'Djamena, ce qui en fait un levier d'intégration économique. 3. CEEAC : 11 membres, intégration économique, paix et sécurité. CEMAC : 6 membres, union économique et monétaire (franc CFA).
\end{corrige}

\coursec{Le Cameroun dans la CEMAC et les institutions communes}

Dans la section précédente, nous avons vu que le Cameroun est membre de la \cle{CEEAC}, une grande famille de onze États qui couvre toute l'Afrique centrale. Mais à l'intérieur de cette grande famille existe un cercle plus restreint et plus uni : six pays qui partagent non seulement des projets politiques, mais aussi \textbf{la même monnaie} et \textbf{les mêmes institutions financières}. Ce cercle, c'est la \cle{CEMAC}. C'est elle que nous étudions maintenant, car c'est l'organisation d'intégration qui touche le plus directement la vie quotidienne du commerçant de Douala, de l'apprenti mécanicien de Yaoundé ou du transporteur qui livre des marchandises jusqu'à Libreville.

\courssub{De l'UDEAC à la CEMAC : naissance d'une communauté}

\textbf{Une intuition pour commencer.} Imaginez six ateliers voisins dans un même quartier. Chacun fabrique ses produits, mais ils achètent leurs matières premières aux mêmes fournisseurs, se prêtent des outils et acceptent les mêmes bons d'achat. Au bout d'un moment, ils comprennent qu'ils gagneraient à fonctionner comme une seule grande entreprise : mêmes règles, même caisse commune, même monnaie interne. C'est exactement le chemin parcouru par les États d'Afrique centrale.

\textbf{Le raisonnement historique.} Dès 1964, les pays d'Afrique centrale nouvellement indépendants créent l'\cle{UDEAC} (Union Douanière et Économique de l'Afrique Centrale) par le traité de Brazzaville, pour faciliter les échanges entre eux. Mais l'UDEAC reste limitée : elle coordonne les droits de douane sans construire un véritable marché commun. En 1994, une grave crise économique (la dévaluation du franc CFA) pousse les chefs d'État à aller plus loin. Le \textbf{traité de N'Djamena} (capitale du Tchad), signé en 1994 et entré en vigueur en 1999, transforme l'UDEAC en une communauté beaucoup plus ambitieuse : la CEMAC.

\begin{definition}[La CEMAC]
La \cle{CEMAC} (Communauté Économique et Monétaire de l'Afrique Centrale) est une union économique et monétaire regroupant six États d'Afrique centrale, créée par le traité de N'Djamena signé en 1994 (entrée en vigueur 1999) en remplacement de l'UDEAC (1964). Son siège est à \textbf{Bangui} (République centrafricaine).
\end{definition}

\textbf{Les six membres} de la CEMAC sont : le \textbf{Cameroun}, le \textbf{Gabon}, le \textbf{Congo} (Brazzaville), le \textbf{Tchad}, la \textbf{République centrafricaine} (RCA) et la \textbf{Guinée équatoriale} (membre depuis 1983 pour l'UDEAC, partie prenante de la CEMAC à sa création).

\textbf{Les deux grands objectifs} de la CEMAC tiennent en deux mots :
\begin{itemize}
\item l'\cle{union économique} : construire un \textbf{marché commun} où les biens, les services, les capitaux et les personnes circulent librement d'un pays à l'autre, sans droits de douane intérieurs ;
\item l'\cle{union monétaire} : utiliser une \textbf{monnaie unique}, le franc CFA d'Afrique centrale, gérée par une banque centrale commune à tous les pays.
\end{itemize}

\begin{aretenir}[L'essentiel sur la CEMAC]
CEMAC = 6 États + 1 monnaie + 1 marché commun. Créée par le traité de N'Djamena (1994, entrée en vigueur 1999), elle succède à l'UDEAC (1964). Siège : Bangui (RCA). Objectifs : union économique et union monétaire.
\end{aretenir}

\courssub{Les institutions communes : la BEAC, la COBAC et la BDEAC}

\textbf{Pourquoi des institutions communes ?} Une monnaie partagée par six pays ne peut pas être imprimée par chacun dans son coin : il faudrait un gouverneur par pays, des décisions contradictoires, et la confiance s'effondrerait. Il faut donc des organes \textbf{supranationaux}, c'est-à-dire au-dessus des États, qui décident pour tous. La CEMAC en possède trois principaux dans le domaine financier.

\begin{definition}[La BEAC]
La \cle{BEAC} (Banque des États de l'Afrique Centrale) est la \textbf{banque centrale commune} aux six pays de la CEMAC. Elle \textbf{émet le franc CFA d'Afrique centrale} (code XAF) et \textbf{garantit la stabilité monétaire} de la zone : elle fixe les taux d'intérêt directeurs, contrôle la quantité de monnaie en circulation et veille à la solidité des banques. Son siège est à \textbf{Yaoundé} (Cameroun).
\end{definition}

\textbf{Comment cela fonctionne-t-il concrètement ?} Prenons une démarche d'observation simple. Regardez un billet de 10 000 francs CFA utilisé à Douala : il porte la mention « Banque des États de l'Afrique Centrale », et non « République du Cameroun ». Ce même billet est accepté à Libreville, à Brazzaville, à N'Djamena, à Bangui et à Malabo. Aucun État membre ne peut décider seul de dévaluer ou d'émettre plus de billets : c'est le conseil d'administration de la BEAC, où chaque pays est représenté, qui décide pour tous. C'est la preuve visible, dans votre portefeuille, que la souveraineté monétaire est \textbf{mutualisée}.

Autour de la BEAC gravitent deux autres institutions :
\begin{itemize}
\item la \cle{COBAC} (Commission Bancaire de l'Afrique Centrale) : c'est le \textbf{gendarme des banques}. Elle agrée les banques commerciales, surveille leur santé financière et peut sanctionner celles qui mettent en danger l'épargne des clients. Son siège est à \textbf{Libreville} (Gabon) ;
\item la \cle{BDEAC} (Banque de Développement des États de l'Afrique Centrale) : c'est la \textbf{banque des grands projets}. Elle finance les infrastructures et les investissements de développement dans la sous-région (routes, énergie, agriculture, industrie). Son siège est à \textbf{Brazzaville} (Congo).
\end{itemize}

\begin{center}
\begin{tabularx}{\linewidth}{|l|X|l|}
\hline
\rowcolor{popblueL} \textbf{Institution} & \textbf{Rôle principal} & \textbf{Siège} \\
\hline
CEMAC (secrétariat exécutif) & Pilotage politique et économique de la communauté & Bangui (RCA) \\
\hline
BEAC & Émet le franc CFA (XAF), garantit la stabilité monétaire & Yaoundé (Cameroun) \\
\hline
COBAC & Surveille et contrôle les banques commerciales & Libreville (Gabon) \\
\hline
BDEAC & Finance les projets de développement de la sous-région & Brazzaville (Congo) \\
\hline
\end{tabularx}
\end{center}

\begin{savaistu}[Yaoundé, cœur financier de l'Afrique centrale]
Le siège de la BEAC se trouve à \textbf{Yaoundé}. Cela signifie que les grandes décisions monétaires qui concernent six pays et des dizaines de millions d'habitants se prennent dans la capitale camerounaise. Le Cameroun est ainsi considéré comme le \textbf{cœur financier de l'Afrique centrale} : les banques, les assurances et les sièges régionaux de nombreuses entreprises s'y concentrent, ce qui renforce son influence économique dans toute la sous-région.
\end{savaistu}

\courssub{Cartographier la CEMAC et ses institutions}

La carte ci-dessous représente de manière schématique les six pays de la CEMAC et localise les sièges des trois grandes institutions communes. Observez la position centrale du Cameroun, qui partage des frontières avec quatre des cinq autres membres.

\begin{popfigure}
\begin{center}
\begin{tikzpicture}[scale=1.05, every node/.style={font=\small}]
% Tchad (nord)
\fill[popgoldL] (1.2,4.2) -- (4.2,4.2) -- (4.4,5.6) -- (1.0,5.6) -- cycle;
\node at (2.7,4.95) {\textbf{TCHAD}};
% RCA (est)
\fill[poppinkL] (4.5,2.4) -- (6.8,2.4) -- (6.8,4.1) -- (4.4,4.1) -- cycle;
\node at (5.6,3.3) {\textbf{RCA}};
% Cameroun (centre-ouest)
\fill[popgreenL] (1.0,2.2) -- (2.6,2.2) -- (2.6,4.1) -- (4.35,4.1) -- (4.45,2.3) -- (3.4,1.2) -- (2.2,0.6) -- (1.0,1.0) -- cycle;
\node at (3.0,3.0) {\textbf{CAMEROUN}};
% Congo (sud-est)
\fill[popblueL] (3.5,0.2) -- (5.6,0.2) -- (5.6,2.3) -- (4.5,2.3) -- (3.45,1.25) -- cycle;
\node at (4.7,1.2) {\textbf{CONGO}};
% Gabon (sud-ouest)
\fill[poporangeL] (1.2,-1.2) -- (3.3,-1.2) -- (3.3,0.5) -- (2.2,0.55) -- (1.0,0.95) -- cycle;
\node at (2.2,-0.45) {\textbf{GABON}};
% Guinée équatoriale (petit, côte)
\fill[poppurpleL] (0.6,-0.5) rectangle (1.1,0.2);
\node[font=\footnotesize, align=center] at (0.35,-0.85) {GUINÉE\\ÉQUATORIALE};
% Villes / sièges
\fill[popink] (2.45,2.55) circle (0.07); \node[anchor=east, font=\footnotesize] at (2.4,2.55) {\textbf{Yaoundé} (BEAC)};
\fill[popink] (5.0,3.85) circle (0.07); \node[anchor=south west, font=\footnotesize] at (5.0,3.9) {\textbf{Bangui} (CEMAC)};
\fill[popink] (4.35,0.55) circle (0.07); \node[anchor=west, font=\footnotesize] at (4.45,0.45) {\textbf{Brazzaville} (BDEAC)};
% Nord
\draw[-{Stealth}, thick, popdark] (7.2,5.0) -- (7.2,5.8); \node[font=\footnotesize] at (7.2,6.0) {N};
\end{tikzpicture}
\end{center}
\legende{Carte schématique des six États de la CEMAC et sièges des institutions communes : BEAC à Yaoundé, secrétariat exécutif de la CEMAC à Bangui, BDEAC à Brazzaville (croquis simplifié, sans précision cartographique).}
\end{popfigure}

\courssub{Le franc CFA (XAF) : une monnaie pour environ 60 millions de personnes}

\begin{exemplebox}[Une monnaie partagée par environ 60 millions de personnes]
Le franc CFA d'Afrique centrale (code international \cle{XAF}) est utilisé par environ \textbf{60 millions de personnes} réparties dans les six pays de la CEMAC (estimations récentes). Concrètement, un grossiste camerounais qui achète du bois au Gabon ou vend du ciment au Tchad règle ses factures dans la même monnaie que ses partenaires : pas de change à calculer, pas de risque de variation entre deux monnaies nationales. C'est un gain de temps, de sécurité et de coût pour tout le commerce sous-régional.
\end{exemplebox}

\begin{attention}[Erreur fréquente : il n'existe pas un, mais deux francs CFA]
Beaucoup d'élèves croient que le franc CFA du Cameroun est le même que celui du Sénégal ou de la Côte d'Ivoire. C'est \textbf{faux}. Il existe \textbf{deux zones CFA distinctes} : le franc CFA d'Afrique centrale (\cle{XAF}, zone CEMAC, émis par la BEAC) et le franc CFA d'Afrique de l'Ouest (\cle{XOF}, zone UEMOA, émis par la BCEAO à Dakar). Ils ont la même valeur faciale et la même parité avec l'euro, mais ils ne sont \textbf{pas interchangeables directement} : un billet XAF de Yaoundé n'est pas accepté tel quel dans les commerces de Dakar. Au Bac, ne confondez jamais les deux zones.
\end{attention}

\courssub{Les avantages concrets de la CEMAC pour les citoyens}

Au-delà des textes, que change la CEMAC dans la vie courante ? Trois acquis sont directement observables :
\begin{itemize}
\item \textbf{la même monnaie de Douala à Libreville} : le franc CFA (XAF) circule partout dans la zone, ce qui simplifie les échanges commerciaux et les voyages ;
\item \textbf{la libre circulation des biens et des personnes} : en principe, un citoyen d'un pays membre peut se déplacer, s'installer et travailler dans un autre pays de la CEMAC, et les marchandises franchissent les frontières sans droits de douane intérieurs ;
\item \textbf{le passeport CEMAC} : un passeport communautaire commun (biométrique depuis 2016) facilite les déplacements des ressortissants dans la sous-région et symbolise l'appartenance à un même espace.
\end{itemize}

\begin{exoresolu}[Appliquer la notion à une situation concrète]
\textbf{Énoncé.} Mme Etoundi, commerçante à Douala, exporte régulièrement des denrées alimentaires vers Libreville (Gabon) et Brazzaville (Congo). Son fournisseur gabonais lui propose d'être payé en francs CFA gabonais. Elle hésite, craignant de devoir changer cet argent en rentrant au Cameroun. À l'aide de vos connaissances sur la CEMAC, rassurez Mme Etoundi en deux arguments.
\tcblower
\textbf{Solution détaillée.}
\begin{enumerate}
\item \textbf{Argument monétaire} : le Gabon et le Cameroun appartiennent tous deux à la CEMAC et utilisent la même monnaie, le franc CFA d'Afrique centrale (XAF), émise par la BEAC. Les francs reçus à Libreville ont exactement la même valeur et le même usage à Douala : il n'y a \textbf{aucune opération de change} à effectuer.
\item \textbf{Argument économique} : grâce à l'union économique de la CEMAC (marché commun), les marchandises circulent entre les pays membres sans droits de douane intérieurs, ce qui réduit les coûts de son commerce et sécurise ses échanges.
\end{enumerate}
\textbf{Conclusion} : Mme Etoundi peut accepter le paiement sans crainte ; la monnaie gabonaise et la monnaie camerounaise sont une seule et même monnaie, le XAF.
\end{exoresolu}

\courssub{Le poids du Cameroun : première économie de la CEMAC}

Tous les membres de la CEMAC ne pèsent pas le même poids. Grâce à la diversité de son économie (agriculture, industrie, commerce, pétrole, services), le \textbf{Cameroun est la première économie de la zone} : il représente à lui seul environ \textbf{40 à 45 \% du PIB de la CEMAC}. Le diagramme ci-dessous illustre la part approximative de chaque pays dans le produit intérieur brut total de la communauté.

\begin{popfigure}
\begin{center}
\begin{tikzpicture}
\begin{axis}[
    ybar, bar width=16pt,
    width=0.95\linewidth, height=6.2cm,
    ymin=0, ymax=55,
    ylabel={Part du PIB de la zone (en \%)},
    symbolic x coords={Cameroun, Gabon, Congo, Tchad, Guinée éq., RCA},
    xtick=data,
    x tick label style={font=\footnotesize, rotate=15, anchor=east},
    nodes near coords, nodes near coords style={font=\footnotesize},
    every axis plot/.append style={fill=popblue, draw=popdark},
]
\addplot coordinates {(Cameroun,43) (Gabon,20) (Congo,15) (Tchad,12) (Guinée éq.,8) (RCA,2)};
\end{axis}
\end{tikzpicture}
\end{center}
\legende{Part approximative de chaque État dans le PIB total de la CEMAC (valeurs arrondies, ordres de grandeur indicatifs). Le Cameroun domine nettement la zone avec environ 43 \% du PIB communautaire.}
\end{popfigure}

Ce poids économique a deux conséquences importantes. D'une part, il donne au Cameroun une \textbf{responsabilité de locomotive} : quand l'économie camerounaise va bien, elle entraîne toute la sous-région ; quand elle ralentit, ses voisins en souffrent. D'autre part, il explique pourquoi Yaoundé accueille le siège de la BEAC et pourquoi le Cameroun est le premier partenaire commercial de plusieurs de ses voisins, notamment le Tchad et la RCA, pays enclavés qui transitent par le port de Douala.

\begin{aretenir}[À retenir pour le Bac]
\begin{itemize}
\item La CEMAC (traité de N'Djamena 1994, entrée en vigueur 1999) succède à l'UDEAC (1964) ; siège à Bangui ; 6 membres : Cameroun, Gabon, Congo, Tchad, RCA, Guinée équatoriale.
\item Double objectif : union économique (marché commun) et union monétaire (monnaie unique XAF).
\item Institutions : BEAC (Yaoundé, émet le XAF), COBAC (Libreville, surveillance des banques), BDEAC (Brazzaville, financement du développement).
\item Le Cameroun est la première économie de la zone (environ 40 à 45 \% du PIB) et le cœur financier de l'Afrique centrale.
\item Piège classique : le XAF (Afrique centrale) n'est pas le XOF (Afrique de l'Ouest).
\end{itemize}
\end{aretenir}

\begin{exercice}[Pour s'entraîner]
\begin{enumerate}
\item Définis la CEMAC en citant sa date de création, son siège et ses six États membres.
\item Explique en quelques lignes le rôle de la BEAC et précise où se trouve son siège.
\item Ton voisin affirme : « Avec mes francs CFA camerounais, je peux payer directement mes achats à Dakar. » Corrige son erreur en t'appuyant sur la distinction entre XAF et XOF.
\item Rédige un court paragraphe argumenté (8 à 10 lignes) montrant en quoi le Cameroun est la « locomotive » économique de la CEMAC.
\end{enumerate}
\end{exercice}

\begin{corrige}[Exercice : pour s'entraîner]
\begin{enumerate}
\item La CEMAC (Communauté Économique et Monétaire de l'Afrique Centrale) est une union économique et monétaire créée par le traité de N'Djamena signé en 1994 (entrée en vigueur 1999), en remplacement de l'UDEAC. Son siège est à Bangui (RCA). Ses six membres sont le Cameroun, le Gabon, le Congo, le Tchad, la République centrafricaine et la Guinée équatoriale.
\item La BEAC est la banque centrale commune aux six pays de la CEMAC. Elle émet le franc CFA d'Afrique centrale (XAF), fixe la politique monétaire et garantit la stabilité de la monnaie. Son siège est à Yaoundé, au Cameroun.
\item L'affirmation est fausse : le franc CFA du Cameroun est le XAF (Afrique centrale, émis par la BEAC), tandis que le Sénégal utilise le XOF (Afrique de l'Ouest, émis par la BCEAO). Ce sont deux monnaies distinctes, non interchangeables directement : il faut passer par une opération de change.
\item Réponse attendue (éléments) : le Cameroun représente environ 40 à 45 \% du PIB de la CEMAC, ce qui en fait la première économie de la zone ; son économie est diversifiée (agriculture, industrie, pétrole, services) ; Yaoundé abrite le siège de la BEAC, ce qui en fait le cœur financier de l'Afrique centrale ; enfin, le port de Douala sert de débouché maritime aux pays enclavés (Tchad, RCA), ce qui renforce le rôle moteur du Cameroun dans le commerce sous-régional.
\end{enumerate}
\end{corrige}

Ainsi, la CEMAC fait du Cameroun bien plus qu'un simple membre : elle en fait le centre de gravité économique et financier de l'Afrique centrale. Mais cette intégration, si prometteuse soit-elle, connaît aussi des limites : c'est ce que nous examinerons dans le bilan de la leçon.

\coursec{Comprendre l'intégration sous-régionale}

\courssub{Définitions et enjeux}

L'\cle{intégration sous-régionale} est un processus par lequel des États géographiquement proches décident de mettre en commun tout ou partie de leurs souverainetés (économique, monétaire, sécuritaire, politique) pour accélérer leur développement et peser davantage sur la scène internationale. En Afrique centrale, ce processus repose sur deux piliers imbriqués : la \cle{CEEAC} (Communauté Économique des États de l'Afrique Centrale) pour la coopération large, et la \cle{CEMAC} (Communauté Économique et Monétaire de l'Afrique Centrale) pour l'union monétaire et économique approfondie.

\begin{definition}[Niveaux d'intégration économique (rappel)]
\begin{enumerate}
\item \textbf{Zone de libre-échange} : suppression des droits de douane entre membres (ex. : ZLECAf).
\item \textbf{Union douanière} : zone de libre-échange + tarif extérieur commun (TEC) vis-à-vis des tiers.
\item \textbf{Marché commun} : union douanière + libre circulation des facteurs de production (travail, capital).
\item \textbf{Union économique et monétaire} : marché commun + monnaie unique + politique économique harmonisée (ex. : CEMAC, UEMOA).
\end{enumerate}
\end{definition}

\courssub{Les motivations fondamentales}

\begin{itemize}
\item \textbf{Le marché élargi} : les États d'Afrique centrale sont de petits marchés pris individuellement ; leur regroupement offre un débouché de près de 200 millions de consommateurs (CEEAC).
\item \textbf{L'enclavement} : le Tchad, la RCA sont des pays enclavés ; l'intégration leur garantit un accès à la mer (corridors vers Douala, Pointe-Noire, Libreville).
\item \textbf{La paix et la sécurité} : les conflits ignorent les frontières (Boko Haram, groupes armés en RCA, Est de la RDC) ; la mutualisation des moyens (FMM, forces de la CEEAC) est une nécessité.
\item \textbf{Le poids diplomatique} : parler d'une seule voix à l'UA, à l'ONU, dans les négociations commerciales (OMC, Accords UE-ACP/UE-Afrique).
\end{itemize}

\begin{savaistu}[Le Traité d'Abuja (1991) et l'Agenda 2063]
Le \textbf{Traité d'Abuja} (Nigeria, 1991) a fixé le cap : créer une \textbf{Communauté Économique Africaine (CEA)} en six étapes sur 34 ans, en s'appuyant sur les Communautés Économiques Régionales (CER) comme la CEEAC. L'\textbf{Agenda 2063} de l'UA (adopté en 2013) réaffirme cette vision : « L'Afrique que nous voulons » — intégrée, prospère, pacifique, pilotée par ses propres citoyens.
\end{savaistu}

\coursec{Le Cameroun et l'Union Africaine (UA)}

\courssub{De l'OUA à l'UA : le Cameroun membre fondateur}

Le Cameroun a adhéré à l'\textbf{Organisation de l'Unité Africaine (OUA)} dès sa création en 1963 (Sommet d'Addis-Abeba). Il a soutenu la transformation en \textbf{Union Africaine (UA)} en 2002 (Sommet de Durban), pour doter le continent d'institutions plus supranationales (Commission, Parlement panafricain, Cour de justice) et d'une vision de développement partagé (NEPAD, puis Agenda 2063).

\courssub{Les organes de l'UA et la participation camerounaise}

\begin{center}
\begin{tabularx}{\linewidth}{|l|X|X|}\hline
\rowcolor{popblueL}
\textbf{Organe} & \textbf{Rôle} & \textbf{Participation / Intérêt pour le Cameroun} \\ \hline
Conférence des Chefs d'État & Organe suprême, définit la politique & Le Président de la République y défend les positions nationales (sécurité, climat, ZLECAf). \\ \hline
Conseil exécutif (Ministres AE) & Prépare les sommets, suit l'exécution & Le Minrexeau y coordonne la diplomatie multilatérale. \\ \hline
Commission de l'UA (Addis-Abeba) & Exécutif permanent, met en œuvre les décisions & Des cadres camerounais y servent (ex. : Commissaires, directeurs). \\ \hline
Parlement panafricain (Midrand, Afrique du Sud) & Organe consultatif, représente les peuples & Députés camerounais élus y siègent ; rôle de contrôle politique. \\ \hline
Cour africaine des droits de l'homme et des peuples (Arusha) & Justice supranationale & Le Cameroun a ratifié la Charte africaine ; justiciable pour les citoyens. \\ \hline
Conseil de paix et de sécurité (CPS) & Prévention, gestion, résolution des conflits & Cameroun membre élu à plusieurs mandats ; engagement troupes (MINUSCA, FMM). \\ \hline
\end{tabularx}
\end{center}

\courssub{Engagements concrets du Cameroun au sein de l'UA}

\begin{itemize}
\item \textbf{Maintien de la paix} : contribution historique aux missions de l'UA/ONU (MISCA puis MINUSCA en RCA, MISMA puis MINUSMA au Mali, FMM au Lac Tchad). Le Cameroun est l'un des premiers fournisseurs de troupes de police et militaires du continent.
\item \textbf{Agenda 2063} : alignement de la \cle{Stratégie Nationale de Développement 2020-2030 (SND30)} sur les aspirations de l'Agenda 2063 (industrialisation, infrastructures, capital humain, gouvernance).
\item \textbf{ZLECAf} : ratification rapide (2019), mise en œuvre du commerce préférentiel (démarrage officiel 1er janvier 2021), désignation de Douala comme hub logistique potentiel.
\item \textbf{Siège d'institutions} : le Cameroun abrite le \textbf{Siège de la BEAC} (Yaoundé), institution spécialisée de l'UA pour la zone CEMAC, et accueille régulièrement des sommets (ex. Sommet UA-CEEAC 2022 sur la sécurité maritime).
\end{itemize}

\begin{attention}[Ne pas confondre UA et CEEAC/CEMAC]
L'\textbf{UA} est l'organisation \textbf{continentale} (55 États). La \textbf{CEEAC} et la \textbf{CEMAC} sont des \textbf{Communautés Économiques Régionales (CER)} reconnues par l'UA comme piliers de l'intégration en Afrique centrale. Le Cameroun y siège simultanément.
\end{attention}

\coursec{Le Cameroun dans la CEEAC}

\courssub{Genèse, membres et objectifs}

La \cle{CEEAC} a été créée par le Traité de Libreville (18 octobre 1983), réactivée en 1999 (Sommet de Malabo). Elle compte \textbf{11 États membres} : les 6 de la CEMAC (Cameroun, Tchad, RCA, Gabon, Congo, Guinée équatoriale) + \textbf{RDC, Sao Tomé-et-Principe, Burundi, Rwanda, Angola}. Son siège est à \textbf{Libreville} (Gabon).

\begin{propriete}[Objectifs prioritaires de la CEEAC (Traité révisé 2009)]
\begin{enumerate}
\item Parvenir à l'union douanière puis au marché commun.
\item Assurer la libre circulation des personnes, des biens, des services et des capitaux.
\item Harmoniser les politiques sectorielles (transports, énergie, agriculture, industrie, éducation).
\item Promouvoir la paix, la sécurité et la stabilité (Mécanisme de prévention, de gestion et de règlement des conflits — MPGR).
\end{enumerate}
\end{propriete}

\courssub{Institutions clés et apport du Cameroun}

\begin{itemize}
\item \textbf{Conférence des Chefs d'État} : le Président Paul Biya y impulse les dossiers sécurité (Lac Tchad, RCA) et infrastructures (corridors, énergie).
\item \textbf{Secrétariat Général (Libreville)} : assure l'administration courante ; des experts camerounais y occupent des postes de direction.
\item \textbf{Parlement de la CEEAC (Malabo)} : députés camerounais y siègent ; rôle de contrôle et d'harmonisation législative.
\item \textbf{Cour de justice de la CEEAC (N'Djamena)} : règlement des différends entre États ou entre États et citoyens.
\item \textbf{Forces de la CEEAC (FOMAC)} : force en attente pour le maintien de la paix ; le Cameroun fournit des contingents et l'appui logistique (port de Douala, aéroports).
\end{itemize}

\courssub{Projets structurants portés par le Cameroun dans la CEEAC}

\begin{itemize}
\item \textbf{Corridors multimodaux} : Douala--Bangui (RCA), Douala--N'Djamena (Tchad), Kribi--Bangui (nouveau port en eau profonde), projet routier/ferroviaire vers la RDC (via RCA ou Congo).
\item \textbf{Énergie} : interconnexion électrique Cameroun--Tchad (ligne Lom Pangar--N'Djamena), projet Grand Inga (RDC) -- Afrique centrale (le Cameroun comme plaque tournante).
\item \textbf{Navigation fluviale} : fleuves Oubangui, Sangha, Congo -- le Cameroun (via le port de Douala et la façade maritime) est la porte de sortie naturelle.
\end{itemize}

\begin{exemplebox}[La CEEAC, cadre de la diplomatie humanitaire]
Lors de la crise centrafricaine (2013-2014), la CEEAC a déployé la \textbf{MISCA} (Mission internationale de soutien à la Centrafrique sous conduite africaine), précurseur de la MINUSCA (ONU). Le Cameroun a accueilli des centaines de milliers de réfugiés centrafricains (Est, Adamaoua, Nord) et servi de base logistique arrière pour la force africaine puis onusienne. C'est l'intégration solidaire dans l'épreuve.
\end{exemplebox}

\coursec{Le Cameroun dans la CEMAC et les institutions communes}

\courssub{Spécificité de la CEMAC : l'union monétaire et économique}

La \cle{CEMAC} (Traité de 1994, révisé 2009) regroupe les \textbf{6 États} partageant le franc CFA : Cameroun, Tchad, RCA, Gabon, Congo, Guinée équatoriale. Elle va plus loin que la CEEAC : \textbf{monnaie unique, banque centrale commune, surveillance bancaire commune, parlement communautaire, cour de justice commune}. Son siège institutionnel est à \textbf{Bangui} (RCA), mais les institutions techniques sont réparties.

\courssub{Les institutions communes et la place du Cameroun}

\begin{center}
\begin{tabularx}{\linewidth}{|l|X|l|}\hline
\rowcolor{popblueL}
\textbf{Institution} & \textbf{Mission} & \textbf{Siège / Lien Cameroun} \\ \hline
Conférence des Chefs d'État & Orientation suprême & Tournante ; Cameroun hôte régulier. \\ \hline
Conseil des Ministres (Éco/Fin) & Décisions économiques & Minfi Cameroun y siège ; poids du 1er PIB zone. \\ \hline
\cle{Commission de la CEMAC} & Exécutif communautaire, mise en œuvre & Siège : \textbf{Bangui}. Secrétaire Général souvent camerounais. \\ \hline
\cle{Parlement communautaire} & Contrôle, avis consultatif & Siège : \textbf{Malabo} (Guinée équatoriale). Députés camerounais élus. \\ \hline
\cle{Cour de justice} & Contentieux communautaire & Siège : \textbf{N'Djamena} (Tchad). Juges camerounais. \\ \hline
\cle{BEAC (Banque des États de l'Afrique Centrale)} & Banque centrale, émission monnaie, politique monétaire & \textbf{Siège : Yaoundé (Cameroun)}. Gouverneur souvent camerounais. \\ \hline
\cle{COBAC (Commission Bancaire de l'Afrique Centrale)} & Supervision prudentielle banques & Siège : \textbf{Libreville} (Gabon). Experts camerounais. \\ \hline
\cle{BDEAC (Banque de Développement des États de l'Afrique Centrale)} & Financement projets d'intégration & Siège : \textbf{Brazzaville} (Congo). Cameroun actionnaire majeur, projets financés (routes, énergie). \\ \hline
\end{tabularx}
\end{center}

\courssub{La monnaie commune : le franc CFA (CEMAC)}

\begin{definition}[Franc CFA - Coopération Financière en Afrique Centrale]
Monnaie commune aux 6 États de la CEMAC. \textbf{Parité fixe} : 1 EUR = 655,957 XAF (garantie par le Trésor français via un compte d'opérations). \textbf{Convertibilité} : illimitée vers l'euro. \textbf{Indépendance monétaire} : la BEAC définit la politique monétaire (taux directeurs, réserves de change) pour la zone. \textbf{Stabilité} : inflation cible \textasciitilde 3\% (souvent respectée hors chocs pétroliers/COVID).
\end{definition}

\begin{exemplebox}[Avantage concret pour l'entrepreneur technique]
Un soudeur à Douala qui fabrique des portails pour un client à N'Djamena (Tchad) facture en francs CFA. Il est payé par virement BEAC en francs CFA. Aucun risque de change, aucun frais de conversion, délai bancaire standard. C'est la \cle{réduction du coût de transaction} que permet l'union monétaire.
\end{exemplebox}

\courssub{Libre circulation et droit d'établissement : textes et réalité}

\begin{itemize}
\item \textbf{Textes} : Convention 1983 (libre circulation), Convention 1990 (droit d'établissement), Supplément 2013 (biométrie, passeport CEMAC).
\item \textbf{Réalité} : un citoyen CEMAC peut entrer dans un autre pays membre avec sa \textbf{CNI biométrique} (ou passeport CEMAC) sans visa pour 90 jours renouvelables. Le droit d'établissement (travailler, créer entreprise) est acquis en droit, mais freiné en pratique par des tracasseries administratives (carte de séjour, permis de travail encore demandés illégalement).
\end{itemize}

\begin{attention}[Piège Bac : « Libre circulation = absence de contrôle »]
La libre circulation ne signifie pas \emph{absence de contrôle aux frontières}. Les États gardent le droit de contrôler l'identité (sécurité, santé publique). Ce qui est interdit, c'est le \textbf{visa obligatoire} et les \textbf{barrières discriminatoires} (rackets, refus d'entrée aux titulaires de CNI valide).
\end{attention}

\coursec{Bilan : acquis et obstacles de l'intégration}

Dans les sections précédentes, nous avons vu que le Cameroun est membre de la CEMAC et qu'il participe à des institutions communes : la \cle{BEAC}, la monnaie commune (le \cle{franc CFA}), la \cle{Commission de la CEMAC}, le \cle{Parlement communautaire}, la \cle{Cour de justice}, la \cle{COBAC}, la \cle{BDEAC}. Mais ces institutions donnent-elles des résultats concrets dans la vie des populations ? L'intégration sous-régionale est-elle un succès ou un échec ? C'est à cette question de \cle{bilan} que nous répondons maintenant. Un bon bilan, comme au Bac, n'est jamais tout noir ni tout blanc : il pèse les \cle{acquis} (ce qui a réussi) et les \cle{obstacles} (ce qui bloque encore), puis ouvre sur les \cle{perspectives}.

\courssub{Les acquis de l'intégration sous-régionale}

\courssub{Une monnaie commune stable}

Le premier acquis, le plus visible, est la \cle{monnaie commune} : le franc CFA d'Afrique centrale, utilisé par les six pays de la CEMAC. Grâce à la BEAC (siège à Yaoundé), cette monnaie est \textbf{stable} : sa valeur est garantie par la parité fixe avec l'euro, elle est librement convertible, et l'inflation reste maîtrisée (objectif \textasciitilde 3\%). Pour un commerçant ou un artisan camerounais qui achète des intrants à Libreville ou vend des produits à N'Djamena, c'est un avantage énorme : il facture et est payé dans la même monnaie, sans risque de change ni frais de conversion.

\begin{exemplebox}[La monnaie commune facilite les échanges et l'investissement]
Mama Béa, commerçante à Douala, achète régulièrement du poisson fumé au Gabon et du tissu au Congo. Elle paie en francs CFA, la même monnaie que ses fournisseurs. Aucun bureau de change, aucun frais de conversion, aucune mauvaise surprise liée à la variation des monnaies. De même, une PME camerounaise de BTP qui remporte un marché à Bata (Guinée équatoriale) n'a pas de risque de change pour rapatrier ses bénéfices. C'est l'avantage concret de l'intégration monétaire pour l'activité économique.
\end{exemplebox}

\courssub{La libre circulation partielle des personnes}

Le deuxième acquis est la \cle{libre circulation} des personnes dans l'espace CEMAC. Elle reste \textbf{partielle} (des contrôles d'identité et de sécurité existent), mais elle est réelle : les citoyens de la zone peuvent voyager plus facilement d'un pays à l'autre.

\begin{savaistu}[Voyager avec sa seule carte d'identité]
Dans le cadre des accords CEMAC (Convention 1983, Supplément 2013), un Camerounais muni de sa \textbf{carte nationale d'identité biométrique} (ou du passeport CEMAC) peut entrer au Tchad, en République centrafricaine, au Gabon, au Congo ou en Guinée équatoriale sans visa pour un séjour de 90 jours renouvelable. C'est un droit concret, souvent ignoré des populations et parfois contesté par des agents zélés aux frontières, qui montre que l'intégration n'est pas qu'une affaire de chefs d'État : elle touche la vie quotidienne de chacun.
\end{savaistu}

\courssub{La coopération sécuritaire}

Face aux menaces qui ignorent les frontières, les États de la sous-région ont compris qu'aucun pays ne peut se défendre seul. Depuis 2014, le groupe terroriste \textbf{Boko Haram} (et sa faction ISWAP) attaque l'Extrême-Nord du Cameroun, le Nigeria, le Niger et le Tchad. La réponse a été commune : la \cle{Force multinationale mixte (FMM)}, créée en 1994 (réactivée/renforcée par l'UA en 2015), qui réunit des soldats du Cameroun, du Tchad, du Niger, du Nigeria et du Bénin. Elle coordonne les opérations militaires, partage les renseignements et mène des actions conjointes. Cette coopération a permis de réduire fortement les capacités du groupe terroriste dans le bassin du Lac Tchad. C'est la preuve que l'union fait la force, y compris dans le domaine de la sécurité.

\courssub{Les infrastructures transnationales}

L'intégration se construit aussi avec du bitume, des rails et des ponts. Plusieurs \cle{infrastructures transnationales} relient le Cameroun à ses voisins :
\begin{itemize}
\item le \textbf{corridor Douala--N'Djamena}, axe routier (N1, N3) et ferroviaire (Camrail jusqu'à Ngaoundéré, puis route) vital pour le Tchad, pays enclavé dont une grande partie des importations et exportations transite par le port de Douala ;
\item le \textbf{corridor Douala--Bangui} (route N1/N4, projet ferroviaire), qui joue le même rôle pour la République centrafricaine ;
\item les \textbf{ponts frontaliers}, comme le pont sur la Logone entre Yagoua (Cameroun) et Bongor (Tchad), inauguré en 2023, ou le projet de pont sur le fleuve Ntem entre le Cameroun (Campo) et la Guinée équatoriale (Ebebiyin) ;
\item le \textbf{barrage de Nachtigal} (420 MW, sur la Sanaga, mis en service progressif 2023-2024), projet national structurant qui, par l'interconnexion électrique prévue (ligne vers le Tchad, vers le Nord-Cameroun/RCA), illustre la logique d'infrastructures d'envergure renforçant l'économie nationale et les interconnexions énergétiques sous-régionales.
\end{itemize}

\courssub{Les acquis citoyens : l'intégration vécue au quotidien}

Au-delà des institutions, l'intégration change la vie des citoyens ordinaires :
\begin{itemize}
\item \textbf{Étudier} : un étudiant camerounais peut suivre une formation dans une université ou une école technique d'un pays voisin (ex. : INJS Yaoundé \leftrightarrow homologues de la zone, écoles d'ingénieurs), et les diplômes tendent à être reconnus dans la zone (accords CAMES).
\item \textbf{Commercer} : les commerçants traversent les frontières pour acheter et vendre (le fameux « commerce frontalier »), faisant vivre des milliers de familles (marchés de Kousseri, Garoua-Boulaï, Kyé-Ossi, Ambam).
\item \textbf{Se soigner} : des malades des zones frontalières se rendent dans les hôpitaux du pays voisin lorsque celui-ci est plus proche ou mieux équipé (ex. : hôpital de Garoua pour le Nord-Tchad, hôpital de Bertoua pour l'Est-RCA).
\item \textbf{Appartenir à une communauté} : les échanges culturels, sportifs (compétitions de football UNIFFAC, festivals) et religieux nourrissent un \cle{sentiment d'appartenance africaine}, l'idée que nous sommes d'abord des frères avant d'être des voisins.
\end{itemize}

\courssub{Les obstacles à l'intégration}

\courssub{Des échanges intra-africains très faibles}

Voici le paradoxe central : les pays de la sous-région commercent beaucoup plus avec l'extérieur (Europe, Chine, Inde) qu'entre eux. Les échanges entre pays africains restent faibles, et ceux entre pays de la CEMAC sont encore plus faibles.

\begin{exemplebox}[Un chiffre qui fait réfléchir]
Les échanges entre pays africains représentent \textbf{environ 15 \%} de leur commerce total (chiffre 2023, en progression grâce à la ZLECAf), contre \textbf{plus de 60 \%} en Europe (intra-UE). Autrement dit : un Européen vend surtout à ses voisins européens, alors qu'un Africain vend surtout hors d'Afrique. Dans la CEMAC, la part des échanges internes à la zone reste \textbf{inférieure à 5 \%} du commerce total des États membres (hors pétrole brut réexporté via pipelines).
\end{exemplebox}

\begin{center}
\begin{popfigure}
\begin{tikzpicture}
\begin{axis}[
    ybar,
    width=0.95\linewidth,
    height=7cm,
    bar width=1.2cm,
    ymin=0, ymax=70,
    ylabel={Part des échanges internes (\% du commerce total)},
    symbolic x coords={CEMAC, Afrique (total), Union Europ\'eeenne, ASEAN (Asie du Sud-Est)},
    xtick=data,
    nodes near coords,
    every node near coord/.append style={font=\small\bfseries},
    ymajorgrids=true,
    grid style={popline!40},
    axis line style={popdark},
    enlarge x limits=0.25,
]
\addplot[fill=poporange,draw=popdark] coordinates {(CEMAC,4) (Afrique (total),15) (Union Europ\'eeenne,63) (ASEAN (Asie du Sud-Est),23)};
\end{axis}
\end{tikzpicture}
\legende{Part approximative des échanges commerciaux réalisés à l'intérieur de chaque zone (en \% du commerce total). La CEMAC et l'Afrique commercent très peu entre partenaires de la même zone, contrairement à l'Europe et à l'ASEAN (ordres de grandeur 2023, sources : BAD, UNCTAD, OMC).}
\end{popfigure}
\end{center}

Pourquoi une telle faiblesse ? La raison structurelle principale est que les économies de la zone sont \cle{concurrentes} au lieu d'être \emph{complémentaires} : presque toutes exportent les mêmes produits bruts (pétrole, bois, cacao, coton, minerais, café) et importent les mêmes produits manufacturés (machines, médicaments, véhicules, denrées alimentaires transformées). Le Gabon n'a pas besoin du pétrole du Tchad, puisqu'il en produit lui-même ; le Cameroun n'achète pas le cacao du Gabon, il en produit. Faute de produits transformés différenciés à s'échanger, le commerce intra-zone reste limité.

\courssub{Les obstacles matériels et administratifs}

Même quand la volonté politique existe, la réalité du terrain freine les échanges :
\begin{itemize}
\item le \textbf{mauvais état des routes} : de nombreux axes frontaliers sont en terre ou dégradés, impraticables en saison des pluies ; un camion met parfois 2 à 3 semaines entre Douala et Bangui (600 km) ou Douala et N'Djamena (1800 km) ;
\item les \textbf{tracasseries aux frontières} : contrôles multiples (police, douane, gendarmerie, eaux-et-forêts, santé), demandes de paiements irréguliers (rackets), lenteurs administratives qui découragent les transporteurs et renchérissent les coûts logistiques de 30 à 50 \% ;
\item la \textbf{non-application des décisions communautaires} : les chefs d'État signent des accords (libre circulation, libre installation, droit de résidence, passeport CEMAC), mais les administrations nationales tardent à transposer ces textes en droit interne et à les appliquer sur le terrain (ex. : exigence illégale de visa ou de carte de séjour pour les ressortissants CEMAC).
\end{itemize}

\courssub{Les obstacles politiques et sécuritaires}

L'intégration suppose des États stables et coopératifs. Or la sous-région connaît :
\begin{itemize}
\item des \textbf{conflits armés} récurrents, comme les crises à répétition en République centrafricaine (depuis 2013) ou à l'Est de la RDC, qui bloquent les corridors commerciaux, détruisent les infrastructures et provoquent des afflux de réfugiés (le Cameroun en accueille plus de 450 000, dont 350 000 Centrafricains à l'Est) ;
\item des \textbf{crises institutionnelles et coups d'État} (Mali, Guinée, Burkina, Niger, Gabon en août 2023, Tchad 2021), qui fragilisent la confiance entre partenaires et paralysent parfois le fonctionnement des organes communautaires (sanctions UA/CEEAC, suspension) ;
\item l'\textbf{attachement des États à leur souveraineté} : chaque pays hésite à transférer une partie de son pouvoir de décision (fiscalité, douane, politique industrielle, défense) à une institution commune, par crainte de perdre son autonomie. C'est l'obstacle politique majeur à toute intégration supranationale dans le monde.
\end{itemize}

\begin{center}
\begin{popfigure}
\begin{tikzpicture}
\node[fill=popgreenL, draw=popgreen!70!black, rounded corners, minimum width=0.47\linewidth, minimum height=1.2cm, font=\bfseries\large, align=center] (A) at (0,0) {Les acquis majeurs};
\node[fill=poporangeL, draw=poporange!80!black, rounded corners, minimum width=0.47\linewidth, minimum height=1.2cm, font=\bfseries\large, align=center, anchor=north west] (B) at ([xshift=0.06\linewidth]A.north east) {Les obstacles persistants};
\node[fill=popcream, draw=popgreen!70!black, rounded corners, text width=0.45\linewidth, anchor=north, font=\small, align=left] at ([yshift=-0.2cm]A.south) {%
$\bullet$ Monnaie commune stable (franc CFA, BEAC Yaoundé)\\[2pt]
$\bullet$ Libre circulation partielle (CNI biométrique suffisante pour 90 j)\\[2pt]
$\bullet$ Coopération sécuritaire efficace (FMM vs Boko Haram)\\[2pt]
$\bullet$ Corridors \& ponts frontaliers (Yagoua-Bongou, Campo-Ebebiyin)\\[2pt]
$\bullet$ Infrastructures énergétiques (Nachtigal, interconnexions)\\[2pt]
$\bullet$ Échanges citoyens : études, commerce, soins, culture\\[2pt]
$\bullet$ Sentiment d'appartenance africaine (CAMES, sport, diasporas)};
\node[fill=popcream, draw=poporange!80!black, rounded corners, text width=0.45\linewidth, anchor=north, font=\small, align=left] at ([yshift=-0.2cm]B.south) {%
$\bullet$ Échanges intra-zone $<$ 5 \% (CEMAC), $\sim$15 \% (Afrique)\\[2pt]
$\bullet$ Économies concurrentes (matières premières brutes)\\[2pt]
$\bullet$ Routes dégradées, corridors lents, coûts logistiques élevés\\[2pt]
$\bullet$ Tracasseries frontalières, rackets, non-application textes\\[2pt]
$\bullet$ Conflits armés (RCA, Est RDC), réfugiés, insécurité\\[2pt]
$\bullet$ Crises politiques, coups d'État, sanctions\\[2pt]
$\bullet$ Souveraineté jalouse, faible transfert de compétences};
\end{tikzpicture}
\legende{Bilan comparé de l'intégration sous-régionale : les acquis sont réels et structurants, mais les obstacles structurels et conjoncturels restent massifs. Un bon bilan au Bac présente toujours les deux colonnes avec des faits précis.}
\end{popfigure}
\end{center}

\begin{attention}[Erreur fréquente au Bac : le bilan à charge unique]
Beaucoup d'élèves présentent l'intégration sous-régionale comme un \textbf{échec total}, en ne citant que les problèmes. C'est une erreur de raisonnement qui coûte des points. Un bilan sérieux est toujours \textbf{nuancé} : il reconnaît les acquis (monnaie stable, sécurité commune, libre circulation partielle, infrastructures) \emph{avant} d'exposer les limites (faibles échanges, tracasseries, crises politiques). Une copie qui ne voit qu'un seul côté manque d'esprit critique et d'honnêteté intellectuelle.
\end{attention}

\courssub{Les perspectives : vers une intégration plus profonde}

Malgré les obstacles, l'avenir de l'intégration africaine s'annonce prometteur, grâce à plusieurs dynamiques convergentes :
\begin{itemize}
\item la \textbf{ZLECAf}, grande ambition continentale (voir ci-dessous) ;
\item l'\textbf{industrialisation locale} : transformer sur place le cacao, le bois, le coton, le pétrole, l'aluminium, au lieu d'exporter les matières brutes, créerait des produits différenciés que les pays voisins pourraient enfin s'échanger (ex. : engrais camerounais pour le coton tchadien, bois transformé gabonais pour le BTP congolais) ;
\item la \textbf{jeunesse et la mobilité} : les jeunes Africains étudient, travaillent et entreprennent de plus en plus au-delà des frontières (startups fintech, e-commerce, transport) ; ils sont les acteurs naturels de l'unité par le bas ;
\item l'\textbf{Agenda 2063} de l'Union Africaine, feuille de route pour faire de l'Afrique un continent prospère, intégré et pacifique à l'horizon 2063 (objectifs : réseau routier/ferroviaire transafricain, marché unique de l'air, passeport africain).
\end{itemize}

\begin{definition}[La ZLECAf (AfCFTA)]
La \cle{ZLECAf} (Zone de libre-échange continentale africaine / African Continental Free Trade Area), adoptée en 2018 (Kigali), entrée en vigueur le 30 mai 2019 et lancement du commerce préférentiel le 1er janvier 2021, est un accord entre les pays de l'UA (54 ratifications sur 55) visant à supprimer progressivement les droits de douane sur \textbf{90 \% des lignes tarifaires} (produits non sensibles), puis 7 \% (produits sensibles) et 3 \% (exclus), afin de créer un \textbf{marché unique} d'environ \textbf{1,4 milliard de consommateurs} (2024) et un PIB cumulé de 3 400 milliards USD. C'est le plus grand projet de libre-échange du monde par le nombre de pays participants. Son objectif : faire passer le commerce intra-africain de \textasciitilde 15 \% à 25-30 \% d'ici 2035 (scénario BAD/FMI), générant des millions d'emplois.
\end{definition}

\courssub{Application à la vie courante : être citoyen de l'intégration}

L'intégration ne se fait pas seulement dans les sommets de chefs d'État : elle se vit dans les comportements quotidiens de chaque citoyen, y compris dans un atelier, une entreprise, un quartier ou sur les réseaux sociaux. Voici les comportements qui font avancer l'unité :

\begin{methode}[Comportements citoyens pour l'intégration sous-régionale]
\begin{enumerate}
\item \textbf{Respecter les étrangers résidents} : le Tchadien, le Centrafricain, le Gabonais, le Congolais, l'Équato-Guinéen qui vit au Cameroun est un frère de la sous-région (droit communautaire) ; lui refuser un service, l'insulter, l'exploiter (loyer excessif, salaire indécent) ou le dénoncer arbitrairement est une faute contre l'unité et la loi.
\item \textbf{Rejeter la xénophobie et le tribalisme} : la haine de l'étranger ou de l'autre ethnie détruit le vivre-ensemble et contredit tous les accords signés par nos États ; sur les réseaux sociaux, il faut \textbf{refuser de partager} les discours de haine, les rumeurs xénophobes ou les appels à la violence contre les ressortissants des pays voisins.
\item \textbf{Consommer les produits locaux et sous-régionaux} : acheter le riz (SEMRY, UNVDA), le tissu (CICAM, coton tchadien/gabonais), le savon, les conserves, les matériaux de construction produits au Cameroun ou dans un pays voisin (ex. ciment CIMENCAM, CIMAF, DANGOTE), c'est créer des emplois ici plutôt qu'ailleurs et renforcer le commerce intra-africain (cible ZLECAf).
\item \textbf{Apprendre à connaître ses voisins} : découvrir la culture, les langues (sango, arabe tchadien, fang, ewondo, bassa, douala...), l'histoire et la géographie des pays de la sous-région nourrit le \cle{sentiment d'appartenance africaine} et lève les préjugés.
\item \textbf{Signaler les tracasseries} : tout citoyen témoin d'un racket à un poste de contrôle frontalier ou d'un refus illégal de droit (entrée avec CNI, travail, soin) doit le signaler à sa hiérarchie, aux ONG de défense des droits (REDHAC, ACDH) ou aux autorités consulaires ; le silence est complicité.
\end{enumerate}
\end{methode}

\begin{methode}[Justifier une réponse au Bac en trois étapes (Méthode « AAC »)]
Au Probatoire ou au Baccalauréat technique, une question de bilan (par exemple : « L'intégration sous-régionale est-elle un succès ? Justifiez. » ou « Faites le bilan de l'intégration du Cameroun en Afrique centrale. ») se traite en trois étapes rigoureuses :
\begin{enumerate}
\item \textbf{Affirmer} clairement un point de vue nuancé : « L'intégration sous-régionale est un \textbf{succès partiel} en construction. » (Éviter « Oui » ou « Non » secs).
\item \textbf{Appuyer} ce point de vue sur \textbf{deux faits ou chiffres précis} (un acquis, un obstacle) : « \emph{Acquis} : la monnaie commune (franc CFA) garantit la stabilité des prix et facilite les transactions pour 60 millions de personnes. \emph{Obstacle} : le commerce intra-CEMAC reste inférieur à 5 \% du commerce total, car les économies exportent les mêmes matières premières brutes. »
\item \textbf{Conclure} en reliant l'argument à la notion d'unité et d'avenir : « L'intégration avance donc concrètement sur le plan monétaire et sécuritaire, mais l'unité économique réelle des peuples exige encore de lever les obstacles structurels (industrialisation, infrastructures, gouvernance frontalière) grâce à la ZLECAf et à l'engagement citoyen de chacun. »
\end{enumerate}
Une justification sans fait ni chiffre daté est une opinion ; une justification avec des faits précis est un raisonnement scientifique.
\end{methode}

\begin{exoresolu}[Sujet type Bac : construire un bilan nuancé]
\textbf{Énoncé.} Un élève déclare : « La CEMAC ne sert à rien, c'est un échec total. » En deux paragraphes argumentés, montrez que ce jugement est excessif.
\tcblower
\textbf{Solution détaillée.}

\emph{Paragraphe 1 — Les acquis réels.} Le jugement de cet élève est excessif car la CEMAC a des acquis majeurs, vérifiables. D'abord, la \textbf{monnaie commune}, le franc CFA, gérée par la BEAC (siège à Yaoundé), garantit depuis 1945 (réforme 1994) la stabilité des prix et la convertibilité, facilitant les échanges et l'investissement entre les six pays membres : un entrepreneur camerounais paie ses fournisseurs gabonais ou tchadiens dans la même monnaie, sans risque de change. Ensuite, la \textbf{libre circulation des personnes} est effective : un Camerounais entre au Tchad, en RCA ou au Gabon avec sa seule carte nationale d'identité biométrique, sans visa, pour 90 jours. Enfin, la \textbf{coopération sécuritaire}, notamment la \textbf{Force multinationale mixte (FMM)} contre Boko Haram dans le bassin du Lac Tchad (Cameroun, Tchad, Niger, Nigeria, Bénin), a permis de réduire drastiquement la capacité de nuisance des groupes terroristes depuis 2015, protégeant des millions de civils.

\emph{Paragraphe 2 — Les limites structurelles et la conclusion.} Cependant, ces acquis ne doivent pas cacher des limites profondes : le \textbf{commerce intra-CEMAC reste inférieur à 5 \%} du commerce total des États membres (contre 63 \% dans l'UE), car les économies sont \textbf{concurrentes} (toutes exportatrices de pétrole, bois, cacao, coton bruts) et non complémentaires. S'y ajoutent les \textbf{obstacles matériels} (routes dégradées, corridors lents Douala--Bangui/N'Djamena), les \textbf{tracasseries frontalières} (rackets, non-application du droit d'établissement) et les \textbf{crises politiques} (RCA, coups d'État au Gabon/Tchad) qui paralysent la coopération. En conclusion, l'intégration n'est ni un échec total ni un succès complet : c'est un \textbf{succès partiel} en construction, dont l'approfondissement passe par la \textbf{ZLECAf} (marché unique continental), l'\textbf{industrialisation locale} (transformation sur place), la \textbf{régulation effective des frontières} et l'\textbf{engagement citoyen} de chacun (respect des étrangers, rejet de la xénophobie, consommation locale).
\end{exoresolu}

\coursec{TP guidé : cartographier l'intégration du Cameroun}

\courssub{Consignes}

À partir des connaissances du cours et de la carte muette de l'Afrique centrale fournie (ou tracée au tableau), réaliser les items suivants. Ce TP évalue la capacité à \textbf{localiser}, \textbf{hiérarchiser} et \textbf{commenter} les données spatiales de l'intégration.

\begin{exercice}[Cartographie de l'intégration camerounaise]
\begin{enumerate}
\item \textbf{Fond de carte} : Tracer le contour de l'Afrique centrale (CEEAC) et positionner les 11 capitales. Colorier différemment les 6 pays CEMAC et les 5 pays CEEAC hors CEMAC. Tracer la flèche du Nord et l'échelle graphique.
\item \textbf{Institutions} : Localiser par un symbole distinct (carré rouge, triangle bleu, losange vert...) et nommer les sièges des 7 institutions CEMAC : BEAC (Yaoundé), Commission (Bangui), Parlement (Malabo), Cour de justice (N'Djamena), COBAC (Libreville), BDEAC (Brazzaville), Secrétariat CEEAC (Libreville).
\item \textbf{Corridors \& Infrastructures} : Tracer en trait rouge épais les deux corridors majeurs : Douala--N'Djamena (via Ngaoundéré, Garoua, Maroua, Kousseri) et Douala--Bangui (via Bertoua, Batouri, Garoua-Boulaï). Indiquer le port de Douala (ancrage) et le port en eau profonde de Kribi (projet corridor Kribi--Bangui). Positionner le barrage de Nachtigal et la ligne d'interconnexion électrique vers le Tchad.
\item \textbf{Flux \& Frontières} : Représenter par des flèches proportionnelles (largeur) les principaux flux d'exportation (pétrole Tchad/Cameroun via pipeline Kribi, bois, coton, cacao) et d'importation (produits manufacturés, hydrocarbures raffinés) via Douala/Kribi. Marquer d'un symbole « X » rouge les 3 principaux points de passage frontaliers à fluidifier (Kousseri, Garoua-Boulaï, Kyé-Ossi/Ambam).
\item \textbf{Légende organisée} : Construire une légende hiérarchisée (Institutions, Infrastructures, Flux, Obstacles) avec les codes couleurs/figures utilisés.
\item \textbf{Commentaire synthétique (10 lignes max)} : En vous appuyant sur votre carte, expliquez pourquoi le Cameroun est le « carrefour » et le « moteur » de l'intégration en Afrique centrale, mais aussi pourquoi les obstacles persistent.
\end{enumerate}
\end{exercice}

\begin{corrige}[TP Cartographie - Éléments attendus]
\begin{enumerate}
\item \textbf{Fond de carte} : Contour approximatif golfe de Guinée -- Lac Tchad -- Oubangui. Capitales : Yaoundé, N'Djamena, Bangui, Libreville, Brazzaville, Malabo, Kinshasa, Sao Tomé, Bujumbura, Kigali, Luanda. Coloriage : CEMAC (vert clair), CEEAC hors CEMAC (jaune clair).
\item \textbf{Institutions} : Symboles bien distincts, lisibles, non chevauchants. Yaoundé (BEAC) = cœur monétaire. Bangui (Commission + Secrétariat CEEAC) = cœur politique/administratif. Malabo = Parlement. N'Djamena = Justice. Libreville = COBAC + Secrétariat CEEAC. Brazzaville = BDEAC.
\item \textbf{Corridors} : Douala--N'Djamena (route + rail jusqu'à Ngaoundéré) = axe vital Tchad. Douala--Bangui = axe vital RCA. Kribi = nouveau pôle (port minéralier, gaz, corridor futur). Nachtigal (Sanaga, au Nord de Yaoundé) + ligne HT vers Ngaoundéré/N'Djamena.
\item \textbf{Flux} : Flèche large sortante : Pétrole brut (pipeline Tchad--Kribi, terminal offshore), Bois (Douala/Kribi), Cacao/Café/Coton (Douala). Flèche large entrante : Produits manufacturés, Riz, Blé, Médicaments, Carburants raffinés (Douala/Kribi). Points « X » : Kousseri (Tchad), Garoua-Boulaï (RCA), Kyé-Ossi/Ambam (Gabon/Guinée équatoriale).
\item \textbf{Légende} : Titre : « Le Cameroun, carrefour de l'intégration en Afrique centrale (CEMAC/CEEAC) ». 4 rubriques. Propreté, lisibilité.
\item \textbf{Commentaire type} : « Le Cameroun est le carrefour de l'intégration centrafricaine car il abrite la BEAC (monnaie), le premier port (Douala) et le premier PIB de la CEMAC, d'où partent les deux corridors vitaux vers les pays enclavés (Tchad, RCA). Il héberge aussi la main-d'œuvre et le marché de consommation les plus importants. Cependant, la carte révèle les freins : la concentration des flux sur Douala (goulot), l'état des routes au-delà de Ngaoundéré/Bertoua, les points de blocage frontaliers (Kousseri, Garoua-Boulaï) et l'absence de liaison ferroviaire continue vers l'Est et le Nord. L'intégration physique reste inachevée. »
\end{enumerate}
\end{corrige}

\begin{aretenir}[L'essentiel de la leçon n°5 : Le Cameroun et les structures d'intégration]
\begin{itemize}
\item \textbf{UA (continentale)} : Cameroun membre fondateur (1963/2002), contributeur majeur aux missions de paix (MINUSCA, FMM), aligné sur Agenda 2063 et ZLECAf. Siège BEAC à Yaoundé.
\item \textbf{CEEAC (11 États, siège Libreville)} : Cadre large de coopération (marché commun visé, paix/sécurité FOMAC, infrastructures). Cameroun = pivot géographique (corridors) et démographique/économique.
\item \textbf{CEMAC (6 États, franc CFA, siège institutions réparties)} : Union monétaire et économique aboutie. Institutions : BEAC (Yaoundé), Commission (Bangui), Parlement (Malabo), Cour (N'Djamena), COBAC (Libreville), BDEAC (Brazzaville).
\item \textbf{Bilan nuancé} : \textbf{Acquis} = Monnaie stable, Libre circulation (CNI), Sécurité (FMM), Infrastructures (corridors, ponts, Nachtigal), Vivre-ensemble citoyens. \textbf{Obstacles} = Commerce intra-zone $<$ 5 \% (économies concurrentes brutes), Routes dégradées, Tracasseries/Non-application textes, Conflits/Coups d'État, Souveraineté jalouse.
\item \textbf{Perspectives} = ZLECAf (marché unique 1,4 Md hab., suppression droits de douane 90 \% lignes), Industrialisation locale, Jeunesse mobile, Agenda 2063.
\item \textbf{Citoyenneté active} : Respect étrangers, Rejet xénophobie (réseaux sociaux), Consommer local/sous-régional, Connaître voisins, Signaler tracasseries.
\item \textbf{Méthode Bac (AAC)} : \textbf{A}ffirmer (nuancé) -- \textbf{A}ppuyer (2 faits/chiffres datés : 1 acquis, 1 obstacle) -- \textbf{C}onclure (lien unité/avenir).
\end{itemize}
\end{aretenir}

\coursec{TP guidé : cartographier l'intégration du Cameroun}

La section précédente nous a permis de dresser le bilan de l'intégration sous-régionale : des acquis réels (monnaie commune, libre circulation partielle, institutions communes) mais aussi des obstacles persistants (faiblesse des échanges intra-zone, retards de paiement des contributions, insécurité aux frontières). Ce bilan resterait toutefois abstrait s'il n'était pas \emph{mis en images}. C'est toute l'utilité de ce travail pratique : transformer les connaissances de la leçon en productions concrètes — une carte, un croquis, un diagramme — comme on vous le demandera à l'examen du Probatoire et du Baccalauréat technique. Cartographier, c'est prouver que l'on a compris.

\courssub{Objectifs et organisation du TP}

\begin{definition}[Travail pratique (TP) en ECM]
Un \cle{travail pratique} est une activité de production guidée par laquelle l'élève applique les notions d'une leçon à une tâche concrète : construire une carte, réaliser un croquis, exploiter des données chiffrées, analyser un document. Il mobilise deux savoir-faire officiels : \emph{appliquer les notions de l'unité à des situations concrètes} et \emph{rédiger et justifier une démarche}.
\end{definition}

La classe est divisée en groupes de quatre à six élèves. Chaque groupe réalise les trois productions suivantes, puis les présente devant la classe :

\begin{enumerate}
\item \textbf{TP 1} : une carte muette complétée des pays de la CEMAC et de la CEEAC ;
\item \textbf{TP 2} : un croquis « Le Cameroun, plaque tournante de l'Afrique centrale » ;
\item \textbf{TP 3} : un diagramme construit à partir des données du PIB des pays de la CEMAC.
\end{enumerate}

Chaque production est évaluée selon trois critères : l'\cle{exactitude} des informations reportées, la \cle{lisibilité} (légende claire, titre explicite) et la \cle{justification} orale des choix effectués.

\courssub{TP 1 : construire la carte de la CEMAC et de la CEEAC}

\textbf{Matériel} : un fond de carte vierge de l'Afrique centrale (photocopie ou fond tracé au tableau), des crayons de couleur, une règle, la liste officielle des États membres.

\textbf{Consigne} : placer et colorier les 6 pays de la \cle{CEMAC} (Cameroun, Congo, Gabon, Guinée équatoriale, République centrafricaine, Tchad) avec une première couleur, puis les 11 pays de la \cle{CEEAC} (les 6 de la CEMAC, plus l'Angola, le Burundi, la République démocratique du Congo, le Rwanda, São Tomé-et-Principe) avec une seconde couleur ou un tramage. Compléter avec un titre et une légende.

\begin{methode}[Réussir la carte muette complétée]
\begin{enumerate}
\item \textbf{Repérer} d'abord le Cameroun, pivot de la zone : il est au centre, entre le golfe de Guinée et le Tchad.
\item \textbf{Colorier} les 6 pays de la CEMAC d'une même couleur : la CEMAC est un sous-ensemble de la CEEAC, il faut que cela se voie.
\item \textbf{Tramer ou hachurer} les 5 pays membres de la CEEAC seulement (Angola, Burundi, RDC, Rwanda, São Tomé-et-Principe).
\item \textbf{Nommer} chaque État et marquer sa capitale par un point ou une étoile.
\item \textbf{Rédiger la légende} : une case colorée = un ensemble d'États ; un figuré = une capitale.
\item \textbf{Donner un titre} qui exprime une idée, par exemple : « La CEMAC, cœur économique de la CEEAC ».
\end{enumerate}
\end{methode}

\begin{attention}[Piège fréquent : confondre les deux organisations]
L'erreur la plus répandue consiste à colorier les 11 pays de la même façon, ce qui efface la différence entre CEMAC et CEEAC. Retenez la hiérarchie : la \cle{CEMAC} (6 États, union économique et monétaire, franc CFA) est \emph{incluse} dans la \cle{CEEAC} (11 États, communauté politique plus large). Tout pays de la CEMAC est membre de la CEEAC, mais l'inverse est faux : le Rwanda ou l'Angola ne sont pas dans la CEMAC et n'utilisent pas le franc CFA.
\end{attention}

\courssub{TP 2 : réaliser le croquis « Le Cameroun, plaque tournante de l'Afrique centrale »}

Le \cle{croquis} est l'exercice roi de la géographie scolaire au Cameroun. Contrairement à la carte complétée, il n'est pas un simple report : c'est un \emph{raisonnement dessiné}. On ne montre que ce qui sert l'idée annoncée dans le titre.

\begin{definition}[Croquis géographique]
Un \cle{croquis} est un schéma cartographique simplifié, réalisé à main levée, qui organise dans l'espace quelques faits essentiels (3 à 5 maximum) pour illustrer une idée exprimée par le titre. Il comprend obligatoirement : un fond simplifié, des figurés, une légende organisée et un titre.
\end{definition}

Ici, l'idée à démontrer est la suivante : grâce à sa position, à son port de Douala et au siège de la BEAC à Yaoundé, le Cameroun est la \cle{plaque tournante} (le carrefour) de l'Afrique centrale. Les faits essentiels à reporter sont donc :

\begin{itemize}
\item le \textbf{port de Douala}, principale porte maritime de la sous-région ;
\item le \textbf{corridor Douala--N'Djamena} (route et voie ferrée jusqu'à Ngaoundéré, puis route), qui désenclave le Tchad ;
\item le \textbf{corridor Douala--Bangui}, qui désenclave la Centrafrique ;
\item \textbf{Yaoundé}, siège de la Banque des États de l'Afrique centrale (BEAC) ;
\item les \textbf{États enclavés} (Tchad, Centrafrique) qui dépendent de ces corridors.
\end{itemize}

\begin{methode}[Les quatre étapes du croquis au Bac]
\begin{enumerate}
\item \textbf{L'esquisse du fond} : tracer à main levée, au crayon, les contours simplifiés du Cameroun et de ses voisins concernés (Tchad au nord-est, Centrafrique à l'est, golfe de Guinée au sud-ouest). Le fond doit être léger, sans détail inutile.
\item \textbf{Le report des faits essentiels} : placer les figurés (ancre pour le port, étoile pour la capitale, trait gras fléché pour les corridors). Se limiter à 3 à 5 informations : chaque figuré doit servir le titre.
\item \textbf{La légende organisée} : regrouper les figurés par rubriques (villes, axes, espaces), chacune précédée de son symbole. La légende se lit comme le croquis : du plus important au plus secondaire.
\item \textbf{Le titre explicite} : il exprime une idée et non une simple liste. « Le Cameroun, plaque tournante de l'Afrique centrale » est un bon titre ; « Le Cameroun et ses voisins » n'en est pas un.
\end{enumerate}
\end{methode}

\begin{popfigure}
\begin{center}
\begin{tikzpicture}[scale=1.0, font=\small]
% Fond : Cameroun simplifié
\draw[thick, popblue, fill=popblueL] (0,0) -- (1.2,1.6) -- (1.0,3.2) -- (1.8,4.6) -- (2.6,5.2) -- (3.6,4.4) -- (4.2,3.0) -- (3.8,1.4) -- (2.6,0.2) -- (1.2,-0.4) -- cycle;
% Tchad simplifié
\draw[thick, popmuted, fill=popgoldL] (2.6,5.2) -- (3.6,4.4) -- (5.2,4.8) -- (5.6,6.2) -- (3.2,6.6) -- cycle;
\node[popdark] at (4.4,5.7) {\textbf{TCHAD}};
% Centrafrique simplifiée
\draw[thick, popmuted, fill=popgoldL] (4.2,3.0) -- (3.8,1.4) -- (5.4,1.2) -- (6.4,2.4) -- (5.8,3.6) -- (5.2,4.8) -- (3.6,4.4) -- cycle;
\node[popdark] at (5.3,2.9) {\textbf{CENTRAFRIQUE}};
% Golfe de Guinée
\fill[popblue!20] (-2.2,-1.6) rectangle (2.6,-0.4);
\node[popblue, font=\itshape] at (0.2,-1.1) {Golfe de Guinée};
% Nom Cameroun
\node[popdark, font=\bfseries] at (2.2,2.4) {CAMEROUN};
% Douala (port)
\node[popink] (douala) at (1.0,0.1) {};
\draw[popink, fill=popink] (1.0,0.1) circle (0.09);
\node[anchor=south east, popink] at (0.9,0.1) {\textbf{Douala}};
\draw[popink] (1.0,0.35) -- (1.0,0.6) (0.85,0.5) -- (1.15,0.5);
% Yaoundé (BEAC)
\draw[poporange, fill=poporange] (1.9,0.9) circle (0.09);
\node[anchor=south west, poporange] at (1.95,0.85) {\textbf{Yaoundé (BEAC)}};
% N'Djamena
\draw[popink, fill=popink] (3.4,5.6) circle (0.09);
\node[anchor=west, popink] at (3.5,5.6) {\textbf{N'Djamena}};
% Bangui
\draw[popink, fill=popink] (5.5,2.0) circle (0.09);
\node[anchor=west, popink] at (5.6,2.0) {\textbf{Bangui}};
% Corridor Douala-N'Djamena
\draw[line width=2.2pt, popgreen, -{Stealth[length=3mm]}] (1.15,0.25) .. controls (1.6,2.2) and (2.2,4.2) .. (3.3,5.5);
% Corridor Douala-Bangui
\draw[line width=2.2pt, poppurple, -{Stealth[length=3mm]}] (1.2,0.1) .. controls (2.8,0.6) and (4.2,1.2) .. (5.4,1.95);
% Nord
\draw[-{Stealth}, popdark] (-1.6,5.6) -- (-1.6,6.4);
\node[popdark] at (-1.6,6.7) {\textbf{N}};
\end{tikzpicture}
\end{center}
\legende{Modèle de croquis : le Cameroun, plaque tournante de l'Afrique centrale. Figurés : ancre = port de Douala ; point orange = Yaoundé, siège de la BEAC ; flèche verte = corridor Douala--N'Djamena ; flèche violette = corridor Douala--Bangui ; zones jaunes = États enclavés désenclavés par le Cameroun. Échelle indicative : 1 cm pour environ 300 km.}
\end{popfigure}

\begin{popfigure}
\begin{center}
\begin{tikzpicture}[font=\small]
\draw[thick, popdark] (0,0) rectangle (7.2,3.4);
\node[anchor=west, font=\bfseries] at (0.2,3.1) {LÉGENDE};
% Rubrique villes
\node[anchor=west, font=\itshape] at (0.2,2.6) {Villes et fonctions};
\draw[popink, fill=popink] (0.5,2.1) circle (0.09);
\draw[popink] (0.5,2.35) -- (0.5,2.6) (0.35,2.5) -- (0.65,2.5);
\node[anchor=west] at (0.9,2.1) {Port international (Douala)};
\draw[poporange, fill=poporange] (0.5,1.6) circle (0.09);
\node[anchor=west] at (0.9,1.6) {Capitale, siège de la BEAC (Yaoundé)};
\draw[popink, fill=popink] (0.5,1.1) circle (0.09);
\node[anchor=west] at (0.9,1.1) {Capitale d'un État voisin};
% Rubrique axes
\node[anchor=west, font=\itshape] at (4.2,2.6) {Axes et espaces};
\draw[line width=2.2pt, popgreen, -{Stealth[length=3mm]}] (4.4,2.1) -- (5.3,2.1);
\node[anchor=west] at (5.5,2.1) {Corridor vers le Tchad};
\draw[line width=2.2pt, poppurple, -{Stealth[length=3mm]}] (4.4,1.6) -- (5.3,1.6);
\node[anchor=west] at (5.5,1.6) {Corridor vers la Centrafrique};
\fill[popgoldL, draw=popmuted] (4.4,0.9) rectangle (5.3,1.3);
\node[anchor=west] at (5.5,1.1) {État enclavé};
\fill[popblueL, draw=popblue] (0.5,0.35) rectangle (1.4,0.75);
\node[anchor=west] at (1.6,0.55) {Cameroun};
\end{tikzpicture}
\end{center}
\legende{Exemple de légende normalisée : rubriques séparées (villes et fonctions / axes et espaces), un symbole par fait, une formulation courte par symbole.}
\end{popfigure}

\begin{attention}[Erreur fréquente : surcharger le croquis]
Beaucoup de candidats croient qu'un bon croquis est un croquis « plein ». C'est faux : un croquis surchargé de rivières, de reliefs ou de villes secondaires devient illisible et perd sa valeur de démonstration. Un bon croquis est \cle{sélectif} : chaque figuré doit répondre à la question « cet élément prouve-t-il l'idée du titre ? ». Si la réponse est non, on le supprime. Trois faits bien choisis et bien légendés valent mieux que dix faits noyés dans le dessin.
\end{attention}

\courssub{TP 3 : construire un diagramme à partir des données du PIB de la CEMAC}

Troisième production : traduire des données chiffrées en image. On fournit aux groupes le tableau suivant (données simplifiées et arrondies à des fins pédagogiques, ordres de grandeur récents) :

\begin{center}
\begin{tabularx}{\linewidth}{|l|X|X|}\hline
\rowcolor{popblueL} \textbf{Pays de la CEMAC} & \textbf{PIB (milliards de dollars)} & \textbf{Part dans le PIB de la zone (\%)} \\ \hline
Cameroun & 47 & 44 \\ \hline
Gabon & 20 & 19 \\ \hline
Congo & 15 & 14 \\ \hline
Tchad & 13 & 12 \\ \hline
Guinée équatoriale & 12 & 11 \\ \hline
Centrafrique & 3 & 3 \\ \hline
\rowcolor{popcream} \textbf{Total CEMAC} & \textbf{110} & \textbf{100} \\ \hline
\end{tabularx}
\end{center}

Les parts en pourcentage du tableau sont arrondies pour que leur somme fasse 100\,\%. Le calcul exact de la \cle{part relative} du Cameroun donne $\dfrac{47}{110} \times 100 \approx 42,7\,\%$. Maîtriser ce calcul permet de comparer des pays de poids très différents.

\begin{methode}[Construire un diagramme à partir de données]
\begin{enumerate}
\item \textbf{Choisir le type adapté} : un diagramme \emph{en barres} pour comparer des valeurs entre elles ; un diagramme \emph{circulaire} (camembert) pour montrer des parts d'un total. Ici, les deux sont possibles ; les barres lisent mieux les écarts.
\item \textbf{Graduer les axes} : axe horizontal = les pays ; axe vertical = le PIB, avec une échelle régulière (par exemple 1 cm pour 5 milliards de dollars).
\item \textbf{Tracer les barres} de hauteur proportionnelle aux valeurs, de largeur égale, espacées régulièrement.
\item \textbf{Titrer et sourcer} : tout diagramme porte un titre (« Le poids économique des pays de la CEMAC ») et mentionne l'origine des données.
\item \textbf{Conclure en une phrase} : un diagramme se termine toujours par une lecture rédigée.
\end{enumerate}
\end{methode}

\begin{popfigure}
\begin{center}
\begin{tikzpicture}
\begin{axis}[
    ybar, bar width=0.55cm,
    width=0.95\linewidth, height=6.5cm,
    ymin=0, ymax=55,
    ylabel={PIB (milliards de dollars)},
    symbolic x coords={Cameroun, Gabon, Congo, Tchad, G. éq., RCA},
    xtick=data,
    x tick label style={rotate=20, anchor=east, font=\small},
    ytick={0,10,20,30,40,50},
    axis line style={popdark},
    tick style={popdark},
]
\addplot[fill=popblue, draw=popdark] coordinates {(Cameroun,47) (Gabon,20) (Congo,15) (Tchad,13) (G. éq.,12) (RCA,3)};
\end{axis}
\end{tikzpicture}
\end{center}
\legende{Diagramme en barres : le poids économique comparé des six pays de la CEMAC (PIB arrondis, données simplifiées à usage pédagogique). Lecture : le Cameroun représente à lui seul plus de 40\,\% du PIB de la zone, ce qui confirme son rôle de locomotive économique de l'Afrique centrale.}
\end{popfigure}

\begin{exoresolu}[Lecture de données : la part du Cameroun dans la CEMAC]
Énoncé : à partir du tableau ci-dessus, calculer la part exacte du Cameroun dans le PIB total de la CEMAC, puis rédiger une phrase de conclusion utilisable dans une copie d'examen.
\tcblower
Solution détaillée : le PIB total de la zone vaut $47 + 20 + 15 + 13 + 12 + 3 = 110$ milliards de dollars. La part exacte du Cameroun est :
\[ \frac{47}{110} \times 100 \approx 42,7\,\% \]
Phrase de conclusion rédigée : « Avec environ 43\,\% du PIB de la CEMAC, le Cameroun est de loin la première économie de la sous-région : sa stabilité et ses infrastructures (port de Douala, corridors) conditionnent le fonctionnement de toute la communauté. » Remarque : la réponse comporte un calcul juste \emph{et} une phrase rédigée qui en tire la signification — c'est exactement ce que le correcteur attend.
\end{exoresolu}

\courssub{Analyser un document et restituer : rédiger et justifier sa démarche}

Le TP ne s'arrête pas à la production : il faut encore \emph{exploiter} des documents et \emph{présenter} son travail. L'\cle{analyse de documents} suit une démarche en trois temps :

\begin{enumerate}
\item \textbf{Identifier} le document : nature (texte officiel, carte, tableau), auteur ou source, date, thème.
\item \textbf{Extraire} les informations utiles : on souligne les chiffres, les noms d'institutions, les dates clés (par exemple, dans un texte sur la BEAC : « créée en 1972, siège à Yaoundé, émet le franc CFA »).
\item \textbf{Reformuler} en phrases rédigées : interdiction de recopier le document mot à mot ; on réécrit avec ses propres mots en reliant les informations entre elles.
\end{enumerate}

Vient ensuite la \cle{restitution} : chaque groupe présente ses trois productions devant la classe en cinq minutes. La présentation orale doit justifier les choix effectués : « Nous avons choisi un diagramme en barres plutôt qu'un camembert parce que les écarts entre pays se lisent mieux » ; « Nous avons limité le croquis à quatre figurés pour rester lisible ». Justifier sa démarche, c'est montrer que chaque choix est volontaire et réfléchi : c'est précisément le savoir-faire officiel « rédiger et justifier une démarche, une réponse ou une conclusion ».

\begin{exercice}[Pour s'entraîner]
\begin{enumerate}
\item Sur un fond de carte vierge, placer de mémoire les 6 capitales des pays de la CEMAC.
\item Rédiger le titre et la légende d'un croquis montrant que « le Tchad et la Centrafrique dépendent du port de Douala ».
\item Calculer la part du Gabon dans le PIB de la CEMAC et rédiger une phrase de conclusion.
\item Expliquer en cinq lignes pourquoi un croquis surchargé est pénalisé à l'examen.
\end{enumerate}
\end{exercice}

\begin{corrige}[Exercice]
\begin{enumerate}
\item Yaoundé (Cameroun), Brazzaville (Congo), Libreville (Gabon), Malabo (Guinée équatoriale), Bangui (Centrafrique), N'Djamena (Tchad).
\item Titre : « Douala, porte maritime du Tchad et de la Centrafrique ». Légende : ancre = port de Douala ; flèches = corridors Douala--N'Djamena et Douala--Bangui ; aplats = États enclavés.
\item $\dfrac{20}{110} \times 100 \approx 18,2\,\%$. Conclusion : « Avec environ 18\,\% du PIB de la zone, le Gabon est la deuxième économie de la CEMAC, loin derrière le Cameroun. »
\item Un croquis surchargé devient illisible : le correcteur ne distingue plus les faits essentiels, la légende s'allonge et le raisonnement disparaît. Le croquis étant un raisonnement dessiné, tout élément qui ne sert pas le titre doit être supprimé.
\end{enumerate}
\end{corrige}

\begin{savaistu}[Le croquis, une compétence d'examen]
Au Baccalauréat technique, le croquis et l'analyse de documents ne sont pas réservés à l'Histoire-Géographie : ils sont régulièrement évalués en ECM, notamment dans les sujets portant sur l'intégration sous-régionale, la citoyenneté ou les institutions. Un candidat capable de schématiser les corridors de Douala ou de lire un tableau du PIB de la CEMAC dispose d'un avantage décisif : il transforme ses connaissances en démonstration visible, et c'est cette compétence que les correcteurs récompensent le plus.
\end{savaistu}

\begin{aretenir}[L'essentiel du TP]
Cartographier l'intégration, c'est rendre visible ce que la leçon a expliqué : la CEMAC (6 États) au cœur de la CEEAC (11 États), le Cameroun comme plaque tournante grâce au port de Douala, aux corridors de N'Djamena et de Bangui et au siège de la BEAC à Yaoundé, et son poids économique (environ 43\,\% du PIB de la zone). Trois règles d'or : une carte se construit avec une légende hiérarchisée ; un croquis est sélectif (3 à 5 faits, un titre-idée) ; un diagramme exige un type adapté, des axes gradués, un titre, une source et une phrase de conclusion. Enfin, toute production se justifie oralement : savoir expliquer ses choix fait partie de la note.
\end{aretenir}

\newpage

\coursec{Activité d'intégration}

\coursec{Le Cameroun et les structures d'intégration sous-régionales}

\courssub{Objectifs de l'activité}

À l'issue de cette activité d'intégration, l'élève sera capable de :
\begin{itemize}
\item mobiliser les connaissances acquises sur l'\cle{Union Africaine} (UA), la \cle{CEEAC} et la \cle{CEMAC} pour analyser une situation concrète de la vie économique camerounaise ;
\item identifier, dans un dossier documentaire, les organisations d'intégration qui facilitent les échanges commerciaux et les obstacles qui les freinent ;
\item construire et interpréter un diagramme simple à partir de données chiffrées ;
\item rédiger un paragraphe argumenté et justifié répondant à une question ouverte, conformément au savoir-faire exigible : \emph{rédiger et justifier une démarche, une réponse ou une conclusion}.
\end{itemize}

\courssub{Situation}

M. Etoundi est un commerçant installé au marché Mboppi de Douala. Depuis 2018, il fabrique et vend du savon artisanal et exporte aussi du plantain et du manioc transformé (bâtons, farine) vers deux destinations : \textbf{Libreville} (Gabon) et \textbf{Bangui} (Centrafrique). Il emploie six salariés camerounais et rêve d'ouvrir un dépôt à Bata, en Guinée équatoriale.

Voici le dossier documentaire remis aux élèves.

\textbf{Document 1 : les organisations dont fait partie le Cameroun.}

Le Cameroun est membre de l'\cle{Union Africaine} (organisation continentale qui a succédé à l'OUA en 2002 et qui promeut l'unité, la paix et le développement de l'Afrique), de la \cle{CEEAC} (Communauté Économique des États de l'Afrique Centrale, 11 États membres, créée en 1983, siège à Libreville) et de la \cle{CEMAC} (Communauté Économique et Monétaire de l'Afrique Centrale, 6 États : Cameroun, Gabon, Congo, Tchad, Centrafrique, Guinée équatoriale, issue de l'UDEAC en 1994, siège à Bangui). La CEMAC comprend notamment la \cle{BEAC} (Banque des États de l'Afrique Centrale, qui émet le franc CFA) et la \cle{CEEAC} vise à terme une libre circulation des personnes et des biens en Afrique centrale.

\textbf{Document 2 : les ventes de M. Etoundi en 2024.}

\begin{center}
\begin{tabularx}{\linewidth}{|l|X|X|X|}
\hline
\rowcolor{popblueL} \textbf{Destination} & \textbf{Produits vendus} & \textbf{Chiffre d'affaires (millions FCFA)} & \textbf{Part du total} \\
\hline
Cameroun (marché national) & Savon, plantain, manioc & 60 & 60 \% \\
\hline
Libreville (Gabon) & Savon, plantain & 25 & 25 \% \\
\hline
Bangui (Centrafrique) & Savon, farine de manioc & 15 & 15 \% \\
\hline
\textbf{Total} & & \textbf{100} & \textbf{100 \%} \\
\hline
\end{tabularx}
\end{center}

\textbf{Document 3 : témoignage de M. Etoundi.}

\begin{quote}
\og Quand je vends à Libreville, je n'ai pas besoin de changer d'argent : le Gabon utilise le même franc CFA que nous, grâce à la BEAC. Mes savons fabriqués à Douala entrent au Gabon sans droits de douane, parce que nous sommes tous les deux dans la CEMAC. Mais à la frontière de Kiossi, je perds parfois deux jours : les contrôles sont longs et certains agents demandent des paiements non officiels. Vers Bangui, la route est mauvaise et peu sûre, ce qui augmente mes frais de transport de près de 30 \%. Pourtant, sans ces marchés, je n'aurais jamais pu embaucher mes six employés. On parle aussi d'un grand marché africain, la ZLECAf, porté par l'Union Africaine : si cela marche, je pourrai vendre un jour jusqu'au Nigeria sans droits de douane. \fg
\end{quote}

\textbf{Problème posé :} l'intégration sous-régionale (CEMAC, CEEAC) et continentale (UA) est-elle une chance pour les Camerounais comme M. Etoundi ?

\courssub{Consignes}

Le travail se fait en groupes de quatre pour les questions 1 à 4, puis individuellement pour la question 5.

\begin{enumerate}
\item À partir du document 1, citer les \textbf{trois organisations d'intégration} dont le Cameroun est membre et préciser, pour chacune, son niveau d'action (continental ou sous-régional) et une de ses missions.
\item En vous appuyant sur le témoignage de M. Etoundi (document 3), relever \textbf{deux éléments qui facilitent} son activité commerciale et dire à quelle organisation chacun est lié.
\item Toujours dans le document 3, relever \textbf{trois obstacles} qui freinent les échanges de M. Etoundi et proposer, pour l'un d'eux, une solution que pourrait apporter une meilleure intégration sous-régionale.
\item À l'aide du document 2, construire un \textbf{diagramme en barres} représentant le chiffre d'affaires de M. Etoundi selon les trois destinations, puis calculer la part (en pourcentage) des ventes réalisées \emph{hors du Cameroun}.
\item \textbf{Production individuelle :} rédiger un paragraphe argumenté de 10 à 15 lignes répondant à la question : \og L'intégration sous-régionale est-elle une chance pour les Camerounais ? \fg. Vous justifierez votre réponse par \textbf{au moins deux arguments tirés des documents} et vous formulerez une conclusion claire.
\end{enumerate}

\courssub{Ressources à mobiliser}

\begin{aretenir}[Connaissances utiles pour l'activité]
\begin{itemize}
\item L'\cle{Union Africaine} : organisation continentale (2002, succède à l'OUA) ; objectifs : unité africaine, paix, développement ; elle porte la \cle{ZLECAf} (Zone de libre-échange continentale africaine).
\item La \cle{CEEAC} : 11 États d'Afrique centrale ; objectif : intégration économique et libre circulation dans la sous-région.
\item La \cle{CEMAC} : 6 États ; union économique et monétaire ; institutions : \cle{BEAC} (monnaie commune, le franc CFA), marché commun, libre circulation des biens et des personnes en principe.
\item Acquis de l'intégration : monnaie commune, suppression des droits de douane entre membres, élargissement des débouchés, paix et solidarité.
\item Obstacles : tracasseries aux frontières, mauvais état des routes, insécurité, faible application des textes, concurrence.
\item Calcul d'une part en pourcentage : $\text{part} = \dfrac{\text{valeur de la partie}}{\text{valeur totale}} \times 100$.
\end{itemize}
\end{aretenir}

\courssub{Démarche et corrigé commenté}

\begin{exoresolu}[Corrigé commenté de l'activité]
\textbf{Question 1.} Identifier les trois organisations et leurs missions.

\tcblower
Le Cameroun est membre de :
\begin{itemize}
\item l'\textbf{Union Africaine} (niveau continental) : promouvoir l'unité, la paix et le développement du continent ; elle porte le projet de ZLECAf ;
\item la \textbf{CEEAC} (niveau sous-régional, Afrique centrale, 11 États) : favoriser l'intégration économique et la libre circulation en Afrique centrale ;
\item la \textbf{CEMAC} (niveau sous-régional, 6 États) : assurer l'union économique et monétaire, avec la BEAC qui émet le franc CFA.
\end{itemize}
\emph{Commentaire :} il faut bien distinguer le niveau continental (UA) du niveau sous-régional (CEEAC, CEMAC), distinction souvent exigée à l'examen.

\textbf{Question 2.} Deux éléments facilitateurs.
\begin{enumerate}
\item La \textbf{monnaie commune} : M. Etoundi vend au Gabon sans changer d'argent car le franc CFA est commun aux pays de la CEMAC grâce à la \textbf{BEAC}.
\item La \textbf{suppression des droits de douane} : ses savons entrent au Gabon sans taxe douanière grâce au \textbf{marché commun de la CEMAC}.
\end{enumerate}

\textbf{Question 3.} Trois obstacles.
\begin{enumerate}
\item Les \textbf{tracasseries à la frontière} de Kiossi (contrôles longs, paiements non officiels) ;
\item le \textbf{mauvais état des routes} vers Bangui ;
\item l'\textbf{insécurité} sur certains axes, qui augmente les frais de transport de 30 \%.
\end{enumerate}
\emph{Solution possible :} appliquer réellement les textes de la CEEAC et de la CEMAC sur la libre circulation (corridors sans contrôles abusifs, guichet unique aux frontières) et construire ensemble les routes transfrontalières.

\textbf{Question 4.} Diagramme et calcul.

\begin{popfigure}
\begin{tikzpicture}
\begin{axis}[
ybar, bar width=0.9cm, width=0.85\linewidth, height=5.5cm,
symbolic x coords={Cameroun, Gabon, Centrafrique},
xtick=data, ymin=0, ymax=70,
ylabel={Chiffre d'affaires (millions FCFA)},
nodes near coords, every node near coord/.append style={font=\small},
]
\addplot[fill=popblue] coordinates {(Cameroun,60) (Gabon,25) (Centrafrique,15)};
\end{axis}
\end{tikzpicture}
\legende{Diagramme en barres du chiffre d'affaires de M. Etoundi par destination en 2024.}
\end{popfigure}

Part des ventes hors du Cameroun :
\[ \text{part} = \frac{25 + 15}{100} \times 100 = 40 \% \]
M. Etoundi réalise donc \textbf{40 \%} de son chiffre d'affaires à l'exportation, dans la sous-région.

\textbf{Question 5.} Exemple de paragraphe argumenté attendu.

\emph{L'intégration sous-régionale est une chance pour les Camerounais, même si des obstacles demeurent. D'abord, elle facilite concrètement les échanges : grâce à la CEMAC et à la BEAC, M. Etoundi vend au Gabon avec la même monnaie, le franc CFA, et sans droits de douane, ce qui réduit ses coûts. Ensuite, elle élargit les débouchés : 40 \% de son chiffre d'affaires provient de Libreville et de Bangui, et ces ventes lui ont permis d'embaucher six salariés camerounais, créant ainsi des emplois. Enfin, l'Union Africaine, avec la ZLECAf, ouvre la perspective d'un marché continental encore plus vaste. Certes, les tracasseries aux frontières, les mauvaises routes et l'insécurité freinent encore les échanges, mais ces problèmes appellent davantage d'intégration, et non moins. En conclusion, l'intégration sous-régionale est bien une chance pour les Camerounais, à condition que les États appliquent pleinement les textes sur la libre circulation.}
\end{exoresolu}

\begin{attention}[Pièges fréquents]
\begin{itemize}
\item Ne pas confondre \textbf{CEEAC} (11 États, intégration politico-économique) et \textbf{CEMAC} (6 États, union économique et monétaire) : le Gabon appartient aux deux, le Nigeria à aucune des deux.
\item Ne pas attribuer le franc CFA à la CEEAC : c'est la \textbf{BEAC}, institution de la \textbf{CEMAC}, qui l'émet.
\item Dans le paragraphe argumenté, une opinion sans argument tiré des documents ne vaut pas de point : chaque affirmation doit être \textbf{justifiée} par une donnée (chiffre, fait, institution).
\item Dans le calcul du pourcentage, bien additionner les deux destinations étrangères avant de diviser par le total.
\end{itemize}
\end{attention}

\courssub{TP guidé : cartographier l'intégration du Cameroun}

\begin{methode}[Réaliser un croquis schématique de l'intégration sous-régionale]
\begin{enumerate}
\item \textbf{Tracer le cadre :} Dessiner un polygone simplifié représentant l'Afrique centrale (golfe de Guinée au sud, lac Tchad au nord). Placer une flèche Nord.
\item \textbf{Positionner les 6 États de la CEMAC :} Repérer et nommer les capitales : Yaoundé (Cameroun), Libreville (Gabon), Brazzaville (Congo), N'Djamena (Tchad), Bangui (RCA, siège CEMAC), Malabo (Guinée équatoriale). Les colorier en \textcolor{popblue}{bleu CEMAC}.
\item \textbf{Ajouter les 5 États CEEAC hors CEMAC :} Kinshasa (RDC), Luanda (Angola), Kigali (Rwanda), Bujumbura (Burundi), São Tomé (São Tomé-et-Principe). Les colorier en \textcolor{popgreen}{vert CEEAC seul}.
\item \textbf{Symboliser les institutions :} Un pictogramme \textbf{banque} à Yaoundé (siège BEAC) et un \textbf{bâtiment} à Bangui (siège CEMAC) et Libreville (siège CEEAC).
\item \textbf{Montrer les flux :} Flèches d'échanges (savon, plantain, manioc) de Douala vers Libreville et Bangui. Trait discontinu pour la route Douala--Bangui (obstacle).
\item \textbf{Légende et titre :} \og Le Cameroun dans l'intégration sous-régionale (CEMAC / CEEAC) \fg. Échelle indicative.
\end{enumerate}
\end{methode}

\begin{popfigure}
\begin{tikzpicture}[scale=0.8, every node/.style={font=\small}]
% Coastline / rough shape
\draw[popline, thick] (0,0) -- (8,0) -- (9,2) -- (8,5) -- (5,6) -- (2,5.5) -- (0,3) -- cycle;
% Countries - CEMAC (Blue)
\node[popblue, fill=popblueL, circle, draw=popblue, thick, minimum size=6pt] (Yde) at (3.5, 2.2) {};
\node[popblue, fill=popblueL, circle, draw=popblue, thick, minimum size=6pt] (Lbv) at (2.5, 1.2) {};
\node[popblue, fill=popblueL, circle, draw=popblue, thick, minimum size=6pt] (Bzv) at (2.8, 2.8) {};
\node[popblue, fill=popblueL, circle, draw=popblue, thick, minimum size=6pt] (Ndj) at (6.5, 3.5) {};
\draw[popblue, thick] (Yde) -- (Lbv) -- (Bzv) -- (Ndj) -- (Yde) -- cycle; % CEMAC link
% Countries - CEEAC only (Green)
\node[popgreen, fill=popgreenL, circle, draw=popgreen, thick, minimum size=6pt] (Kin) at (4.5, 3.5) {};
\node[popgreen, fill=popgreenL, circle, draw=popgreen, thick, minimum size=6pt] (Lua) at (3.0, 0.5) {};
\node[popgreen, fill=popgreenL, circle, draw=popgreen, thick, minimum size=6pt] (Kig) at (5.5, 2.0) {};
\node[popgreen, fill=popgreenL, circle, draw=popgreen, thick, minimum size=6pt] (Buj) at (5.8, 1.5) {};
% Capitals labels
\node[above right] at (Yde) {Yaoundé};
\node[below left] at (Lbv) {Libreville};
\node[above left] at (Bzv) {Bangui};
\node[above right] at (Ndj) {N'Djaména};
\node[above right] at (Kin) {Kinshasa};
\node[below right] at (Lua) {Luanda};
\node[above] at (Kig) {Kigali};
\node[below] at (Buj) {Bujumbura};
% Seats
\node[popblue, star, star points=5, star point ratio=2.25, fill=popgold, draw=poporange, inner sep=1.5pt] at (Bzv) {};
\node[popblue, star, star points=5, star point ratio=2.25, fill=popgold, draw=poporange, inner sep=1.5pt] at (Yde) {};
\node[popgreen, star, star points=5, star point ratio=2.25, fill=popgold, draw=poporange, inner sep=1.5pt] at (Lbv) {};
% Flux Etoundi
\draw[poporange, very thick, ->, >=Stealth] (Yde) .. controls (3,1.5) .. (Lbv) node[midway, above, sloped, font=\tiny] {Savon, Plantain};
\draw[poporange, very thick, ->, >=Stealth] (Yde) .. controls (3.5,2.8) .. (Bzv) node[midway, above, sloped, font=\tiny] {Savon, Farine};
% Legend
\begin{scope}[shift={(9.5, 5)}]
\node[popblue, fill=popblueL, circle, draw=popblue, thick, minimum size=6pt] at (0,0) {}; \node[right] at (0.4,0) {Pays CEMAC (6)};
\node[popgreen, fill=popgreenL, circle, draw=popgreen, thick, minimum size=6pt] at (0,-0.7) {}; \node[right] at (0.4,-0.7) {Pays CEEAC seulement (5)};
\node[popblue, star, star points=5, star point ratio=2.25, fill=popgold, draw=poporange, inner sep=1.5pt] at (0,-1.4) {}; \node[right] at (0.4,-1.4) {Sièges institutionnels};
\draw[poporange, very thick, ->, >=Stealth] (0,-2.1) -- (0.7,-2.1); \node[right] at (0.8,-2.1) {Flux commerciaux M. Etoundi};
\draw[popline, dashed, thick] (0,-2.8) -- (0.7,-2.8); \node[right] at (0.8,-2.8) {Axe routier difficile};
\end{scope}
% North arrow
\node[draw, circle, fill=popcream, inner sep=2pt] at (7.5, 5.5) {N};
\draw[->, thick] (7.5, 5.5) -- (7.5, 5.9);
\end{tikzpicture}
\legende{Croquis schématique : le Cameroun au cœur de la CEMAC (bleu) et de la CEEAC (vert). Les étoiles dorées indiquent les sièges (BEAC à Yaoundé, CEMAC à Bangui, CEEAC à Libreville). Les flèches orange montrent les exportations de M. Etoundi.}
\end{popfigure}

\courssub{Bilan de l'activité}

Cette activité a permis de mettre en évidence que l'intégration sous-régionale n'est pas une notion abstraite : elle touche directement la vie d'un commerçant de Douala, ses coûts, sa monnaie et ses employés. Les élèves ont appris à \textbf{mobiliser ensemble} les trois organisations étudiées dans la leçon (UA, CEEAC, CEMAC) pour analyser un dossier documentaire, à distinguer les \textbf{acquis} de l'intégration (monnaie commune, marché commun, élargissement des débouchés, perspective de la ZLECAf) de ses \textbf{obstacles} (tracasseries frontalières, routes dégradées, insécurité, application insuffisante des textes). Ils ont enfin pratiqué deux savoir-faire exigibles au Probatoire et au Baccalauréat technique : exploiter des données chiffrées pour construire un diagramme et calculer un pourcentage, puis \textbf{rédiger une conclusion argumentée et justifiée} sur une question de citoyenneté économique. La leçon débouche ainsi sur une conviction citoyenne : l'intégration est une chance pour le Cameroun, à condition que chacun — États, commerçants, citoyens — œuvre à la rendre effective.

\newpage

\coursec{Méthodes et exercices résolus}

\coursec{Le Cameroun et les structures d'intégration sous-régionales}

\courssub{Savoir-faire 1 : Définir l'intégration sous-régionale et l'illustrer par des exemples camerounais}

\begin{methode}[Définir et illustrer une notion d'intégration]
Pour traiter une question de définition ou d'explication au Bac technique :
\begin{enumerate}
\item \textbf{Définir} la notion avec ses mots-clés : l'\cle{intégration sous-régionale} est le processus par lequel des États voisins créent des institutions communes et abaissent les barrières (douanières, monétaires, politiques) pour renforcer leurs échanges et leur solidarité.
\item \textbf{Dégager les objectifs} : libre circulation des personnes et des biens, monnaie commune, paix et sécurité, développement économique partagé.
\item \textbf{Illustrer} par au moins deux organisations auxquelles le Cameroun appartient : la \cle{CEMAC} et la \cle{CEEAC}, sans oublier le niveau continental (l'\cle{Union Africaine}).
\item \textbf{Ancrer dans le concret} : citer un fait de la vie courante (le franc CFA utilisé à l'atelier ou au marché, le marché de Douala approvisionné par le Tchad ou le Gabon, le déplacement sans visa dans l'espace CEMAC).
\item \textbf{Conclure} par une phrase qui relie la notion à l'intérêt pour le Cameroun.
\end{enumerate}
Réflexe rédactionnel : une définition = une phrase complète + un exemple + un intérêt. Jamais une simple liste de sigles.
\end{methode}

\begin{exoresolu}[Définir l'intégration sous-régionale]
\textbf{Énoncé.} Après avoir défini l'intégration sous-régionale, citez deux organisations d'intégration auxquelles le Cameroun appartient et montrez, à l'aide d'un exemple concret, l'intérêt de cette intégration pour un jeune technicien camerounais.
\tcblower
\textbf{Solution détaillée.}

\emph{Étape 1 : la définition.} L'intégration sous-régionale est le processus par lequel des États géographiquement proches s'unissent au sein d'institutions communes afin de faciliter la libre circulation des personnes, des biens et des capitaux, d'harmoniser leurs politiques économiques et de garantir la paix dans leur espace commun.

\emph{Étape 2 : les organisations.} Le Cameroun est membre :
\begin{itemize}
\item de la \textbf{CEMAC} (Communauté Économique et Monétaire de l'Afrique Centrale), créée par le traité de N'Djamena en 1994 (entrée en vigueur 1999), qui regroupe six États : Cameroun, Tchad, Centrafrique, Gabon, Congo-Brazzaville et Guinée équatoriale ;
\item de la \textbf{CEEAC} (Communauté Économique des États de l'Afrique Centrale), créée en 1983 à Libreville, qui regroupe onze États d'Afrique centrale.
\end{itemize}
On peut ajouter qu'au niveau continental, le Cameroun est membre de l'Union Africaine (UA), qui a succédé en 2002 à l'OUA.

\emph{Étape 3 : l'exemple concret.} Un jeune titulaire d'un Baccalauréat technique en froid et climatisation peut aller travailler au Gabon ou au Tchad : dans l'espace CEMAC, la libre circulation des personnes est reconnue et la monnaie utilisée partout est le franc CFA (BEAC), ce qui supprime les problèmes de change. De même, un grossiste de Douala achète du bétail tchadien et revend des produits manufacturés camerounais dans toute la sous-région grâce au corridor Douala--N'Djamena.

\emph{Conclusion.} L'intégration sous-régionale élargit le marché, facilite l'emploi et renforce la solidarité entre peuples voisins ; elle est donc un atout direct pour la vie professionnelle du jeune Camerounais.
\end{exoresolu}

\courssub{Savoir-faire 2 : Présenter le Cameroun au sein de l'Union Africaine}

\begin{methode}[Construire un exposé sur le Cameroun et l'UA]
\begin{enumerate}
\item \textbf{Situer l'organisation} : l'\cle{Union Africaine}, créée en 2002 à Durban (Afrique du Sud), succède à l'OUA (1963) ; son siège est à Addis-Abeba (Éthiopie).
\item \textbf{Rappeler l'adhésion du Cameroun} : membre fondateur de l'OUA en 1963, le Cameroun est membre de l'UA depuis sa création.
\item \textbf{Énumérer les objectifs de l'UA} : unité et solidarité africaines, intégration politique et économique du continent, paix et sécurité, promotion des droits humains et de la bonne gouvernance.
\item \textbf{Montrer la participation concrète du Cameroun} : présence aux sommets des chefs d'État, contribution aux missions de paix, accueil d'institutions, application des décisions de l'UA.
\item \textbf{Évaluer} : citer un acquis (coopération, médiation dans les crises) et une limite (faible application de certaines décisions, financement).
\end{enumerate}
Réflexe : toujours distinguer l'UA (échelle continentale) des organisations sous-régionales (CEMAC, CEEAC) : ne pas les confondre dans la copie.
\end{methode}

\begin{exoresolu}[Le Cameroun dans l'Union Africaine]
\textbf{Énoncé.} Présentez l'Union Africaine (date de création, siège, deux objectifs) puis montrez en deux points comment le Cameroun participe à la vie de cette organisation.
\tcblower
\textbf{Solution détaillée.}

\emph{Étape 1 : présentation de l'UA.}
\begin{itemize}
\item \textbf{Date de création} : l'Union Africaine est créée en 2002 à Durban (Afrique du Sud), en remplacement de l'Organisation de l'Unité Africaine (OUA) fondée en 1963 à Addis-Abeba.
\item \textbf{Siège} : Addis-Abeba (Éthiopie).
\item \textbf{Deux objectifs} : (1) promouvoir l'unité, la solidarité et l'intégration politique et économique du continent africain ; (2) défendre la paix, la sécurité et la stabilité, ainsi que les droits humains et la démocratie.
\end{itemize}

\emph{Étape 2 : la participation du Cameroun.}
\begin{itemize}
\item \textbf{Participation diplomatique} : le Cameroun, membre fondateur de l'OUA en 1963, prend part aux sommets des chefs d'État et de gouvernement de l'UA et y défend les positions de l'Afrique centrale ; il applique les décisions de l'organisation dans sa politique nationale.
\item \textbf{Contribution à la paix et à la coopération} : le Cameroun contribue aux efforts de règlement des crises africaines (médiations, coopération sécuritaire contre le terrorisme dans le bassin du lac Tchad) et accueille sur son sol des réunions et des organes de coopération continentale.
\end{itemize}

\emph{Conclusion.} Membre actif depuis l'indépendance, le Cameroun fait de l'UA le cadre privilégié de sa diplomatie africaine : il y défend ses intérêts tout en contribuant à l'unité et à la paix du continent.
\end{exoresolu}

\courssub{Savoir-faire 3 : Expliquer le rôle du Cameroun dans la CEEAC}

\begin{methode}[Traiter une question sur la CEEAC]
\begin{enumerate}
\item \textbf{Identifier} : la \cle{CEEAC} (Communauté Économique des États de l'Afrique Centrale) est créée en 1983 par le traité de Libreville ; son siège est à Libreville (Gabon) ; elle compte onze États membres.
\item \textbf{Préciser ses missions} : promouvoir la coopération économique, la paix et la sécurité en Afrique centrale ; elle est reconnue par l'UA comme la communauté économique régionale de l'Afrique centrale.
\item \textbf{Situer le Cameroun} : membre fondateur, État le plus peuplé et à l'économie la plus diversifiée de la sous-région ; il joue un rôle moteur (port de Douala, corridor vers le Tchad et la Centrafrique).
\item \textbf{Citer un organe de paix} : le Conseil de paix et de sécurité de l'Afrique centrale (COPAX), instrument de prévention et de règlement des conflits.
\item \textbf{Conclure} sur l'importance stratégique du Cameroun (poids démographique, économique et géographique).
\end{enumerate}
Réflexe : apprendre les sigles avec leur signification complète ; au Bac, un sigle non développé au moins une fois coûte des points.
\end{methode}

\begin{exoresolu}[Le Cameroun, pivot de la CEEAC]
\textbf{Énoncé.} Présentez la CEEAC (date, siège, nombre de membres) puis justifiez en deux arguments l'affirmation suivante : « Le Cameroun est un État pivot de l'intégration en Afrique centrale. »
\tcblower
\textbf{Solution détaillée.}

\emph{Étape 1 : présentation de la CEEAC.}
\begin{itemize}
\item \textbf{Date de création} : 1983, par le traité de Libreville.
\item \textbf{Siège} : Libreville (Gabon).
\item \textbf{Membres} : onze États d'Afrique centrale, dont le Cameroun, le Tchad, la Centrafrique, le Gabon, le Congo, la Guinée équatoriale, la République démocratique du Congo, l'Angola, le Burundi, le Rwanda, Sao Tomé-et-Principe.
\end{itemize}
La CEEAC œuvre pour la coopération économique et la paix ; elle dispose notamment du COPAX pour la prévention et la gestion des conflits.

\emph{Étape 2 : le Cameroun, État pivot.}

\emph{Argument 1 — poids économique et logistique.} Le Cameroun possède l'économie la plus diversifiée de la sous-région (agriculture, industrie, bois, pétrole, services). Le port de Douala est la principale porte d'entrée maritime des pays enclavés voisins : le Tchad et la Centrafrique reçoivent l'essentiel de leurs importations par le corridor Douala--N'Djamena et Douala--Bangui.

\emph{Argument 2 — position géographique et démographique.} Charnière entre l'Afrique centrale et l'Afrique de l'Ouest, le Cameroun partage des frontières avec six pays ; avec environ 28 millions d'habitants, il constitue le plus grand marché de la CEEAC, ce qui en fait un partenaire commercial incontournable.

\emph{Conclusion.} Par son économie, son port et sa position géographique, le Cameroun est bien le pivot de l'intégration en Afrique centrale : sans lui, les échanges sous-régionaux seraient fortement ralentis.
\end{exoresolu}

\courssub{Savoir-faire 4 : Décrire le Cameroun dans la CEMAC et ses institutions communes}

\begin{methode}[Présenter la CEMAC et ses institutions]
\begin{enumerate}
\item \textbf{Identifier} : la \cle{CEMAC} (Communauté Économique et Monétaire de l'Afrique Centrale) est créée en 1994 par le traité de N'Djamena (entrée en vigueur en 1999) ; siège à Bangui (Centrafrique) ; six membres : Cameroun, Tchad, Centrafrique, Gabon, Congo, Guinée équatoriale.
\item \textbf{Citer les institutions communes et leurs sièges} :
\begin{itemize}
\item la \cle{BEAC} (Banque des États de l'Afrique Centrale) : Yaoundé, Cameroun — émet le franc CFA ;
\item la \cle{COBAC} (Commission Bancaire de l'Afrique Centrale) : Libreville, Gabon — contrôle les banques ;
\item le Parlement communautaire : Malabo (Guinée équatoriale) ;
\item la Cour de justice communautaire : N'Djamena (Tchad).
\end{itemize}
\item \textbf{Expliquer la monnaie commune} : le franc CFA de la zone CEMAC (XAF), garant de la stabilité des échanges.
\item \textbf{Dégager les avantages pour le Cameroun} : libre circulation, marché élargi, stabilité monétaire, siège de la BEAC à Yaoundé.
\end{enumerate}
Réflexe : faire un tableau « institution -- rôle -- siège » dans son brouillon pour ne pas mélanger les villes.
\end{methode}

\begin{exoresolu}[Les institutions communes de la CEMAC]
\textbf{Énoncé.} Complétez le tableau suivant puis répondez à la question.

\begin{center}
\begin{tabularx}{\linewidth}{|l|X|X|}
\hline
\rowcolor{popblueL} \textbf{Institution} & \textbf{Rôle} & \textbf{Siège} \\
\hline
BEAC & \ldots & \ldots \\
\hline
COBAC & \ldots & \ldots \\
\hline
Parlement communautaire & \ldots & \ldots \\
\hline
\end{tabularx}
\end{center}

Question : pourquoi la présence du siège de la BEAC à Yaoundé est-elle un atout pour le Cameroun ?
\tcblower
\textbf{Solution détaillée.}

\emph{Étape 1 : tableau complété.}

\begin{center}
\begin{tabularx}{\linewidth}{|l|X|X|}
\hline
\rowcolor{popblueL} \textbf{Institution} & \textbf{Rôle} & \textbf{Siège} \\
\hline
BEAC & Banque centrale commune : émet le franc CFA et conduit la politique monétaire de la zone & Yaoundé (Cameroun) \\
\hline
COBAC & Contrôle et surveillance des banques et établissements financiers de la sous-région & Libreville (Gabon) \\
\hline
Parlement communautaire & Organe délibérant représentant les peuples de la CEMAC & Malabo (Guinée équatoriale) \\
\hline
\end{tabularx}
\end{center}

\emph{Étape 2 : l'atout pour le Cameroun.} La présence du siège de la BEAC à Yaoundé place le Cameroun au cœur des décisions monétaires de la sous-région : cela crée des emplois qualifiés, attire les institutions financières, renforce le prestige diplomatique du pays et lui donne une voix privilégiée dans la gestion du franc CFA, monnaie utilisée par les six États membres.

\emph{Conclusion.} En abritant la banque centrale commune, le Cameroun confirme sa position de leader économique de la CEMAC et bénéficie directement de la stabilité monétaire qu'elle garantit.
\end{exoresolu}

\courssub{Savoir-faire 5 : Rédiger et justifier un bilan (acquis et obstacles)}

\begin{methode}[Construire un bilan argumenté en deux parties]
Pour une question du type « Faites le bilan de l'intégration sous-régionale au Cameroun » :
\begin{enumerate}
\item \textbf{Annoncer le plan} en introduction : on présentera d'abord les acquis, puis les obstacles, avant de conclure.
\item \textbf{Acquis} (au moins trois) : monnaie commune stable (franc CFA) ; libre circulation des personnes dans la CEMAC ; corridors commerciaux (Douala--N'Djamena) ; coopération sécuritaire (lutte contre Boko Haram dans le bassin du lac Tchad) ; institutions communes fonctionnelles (BEAC, COBAC).
\item \textbf{Obstacles} (au moins trois) : faiblesse des échanges intra-africains ; économies concurrentes (toutes exportent des matières premières) ; mauvais état des routes et frontières lentes ; non-application de certaines décisions communautaires ; crises sécuritaires.
\item \textbf{Justifier chaque point} par un fait ou un exemple, pas par une simple affirmation.
\item \textbf{Conclure} par un jugement nuancé et une perspective (ZLECAf, industrialisation, infrastructures).
\end{enumerate}
Réflexe : un bilan équilibré = autant de développement pour les acquis que pour les limites ; la conclusion doit trancher avec nuance.
\end{methode}

\begin{exoresolu}[Bilan de l'intégration sous-régionale]
\textbf{Énoncé.} « L'intégration sous-régionale profite au Cameroun, mais elle reste inachevée. » En vous appuyant sur trois acquis et trois obstacles, discutez cette affirmation et concluez.
\tcblower
\textbf{Solution détaillée.}

\emph{Introduction.} Depuis son indépendance, le Cameroun s'est engagé dans plusieurs organisations d'intégration : l'UA au niveau continental, la CEEAC et la CEMAC au niveau sous-régional. On montrera que cette intégration a produit des acquis réels, mais qu'elle se heurte encore à des obstacles qui la rendent inachevée.

\emph{I. Les acquis.}
\begin{enumerate}
\item \textbf{La stabilité monétaire} : le franc CFA, émis par la BEAC dont le siège est à Yaoundé, est une monnaie commune stable qui facilite les échanges entre les six pays de la CEMAC sans risque de change.
\item \textbf{La libre circulation et le commerce} : les ressortissants de la CEMAC circulent librement ; le port de Douala et le corridor Douala--N'Djamena font du Cameroun le fournisseur et le débouché maritime du Tchad et de la Centrafrique.
\item \textbf{La coopération sécuritaire} : face au terrorisme de Boko Haram, le Cameroun coopère avec le Tchad et le Nigeria au sein de la Force multinationale mixte, montrant que l'intégration protège aussi les populations.
\end{enumerate}

\emph{II. Les obstacles.}
\begin{enumerate}
\item \textbf{Des échanges intra-sous-régionaux faibles} : les pays d'Afrique centrale commercent davantage avec l'Europe ou l'Asie qu'entre eux, car leurs économies se ressemblent (exportation de pétrole, bois, cacao).
\item \textbf{Des infrastructures insuffisantes} : routes dégradées, frontières lentes, tracasseries douanières qui renchérissent le transport entre Douala et N'Djamena.
\item \textbf{L'application incomplète des décisions} : la libre circulation des biens et des personnes est souvent entravée sur le terrain, et les États tardent à harmoniser leurs législations.
\end{enumerate}

\emph{Conclusion.} L'affirmation est donc fondée : l'intégration profite au Cameroun (monnaie stable, marchés, sécurité), mais elle reste inachevée faute d'échanges suffisants et d'infrastructures. Des projets comme la Zone de libre-échange continentale africaine (ZLECAf) et l'industrialisation locale pourraient lui donner un nouvel élan.
\end{exoresolu}

\courssub{Savoir-faire 6 : Réaliser un croquis ou une carte schématique de l'intégration du Cameroun (TP guidé)}

\begin{methode}[Cartographier l'intégration du Cameroun]
Pour produire un croquis lisible au Bac ou en TP :
\begin{enumerate}
\item \textbf{Tracer le contour simplifié} du Cameroun (forme triangulaire caractéristique) et placer les pays voisins : Nigeria, Tchad, Centrafrique, Gabon, Congo, Guinée équatoriale.
\item \textbf{Placer les repères essentiels} : Yaoundé (capitale politique, siège de la BEAC), Douala (port et capitale économique), la frontière avec chaque voisin.
\item \textbf{Matérialiser les flux} : flèches pour le corridor Douala--N'Djamena (vers le Tchad) et Douala--Bangui (vers la Centrafrique), épaisseur proportionnelle à l'importance du flux.
\item \textbf{Ajouter les éléments obligatoires} : titre, légende numérotée ou fléchée, orientation (flèche du nord), échelle indicative.
\item \textbf{Rédiger un court commentaire} (5 lignes) qui explique ce que montre le croquis.
\end{enumerate}
Réflexe : un croquis sans titre, sans légende ou sans nord est sanctionné, même s'il est exact.
\end{methode}

\begin{exoresolu}[Croquis : le Cameroun, plaque tournante de l'Afrique centrale]
\textbf{Énoncé.} Réalisez un croquis schématique montrant la position du Cameroun au sein de la CEMAC : vous y ferez figurer les pays membres, les villes de Yaoundé et Douala, ainsi que les deux grands corridors commerciaux vers les pays enclavés. Accompagnez votre croquis d'un commentaire de cinq lignes.
\tcblower
\textbf{Solution détaillée.}

\emph{Étape 1 : le croquis.}

\begin{popfigure}
\begin{tikzpicture}[scale=1.0, every node/.style={font=\small}]
% Pays voisins (blocs simplifiés)
\fill[popmuted!25] (-4.5,2.2) rectangle (-1.4,4.6); % Nigeria
\node at (-2.95,3.4) {NIGERIA};
\fill[popmuted!25] (1.6,3.0) rectangle (5.0,4.8); % Tchad
\node at (3.3,3.9) {TCHAD};
\fill[popmuted!25] (2.6,0.6) rectangle (5.2,2.4); % Centrafrique
\node at (3.9,1.5) {CENTRAFRIQUE};
\fill[popmuted!25] (-1.2,-3.4) rectangle (1.6,-1.8); % Gabon
\node at (0.2,-2.6) {GABON};
\fill[popmuted!25] (2.0,-3.4) rectangle (4.6,-1.6); % Congo
\node at (3.3,-2.5) {CONGO};
\fill[popmuted!25] (-3.4,-2.4) rectangle (-1.6,-1.4); % Guinée équatoriale
\node[align=center] at (-2.5,-1.9){\shortstack{GUINÉE\\ÉQUATORIALE}};
% Cameroun (polygone simplifié)
\fill[poporange!45] (-1.2,4.2) -- (0.4,4.6) -- (1.6,3.4) -- (2.6,2.6) -- (2.4,0.8) -- (1.4,-1.6) -- (-0.6,-1.8) -- (-1.6,-0.6) -- (-2.2,1.0) -- (-1.8,2.6) -- cycle;
\draw[thick,popdark] (-1.2,4.2) -- (0.4,4.6) -- (1.6,3.4) -- (2.6,2.6) -- (2.4,0.8) -- (1.4,-1.6) -- (-0.6,-1.8) -- (-1.6,-0.6) -- (-2.2,1.0) -- (-1.8,2.6) -- cycle;
\node[font=\small\bfseries] at (0.1,2.9) {CAMEROUN};
% Villes avec numéros pour légende
\fill[popink] (0.2,0.6) circle (0.09); \node[right,font=\small] at (0.3,0.6) {Yaoundé (BEAC)};
\node[popink, font=\small\bfseries] at (0.5,0.9) {1};
\fill[popink] (-1.0,-0.9) circle (0.09); \node[left,font=\small] at (-1.1,-0.9) {Douala (port)};
\node[popink, font=\small\bfseries] at (-1.4,-0.6) {2};
% Corridors
\draw[-{Stealth[length=3mm]},very thick,popblue] (-0.9,-0.8) .. controls (0.6,1.4) and (2.0,2.8) .. (3.2,3.6);
\node[popblue,font=\small, align=center] at (2.6,2.2){\shortstack{Corridor\\Douala--N'Djamena}};
\draw[-{Stealth[length=3mm]},very thick,popgreen] (-0.8,-0.9) .. controls (1.2,-0.2) and (2.6,0.6) .. (3.8,1.2);
\node[popgreen,font=\small, align=center] at (4.4,0.2){\shortstack{Corridor\\Douala--Bangui}};
% Nord
\draw[-{Stealth[length=2.5mm]},thick] (-4.2,-3.0) -- (-4.2,-2.0); \node[below] at (-4.2,-3.0) {N};
\end{tikzpicture}
\legende{Croquis schématique (sans échelle exacte) : le Cameroun (en orange) au centre de la CEMAC ; 1 : Yaoundé, siège de la BEAC ; 2 : port de Douala ; flèche bleue : corridor vers le Tchad ; flèche verte : corridor vers la Centrafrique.}
\end{popfigure}

\emph{Étape 2 : le commentaire (cinq lignes).} Ce croquis montre que le Cameroun occupe une position centrale dans la CEMAC : il touche cinq des six pays de la sous-région. Son port de Douala est le débouché maritime des pays enclavés, le Tchad et la Centrafrique, desservis par deux corridors commerciaux. Yaoundé, capitale politique, abrite le siège de la BEAC, la banque centrale commune. Le Cameroun apparaît ainsi comme la plaque tournante économique et monétaire de l'Afrique centrale. Cette position explique son rôle moteur dans l'intégration sous-régionale.

\emph{Conclusion.} Le croquis, complété par sa légende, son titre et son orientation, illustre clairement le rôle de pivot du Cameroun dans la sous-région.
\end{exoresolu}

\begin{attention}[Les 5 erreurs qui coûtent des points au Bac]
\begin{enumerate}
\item \textbf{Confondre les organisations} : l'UA est continentale (créée en 2002, siège à Addis-Abeba) ; la CEEAC et la CEMAC sont sous-régionales. Écrire « la CEMAC, organisation continentale » est une erreur éliminatoire.
\item \textbf{Mélanger les sièges des institutions} : BEAC à Yaoundé, COBAC à Libreville, Parlement communautaire à Malabo, Cour de justice à N'Djamena, siège de la CEMAC à Bangui. Apprenez le tableau par cœur et vérifiez-le au brouillon.
\item \textbf{Se tromper de dates} : OUA en 1963, CEEAC en 1983, CEMAC en 1994 (traité de N'Djamena), UA en 2002. Une date fausse dans une définition coûte le point de précision.
\item \textbf{Affirmer sans justifier} : écrire « l'intégration profite au Cameroun » sans exemple (franc CFA, corridor Douala--N'Djamena, libre circulation) ne vaut rien. Chaque argument doit être suivi d'un fait concret.
\item \textbf{Négliger la forme} : sigles non développés à la première occurrence, croquis sans titre, sans légende ni flèche du nord, conclusion absente ou copie non structurée (introduction -- développement -- conclusion). La méthode compte autant que le contenu.
\end{enumerate}
\end{attention}

\newpage

\coursec{Exercices}

\courssub{Je vérifie mes connaissances}

\begin{exercice}[QCM : les organisations d'intégration]
Pour chaque question, une seule réponse est exacte. Recopie la lettre correspondante.
\begin{enumerate}
\item Le siège de l'Union Africaine (UA) se trouve à :
\begin{enumerate}
\item[a)] Abuja (Nigeria) \quad b)] Addis-Abeba (Éthiopie) \quad c)] Nairobi (Kenya) \quad d)] Le Caire (Égypte)
\end{enumerate}
\item La CEMAC compte :
\begin{enumerate}
\item[a)] 4 États membres \quad b)] 5 États membres \quad c)] 6 États membres \quad d)] 11 États membres
\end{enumerate}
\item Le siège de la Banque des États de l'Afrique Centrale (BEAC) est situé à :
\begin{enumerate}
\item[a)] Bangui \quad b)] Libreville \quad c)] Yaoundé \quad d)] Brazzaville
\end{enumerate}
\item La monnaie commune des pays de la CEMAC est :
\begin{enumerate}
\item[a)] le naira \quad b)] le franc CFA (BEAC) \quad c)] le rand \quad d)] l'euro
\end{enumerate}
\end{enumerate}
\end{exercice}

\begin{exercice}[Vrai ou faux, avec justification]
Réponds par \emph{vrai} ou \emph{faux}, puis justifie ta réponse en une ou deux phrases.
\begin{enumerate}
\item L'Union Africaine a succédé à l'Organisation de l'Unité Africaine (OUA) en 2002.
\item Le Cameroun n'appartient qu'à une seule organisation sous-régionale.
\item La libre circulation des personnes et des biens est un objectif de la CEMAC.
\item La ZLECAf (Zone de libre-échange continentale africaine) concerne uniquement l'Afrique centrale.
\end{enumerate}
\end{exercice}

\begin{exercice}[Définitions]
Définis en deux ou trois phrases chacun des termes suivants, en donnant un exemple tiré de la sous-région du Cameroun :
\begin{enumerate}
\item l'intégration sous-régionale ;
\item la CEEAC ;
\item la libre circulation des personnes et des biens.
\end{enumerate}
\end{exercice}

\begin{exercice}[Calculs directs : parts dans la CEMAC]
Le tableau ci-dessous donne le PIB approximatif (en milliards de dollars) des six pays de la CEMAC en 2022 (données simplifiées).

\begin{center}
\begin{tabularx}{\linewidth}{|l|X|}\hline
\rowcolor{popblueL} \textbf{Pays} & \textbf{PIB (milliards de dollars)} \\ \hline
Cameroun & 45 \\ \hline
Gabon & 20 \\ \hline
Congo & 14 \\ \hline
Tchad & 12 \\ \hline
Guinée équatoriale & 12 \\ \hline
Centrafrique & 2,5 \\ \hline
\end{tabularx}
\end{center}

\begin{enumerate}
\item Calcule le PIB total de la CEMAC.
\item Calcule, en pourcentage arrondi à l'unité, la part du Cameroun dans ce PIB total.
\item Que représente la part cumulée du Cameroun et du Gabon ? Conclus en une phrase.
\end{enumerate}
\end{exercice}

\courssub{Je m'entraîne}

\begin{exercice}[Associer institutions et sièges]
Recopie et complète le tableau suivant en associant chaque institution à son siège et à son domaine d'action principal.

\begin{center}
\begin{tabularx}{\linewidth}{|l|X|X|}\hline
\rowcolor{popblueL} \textbf{Institution} & \textbf{Siège} & \textbf{Domaine d'action} \\ \hline
BEAC & \ldots & \ldots \\ \hline
CEMAC (secrétariat exécutif) & \ldots & \ldots \\ \hline
CEEAC (secrétariat général) & \ldots & \ldots \\ \hline
Union Africaine (Commission) & \ldots & \ldots \\ \hline
\end{tabularx}
\end{center}

Sièges à placer : Addis-Abeba, Bangui, Libreville, Yaoundé. Domaines à placer : monnaie commune ; intégration politique continentale ; intégration économique et monétaire de la CEMAC ; intégration de l'Afrique centrale au sens large.
\end{exercice}

\begin{exercice}[Situation concrète : un commerçant camerounais]
M. Etoundi, commerçant à Douala, souhaite exporter du savon produit dans son entreprise vers le Tchad et le Gabon.
\begin{enumerate}
\item Cite deux avantages que l'appartenance du Cameroun à la CEMAC peut lui apporter.
\item Cite deux difficultés concrètes qu'il risque de rencontrer sur le terrain (routes, tracasseries, monnaie, délais).
\item Propose une solution réaliste pour chacune de ces deux difficultés.
\end{enumerate}
\end{exercice}

\begin{exercice}[Construction graphique : le PIB de la CEMAC]
En reprenant le tableau des PIB de l'exercice « Calculs directs » :
\begin{enumerate}
\item Construis sur ton cahier un diagramme en barres représentant le PIB des six pays (échelle conseillée : \SI{1}{cm} pour 5 milliards de dollars ; n'oublie ni le titre ni les légendes des axes).
\item Dégage deux enseignements de ton graphique et rédige-les en deux phrases complètes (inégalités entre pays, poids du Cameroun).
\end{enumerate}
\end{exercice}

\begin{exercice}[Analyse de document : la ZLECAf]
\textbf{Document.} « Entrée en vigueur en 2021, la Zone de libre-échange continentale africaine (ZLECAf) vise à créer un marché unique de plus d'un milliard de consommateurs. Pour le Cameroun, elle ouvre des débouchés au-delà de la CEMAC, mais expose aussi ses entreprises à la concurrence de pays plus industrialisés comme le Nigeria, l'Afrique du Sud ou l'Égypte. »
\begin{enumerate}
\item Présente la ZLECAf en deux phrases (nature, objectif).
\item Releve dans le texte un atout et une menace de la ZLECAf pour le Cameroun.
\item À ton avis, que doit faire l'État camerounais pour que les entreprises locales profitent de la ZLECAf ? Justifie ta réponse en trois lignes.
\end{enumerate}
\end{exercice}

\courssub{Je me prépare au Bac}

\begin{exercice}[Analyse de documents et conclusion rédigée]
\textbf{Document 1.} « Créée en 1994, la CEMAC regroupe le Cameroun, la Centrafrique, le Congo, le Gabon, la Guinée équatoriale et le Tchad. Elle vise une union économique et monétaire : monnaie commune (le franc CFA), banque centrale commune (BEAC, siège à Yaoundé) et libre circulation des personnes et des biens. »

\textbf{Document 2.} Carte de l'Afrique centrale montrant les six États de la CEMAC, les onze États de la CEEAC, les principaux axes routiers et le port de Douala, première porte d'entrée maritime du Tchad et de la Centrafrique.

\textbf{Document 3.} « Malgré les textes, les échanges entre pays d'Afrique centrale restent faibles : moins de 15 \% de leur commerce total. Les causes : mauvais état des routes, productions similaires (pétrole, bois, cacao), lenteurs douanières et conflits dans certains pays. »

\textbf{Consignes :}
\begin{enumerate}
\item Présente la CEMAC (date de création, membres, objectifs) en cinq lignes. (4 points)
\item À partir du document 2, montre le rôle central du Cameroun dans la sous-région. (4 points)
\item À partir du document 3, dégage deux obstacles à l'intégration en Afrique centrale. (4 points)
\item Rédige une conclusion de dix lignes répondant à la question : l'intégration sous-régionale est-elle une réalité ou un simple projet pour le Cameroun ? (8 points)
\end{enumerate}
\end{exercice}

\begin{exercice}[Croquis noté : le Cameroun dans la CEMAC]
Sur le fond de carte de l'Afrique centrale fourni par l'enseignant, réalise un croquis intitulé \emph{« Le Cameroun dans la CEMAC »}.

\textbf{Informations à placer obligatoirement :}
\begin{enumerate}
\item les six États membres de la CEMAC (hachurés ou coloriés d'une même couleur) ;
\item le siège de la BEAC (Yaoundé) et le siège du secrétariat exécutif de la CEMAC (Bangui) ;
\item le port de Douala et l'axe Douala--N'Djamena et Douala--Bangui (corridors d'approvisionnement des pays enclavés) ;
\item une légende organisée et un titre en majuscules.
\end{enumerate}

\textbf{Barème indicatif :} exactitude du fond de carte (4 points) ; sélection et localisation des informations (8 points) ; légende et titre (4 points) ; propreté et soin (4 points).
\end{exercice}

\begin{exercice}[Dissertation type Baccalauréat technique]
\textbf{Sujet :} « L'intégration sous-régionale : atout ou illusion pour le développement du Cameroun ? »

\textbf{Consignes :} après avoir bien analysé le sujet et défini les termes clés (\emph{intégration sous-régionale}, \emph{développement}), tu rédigeras une copie comportant :
\begin{enumerate}
\item une \textbf{introduction} (accroche, définition des termes, annonce du problème et du plan) ;
\item une \textbf{première partie} montrant que l'intégration est un atout pour le Cameroun (marché élargi, monnaie commune, libre circulation, rôle du port de Douala, paix et sécurité) ;
\item une \textbf{deuxième partie} montrant les limites qui font douter de cette intégration (faiblesse des échanges, routes dégradées, tracasseries, concurrence, conflits) ;
\item une \textbf{conclusion} dans laquelle tu prendras clairement position et proposeras deux solutions pour renforcer l'intégration.
\end{enumerate}
Ta copie sera notée sur la pertinence des arguments, la correction de la langue et l'organisation du devoir. (20 points)
\end{exercice}

\newpage

\coursec{Corrigés des exercices}

\coursec{Corrigés des activités et exercices — Le Cameroun et les structures d'intégration sous-régionales}

\begin{corrige}[Exercice 1 — QCM : les organisations d'intégration]
\begin{enumerate}
\item \textbf{b) Addis-Abeba (Éthiopie).} Le siège de la Commission de l'Union Africaine est à Addis-Abeba, capitale de l'Éthiopie, qui abritait déjà le siège de l'OUA.
\item \textbf{c) 6 États membres.} La CEMAC regroupe le Cameroun, la Centrafrique, le Congo, le Gabon, la Guinée équatoriale et le Tchad.
\item \textbf{c) Yaoundé.} La Banque des États de l'Afrique Centrale (BEAC) a son siège à Yaoundé, capitale du Cameroun.
\item \textbf{b) le franc CFA (BEAC).} Les six pays de la CEMAC utilisent le franc CFA émis par la BEAC.
\end{enumerate}
\textbf{Point de vigilance :} ne pas confondre la CEMAC (6 États) et la CEEAC (11 États) ; ne pas confondre le siège de la BEAC (Yaoundé) et celui du secrétariat exécutif de la CEMAC (Bangui).
\end{corrige}

\begin{corrige}[Exercice 2 — Vrai ou faux, avec justification]
\begin{enumerate}
\item \textbf{Vrai.} L'Union Africaine a été créée en 2002 à Durban (Afrique du Sud) pour succéder à l'Organisation de l'Unité Africaine (OUA), fondée en 1963.
\item \textbf{Faux.} Le Cameroun appartient à plusieurs organisations sous-régionales : il est membre de la CEMAC (6 États) et de la CEEAC (11 États), sans compter l'Union Africaine au niveau continental.
\item \textbf{Vrai.} La libre circulation des personnes et des biens fait partie des objectifs de l'union économique et monétaire de la CEMAC, même si sa mise en œuvre reste incomplète sur le terrain.
\item \textbf{Faux.} La ZLECAf est une zone de libre-échange \emph{continentale} : elle concerne l'ensemble des États africains signataires, et non la seule Afrique centrale.
\end{enumerate}
\end{corrige}

\begin{corrige}[Exercice 3 — Définitions]
\begin{enumerate}
\item \textbf{L'intégration sous-régionale} est le processus par lequel des États voisins d'une même région s'associent pour mettre en commun leurs politiques économiques, monétaires ou sécuritaires (marché commun, monnaie commune, libre circulation). \emph{Exemple :} le Cameroun et ses voisins du bassin du Congo au sein de la CEMAC et de la CEEAC.
\item \textbf{La CEEAC} (Communauté Économique des États de l'Afrique Centrale) est une organisation créée en 1983 qui regroupe 11 États d'Afrique centrale, dont le Cameroun, le Gabon, le Tchad, la République démocratique du Congo et l'Angola. Elle œuvre pour la coopération économique, la paix et la sécurité dans la sous-région ; son secrétariat général est à Libreville.
\item \textbf{La libre circulation des personnes et des biens} est le droit, pour les citoyens et les marchandises des États membres, de se déplacer et de commercer à l'intérieur de l'espace communautaire sans visa ni droits de douane excessifs. \emph{Exemple :} un commerçant camerounais peut, en principe, vendre ses produits au Tchad ou au Gabon dans le cadre de la CEMAC.
\end{enumerate}
\end{corrige}

\begin{corrige}[Exercice 4 — Calculs directs : parts dans la CEMAC]
\begin{enumerate}
\item PIB total de la CEMAC :
\[ 45 + 20 + 14 + 12 + 12 + 2{,}5 = 105{,}5 \text{ milliards de dollars.} \]
\item Part du Cameroun :
\[ \frac{45}{105{,}5} \times 100 \approx 42{,}65 \approx 43\,\%. \]
Le Cameroun représente environ \textbf{43 \%} du PIB de la CEMAC.
\item Part cumulée Cameroun + Gabon :
\[ \frac{45 + 20}{105{,}5} \times 100 = \frac{65}{105{,}5} \times 100 \approx 61{,}6 \approx 62\,\%. \]
\textbf{Conclusion :} à eux deux, le Cameroun et le Gabon pèsent environ 62 \% du PIB de la CEMAC, ce qui montre une forte concentration économique et le poids dominant du Cameroun dans la sous-région.
\end{enumerate}
\textbf{Point de vigilance :} bien additionner les six valeurs avant de calculer les pourcentages ; arrondir uniquement à la fin du calcul.
\end{corrige}

\begin{corrige}[Exercice 5 — Associer institutions et sièges]
\begin{center}
\begin{tabularx}{\linewidth}{|l|X|X|}\hline
\rowcolor{popblueL} \textbf{Institution} & \textbf{Siège} & \textbf{Domaine d'action} \\ \hline
BEAC & Yaoundé & Monnaie commune \\ \hline
CEMAC (secrétariat exécutif) & Bangui & Intégration économique et monétaire de la CEMAC \\ \hline
CEEAC (secrétariat général) & Libreville & Intégration de l'Afrique centrale au sens large \\ \hline
Union Africaine (Commission) & Addis-Abeba & Intégration politique continentale \\ \hline
\end{tabularx}
\end{center}
\textbf{Point de vigilance :} le siège de la BEAC est à Yaoundé, mais le secrétariat exécutif de la CEMAC est à Bangui : ne pas les inverser, c'est une erreur classique à l'examen.
\end{corrige}

\begin{corrige}[Exercice 6 — Situation concrète : un commerçant camerounais]
\begin{enumerate}
\item \textbf{Deux avantages :}
\begin{itemize}
\item la \cle{monnaie commune} (franc CFA) : pas de problème de change entre le Cameroun, le Tchad et le Gabon ;
\item la \cle{libre circulation des biens} dans la CEMAC : suppression ou réduction des droits de douane, accès à un marché élargi de plusieurs millions de consommateurs.
\end{itemize}
\item \textbf{Deux difficultés possibles :}
\begin{itemize}
\item le mauvais état des routes (axe Douala--N'Djamena difficile en saison des pluies), qui rallonge les délais et augmente les coûts ;
\item les tracasseries aux frontières et les lenteurs douanières (contrôles multiples, demandes indues de paiement).
\end{itemize}
\item \textbf{Solutions réalistes :}
\begin{itemize}
\item pour les routes : grouper ses envois avec d'autres commerçants (transport mutualisé) ou utiliser le corridor ferroviaire/route quand il est opérationnel, et planifier les expéditions en saison sèche ;
\item pour les tracasseries : se munir de tous les documents communautaires en règle (certificat d'origine CEMAC, factures) et se rapprocher de la chambre de commerce ou des guichets uniques pour signaler les abus.
\end{itemize}
\end{enumerate}
\end{corrige}

\begin{corrige}[Exercice 7 — Construction graphique : le PIB de la CEMAC]
\begin{enumerate}
\item \textbf{Construction attendue.} Axe horizontal : les six pays ; axe vertical : le PIB en milliards de dollars, gradué de 5 en 5 (échelle \SI{1}{cm} pour 5 milliards). Hauteurs des barres : Cameroun \SI{9}{cm} ; Gabon \SI{4}{cm} ; Congo \SI{2,8}{cm} ; Tchad \SI{2,4}{cm} ; Guinée équatoriale \SI{2,4}{cm} ; Centrafrique \SI{0,5}{cm}. Titre obligatoire : « PIB des pays de la CEMAC en 2022 (milliards de dollars) ».
\item \textbf{Deux enseignements possibles :}
\begin{itemize}
\item Le graphique montre de fortes inégalités entre les pays : le PIB du Cameroun (45) est 18 fois celui de la Centrafrique (2,5).
\item Le Cameroun est la première économie de la CEMAC : sa barre est plus de deux fois plus haute que celle du Gabon, le deuxième, ce qui confirme son poids dominant (environ 43 \% du total).
\end{itemize}
\end{enumerate}
\textbf{Point de vigilance :} une copie sans titre ou sans légende des axes perd des points ; les barres doivent avoir la même largeur et être régulièrement espacées.
\end{corrige}

\begin{corrige}[Exercice 8 — Analyse de document : la ZLECAf]
\begin{enumerate}
\item La ZLECAf est une zone de libre-échange à l'échelle du continent africain, entrée en vigueur en 2021. Son objectif est de créer un marché unique de plus d'un milliard de consommateurs en supprimant progressivement les barrières douanières entre les pays africains.
\item \textbf{Atout :} « elle ouvre des débouchés au-delà de la CEMAC » (nouveaux marchés pour les produits camerounais). \textbf{Menace :} elle « expose ses entreprises à la concurrence de pays plus industrialisés comme le Nigeria, l'Afrique du Sud ou l'Égypte ».
\item \textbf{Réponse attendue (exemple) :} l'État camerounais doit aider les entreprises locales à devenir compétitives : améliorer les routes et l'électricité pour réduire les coûts de production, former les jeunes aux métiers de l'industrie, faciliter l'accès au crédit et protéger temporairement les secteurs fragiles. Sans ces efforts, la ZLECAf profitera surtout aux pays déjà industrialisés.
\end{enumerate}
\end{corrige}

\begin{corrige}[Exercice 9 — Analyse de documents et conclusion rédigée]
\textbf{1. Présentation de la CEMAC (4 points).} Créée en 1994, la CEMAC regroupe six États : le Cameroun, la Centrafrique, le Congo, le Gabon, la Guinée équatoriale et le Tchad. Elle vise une union économique et monétaire entre ses membres. Ses instruments principaux sont la monnaie commune (le franc CFA), une banque centrale commune (la BEAC, dont le siège est à Yaoundé) et la libre circulation des personnes et des biens.

\textbf{2. Rôle central du Cameroun (4 points).} D'après le document 2, le Cameroun occupe une position centrale : le port de Douala est la première porte d'entrée maritime du Tchad et de la Centrafrique, deux pays enclavés qui dépendent des corridors Douala--N'Djamena et Douala--Bangui pour leur approvisionnement. Le Cameroun est donc le poumon commercial de la sous-région.

\textbf{3. Deux obstacles (4 points).} Le document 3 permet de dégager : le mauvais état des routes et les lenteurs douanières ; la similitude des productions (pétrole, bois, cacao), qui limite les échanges entre pays, et les conflits dans certains États. Résultat : moins de 15 \% du commerce total se fait entre pays d'Afrique centrale.

\textbf{4. Conclusion rédigée (8 points) — exemple de réponse.} L'intégration sous-régionale est pour le Cameroun une réalité institutionnelle mais encore un projet inachevé sur le plan économique. Elle est une réalité parce que les institutions existent et fonctionnent : la monnaie commune, la BEAC à Yaoundé, les textes sur la libre circulation. Le Cameroun y joue un rôle moteur grâce au port de Douala et au poids de son économie. Mais elle reste un projet car les échanges entre pays de la sous-région demeurent faibles, les routes sont dégradées, les frontières restent lourdes à franchir et les productions se ressemblent. Pour que l'intégration devienne pleinement réelle, il faut investir dans les infrastructures, appliquer réellement les textes communautaires et diversifier les économies. L'intégration est donc un chantier en cours, dont le Cameroun a tout intérêt à être le bâtisseur.

\textbf{Points de vigilance :} citer les documents dans la réponse ; respecter le nombre de lignes demandé ; la conclusion doit trancher clairement la question posée.
\end{corrige}

\begin{corrige}[Exercice 10 — Croquis noté : le Cameroun dans la CEMAC]
\textbf{Corrigé-type attendu :}
\begin{enumerate}
\item Les six États de la CEMAC (Cameroun, Gabon, Congo, Guinée équatoriale, Tchad, Centrafrique) sont correctement localisés et coloriés d'une même couleur, ce qui les distingue des États voisins non membres (Nigeria, RDC, Angola\ldots).
\item Yaoundé est marquée d'un symbole (étoile ou point) avec la mention « siège de la BEAC » ; Bangui avec la mention « siège du secrétariat exécutif de la CEMAC ».
\item Le port de Douala est placé sur la côte ; les axes Douala--N'Djamena (vers le nord) et Douala--Bangui (vers l'est) sont tracés en traits ou flèches, avec la mention « corridors d'approvisionnement des pays enclavés ».
\item Le titre est en majuscules (« LE CAMEROUN DANS LA CEMAC ») et la légende est organisée en cadre : couleur des États membres, symboles des sièges, traits des corridors, orientation (flèche du nord).
\end{enumerate}
\textbf{Barème indicatif :} fond de carte exact (4 points) ; sélection et localisation des informations (8 points) ; légende et titre (4 points) ; propreté et soin (4 points).

\textbf{Points de vigilance :} toute information sur le croquis doit figurer dans la légende ; ne pas surcharger le croquis d'informations non demandées ; écrire lisiblement et horizontalement les noms des pays.
\end{corrige}

\begin{corrige}[Exercice 11 — Dissertation type Baccalauréat technique]
\textbf{Corrigé-type (plan détaillé et éléments de rédaction).}

\textbf{Introduction.} \emph{Accroche :} depuis les indépendances, les États africains cherchent à s'unir pour peser davantage dans l'économie mondiale. \emph{Définitions :} l'intégration sous-régionale est le rapprochement économique et politique d'États voisins (CEMAC, CEEAC) ; le développement est l'amélioration durable des conditions de vie des populations. \emph{Problématique :} l'intégration sous-régionale est-elle un véritable atout pour le développement du Cameroun, ou une simple illusion ? \emph{Annonce du plan :} nous montrerons d'abord qu'elle est un atout, puis qu'elle connaît des limites qui font douter de sa réalité.

\textbf{Première partie : un atout pour le Cameroun.}
\begin{itemize}
\item Un marché élargi : les produits camerounais accèdent à plusieurs millions de consommateurs de la CEMAC et de la CEEAC.
\item La monnaie commune (franc CFA) et la BEAC (siège à Yaoundé) facilitent les échanges et garantissent la stabilité monétaire.
\item La libre circulation des personnes et des biens réduit les coûts du commerce.
\item Le rôle du port de Douala, porte d'entrée du Tchad et de la Centrafrique, fait du Cameroun le pivot économique de la sous-région.
\item La coopération sécuritaire (lutte contre Boko Haram, piraterie dans le golfe de Guinée) renforce la paix, condition du développement.
\end{itemize}

\textbf{Deuxième partie : des limites qui font douter.}
\begin{itemize}
\item La faiblesse des échanges : moins de 15 \% du commerce des pays d'Afrique centrale se fait entre eux.
\item Les routes dégradées et les corridors mal entretenus renchérissent les transports.
\item Les tracasseries et lenteurs douanières contredisent la libre circulation annoncée.
\item Les productions similaires (pétrole, bois, cacao) limitent la complémentarité ; la concurrence des grands pays (Nigeria) menace les entreprises locales.
\item Les conflits et l'insécurité dans certains États freinent les échanges.
\end{itemize}

\textbf{Conclusion.} Prise de position : l'intégration est un atout réel mais encore incomplet ; elle reste en partie une illusion tant que les textes ne sont pas appliqués. \emph{Deux solutions :} investir dans les infrastructures (routes, corridors Douala--N'Djamena et Douala--Bangui) et appliquer effectivement la libre circulation en supprimant les tracasseries aux frontières, tout en diversifiant l'économie camerounaise.

\textbf{Barème indicatif (20 points) :} introduction (accroche, définitions, problématique, plan) : 4 points ; première partie argumentée : 5 points ; deuxième partie argumentée : 5 points ; conclusion avec prise de position et deux solutions : 4 points ; langue et organisation : 2 points.

\textbf{Points de vigilance :} définir les deux termes clés du sujet dès l'introduction ; illustrer chaque argument par un exemple précis (BEAC, Douala, franc CFA) ; éviter les catalogues d'idées sans lien logique ; la conclusion doit répondre clairement à la question posée.
\end{corrige}

\newpage

\coursec{Fiche bilan}

\begin{popfigure}
\begin{tikzpicture}[mindmap, grow cyclic, every node/.style={concept, font=\small},
root concept/.append style={concept color=popblue, font=\bfseries},
level 1/.append style={level distance=3.4cm, sibling angle=51.4},
level 2/.append style={level distance=2.2cm, font=\scriptsize}]
\node[root concept]{Le Cameroun et l'intégration sous-régionale}
  child[concept color=poporange]{ node {Comprendre l'intégration}
    child { node {Définition : regroupement d'États voisins} }
    child { node {Buts : paix, économie, libre circulation} }
  }
  child[concept color=popgreen]{ node {Union Africaine (UA)}
    child { node {Créée en 2002 à Durban} }
    child { node {Siège : Addis-Abeba} }
    child { node {Succède à l'OUA (1963)} }
  }
  child[concept color=popblue]{ node {CEEAC}
    child { node {Créée en 1983 (Libreville)} }
    child { node {11 États d'Afrique centrale} }
  }
  child[concept color=popgold]{ node {CEMAC}
    child { node {Trait\u00e9 de N'Djamena (1994)} }
    child { node {6 États, monnaie : FCFA} }
  }
  child[concept color=poppurple]{ node {Institutions communes}
    child { node {BEAC, BDEAC, COBAC} }
    child { node {Parlement et Cour CEMAC} }
  }
  child[concept color=poppink]{ node {Acquis et obstacles}
    child { node {Acquis : monnaie unique, paix} }
    child { node {Obstacles : frontières, faibles échanges} }
  }
  child[concept color=popmuted]{ node {Acteurs et citoyenneté}
    child { node {Rôle de l'État et des citoyens} }
    child { node {Solidarité, paix, développement} }
  };
\end{tikzpicture}
\legende{Carte mentale de la leçon : le Cameroun et les structures d'intégration sous-régionales.}
\end{popfigure}

\begin{bilan}
\textbf{Définitions clés}
\begin{itemize}
\item \cle{Intégration sous-régionale} : processus par lequel des États voisins s'unissent pour coopérer dans les domaines politique, économique, sécuritaire et culturel.
\item \cle{Union Africaine (UA)} : organisation panafricaine créée en 2002 à Durban (Afrique du Sud), siège à Addis-Abeba ; elle succède à l'OUA (1963). Objectifs : unité, paix, développement du continent.
\item \cle{CEEAC} : Communauté Économique des États de l'Afrique Centrale, créée en 1983 à Libreville ; 11 membres ; siège à Libreville (Gabon).
\item \cle{CEMAC} : Communauté Économique et Monétaire de l'Afrique Centrale, créée par le traité de N'Djamena (1994) ; 6 membres : Cameroun, Gabon, Congo, Tchad, Centrafrique, Guinée équatoriale.
\end{itemize}

\textbf{Repères à connaître par c\oe ur}
\begin{center}
\begin{tabularx}{\linewidth}{|l|X|X|}\hline
\rowcolor{popblueL} \textbf{Organisation} & \textbf{Création / siège} & \textbf{Rôle principal} \\ \hline
UA & 2002, Addis-Abeba & Unité, paix et développement de l'Afrique \\ \hline
CEEAC & 1983, Libreville & Coopération régionale des 11 États d'Afrique centrale \\ \hline
CEMAC & 1994, Bangui & Intégration économique et monétaire (6 États) \\ \hline
BEAC & 1972, Yaoundé & Banque centrale, émission du FCFA \\ \hline
BDEAC & 1975, Brazzaville & Financement du développement \\ \hline
\end{tabularx}
\end{center}

\textbf{Méthodes express}
\begin{itemize}
\item Pour analyser un texte sur l'intégration : identifier l'organisation, sa date, ses objectifs, puis le rôle du Cameroun.
\item Pour rédiger une conclusion : rappeler les acquis (monnaie unique, paix, libre circulation) puis nuancer avec les obstacles (faibles échanges, frontières, conflits).
\item Toujours justifier sa réponse par un exemple concret (FCFA, corridor Douala--N'Djamena, force multinationale).
\end{itemize}
\end{bilan}

\begin{aretenir}[Checklist avant le Bac]
\begin{itemize}
\item $\square$ Je sais définir l'intégration sous-régionale en une phrase.
\item $\square$ Je connais la date et le lieu de création de l'UA (2002, Durban).
\item $\square$ Je sais que l'UA succède à l'OUA (1963).
\item $\square$ Je cite les 6 pays membres de la CEMAC.
\item $\square$ Je connais le siège de la CEEAC (Libreville) et de la CEMAC (Bangui).
\item $\square$ Je sais que la BEAC, à Yaoundé, émet le franc CFA.
\item $\square$ Je cite au moins deux acquis de l'intégration (monnaie unique, libre circulation, sécurité).
\item $\square$ Je cite au moins deux obstacles (faibles échanges intra-africains, conflits, frontières).
\item $\square$ Je sais expliquer le rôle du Cameroun dans sa sous-région.
\item $\square$ Je sais rédiger une conclusion structurée et justifiée.
\end{itemize}
\end{aretenir}

\findecours

\end{document}
